% concatenated from Arxiv_v2.tex
%\documentclass[11pt,showkeys]{amsart}
\documentclass[11pt]{amsart}


% \usepackage[notcite, notref, color]{showkeys}
\usepackage[english]{babel}
 \usepackage{enumerate}
 \usepackage{xcolor}
%\usepackage{fourier}
\def\X{\mathcal{X}}\usepackage{amscd}
\usepackage{amsthm}
\usepackage[centertags]{amsmath}

\usepackage{newlfont}
\usepackage{graphicx}
 % \usepackage[usenames]{xcolor}
\usepackage{amsfonts, amssymb,accents}
\usepackage{mathrsfs}
\usepackage{latexsym}
\usepackage{bbold}
\usepackage{enumitem}
%\usepackage[french]{babel}
%\usepackage{pstricks,pst-node,pst-fill}
% \usepackage{tikz}\usepackage{verbatim}
% \usepackage{tabulary,tabularx}
% \usepackage[all]{xy}
% \usepackage{tikz}
% \usetikzlibrary{positioning}
% \usetikzlibrary{shapes.multipart}
  \definecolor{colorcita}{RGB}{21,86,130}
  \definecolor{colorref}{RGB}{5,10,177}
  \definecolor{colorweb}{RGB}{177,6,38}
 \usepackage[colorlinks=true,linkcolor=colorref,citecolor=colorcita,urlcolor=colorweb]{hyperref}

%\usepackage{hyperref}


%\usepackage[spanish]{babel}
% \usepackage[latin1]{inputenc}
% \usepackage[notcite, notref, color]{showkeys}

\newtheorem{hip}{Hypothesis}
\newtheorem{theorem}{Theorem}[section]
\newtheorem{problem}[theorem]{Problem}
\newtheorem{proposition}[theorem]{Proposition}
\newtheorem{corollary}[theorem]{Corollary}
\newtheorem{lemma}[theorem]{Lemma}
\newtheorem{remark}[theorem]{Remark}
\newtheorem{definition}[theorem]{Definition}
\newtheorem{example}[theorem]{Example}
\newtheorem{criterion}[theorem]{Criterion}
\newtheorem{application}[theorem]{Application}
\newtheorem{question}[theorem]{Question}
\newcommand{\dis}{\displaystyle}

\newcommand*{\nota}[1]{\marginpar{\tiny #1}} %% para la versi\'{o}n definitiva sustituir por \newcommand*{\nota}[1]{}

\newcommand{\kkk}{\mathbf k}
\newcommand{\PP}{\mathscr P}
\newcommand{\zC}{\mathbb K}
\newcommand{\zR}{\mathbb R}
\newcommand{\zN}{\mathbb N}
\newcommand{\zK}{\mathbb K}
\renewcommand{\u}{\mathfrak A}
\renewcommand{\H}{H(\mathbb{C})}
\newcommand{\C}{\mathbb C}
\newcommand{\zZ}{\mathbb Z}
\newcommand{\sub}{\subseteq}
\newcommand{\zT}{\mathbb T}
\newcommand{\zD}{\mathbb D}
\newcommand{\cupdot}{\mathbin{\mathaccent\cdot\cup}}
\newcommand{\f}{\frac}
\newcommand{\eps}{\epsilon}
\newcommand{\AP}{\mathcal {AP}}
\newcommand{\kAP}{\mathcal {AP}\!_b}%{\mathcal {AP}_k}
\newcommand{\A}{\mathcal {F}}
\newcommand{\PF}{\PP_{\mathcal {F}}}
\newcommand{\PB}{\mathcal {PB}}
\newcommand{\BD}{\overline{\mathcal {BD}}}
\newcommand{\UD}{\overline{\mathcal {D}}}
\newcommand{\LD}{\overline{\mathcal {D}}}
\newcommand{\F}{\mathcal {F}}

\newcommand{\zS}{\mathcal {S}}
\newcommand{\zM}{M_\varphi^*}
\newcommand{\interior}[1]{\accentset{\smash{\raisebox{-0.12ex}{$\scriptstyle\circ$}}}{#1}\rule{0pt}{2.3ex}}
\fboxrule0.0001pt \fboxsep0pt
 
\renewcommand{\l}{\left(}
\renewcommand{\r}{\right)}
%\newcommand{\z1}{\textbb\mathbb 1}
\renewcommand{\baselinestretch}{1.3}

\textwidth=17.4cm \textheight=23cm \hoffset=-20.5mm \voffset=-5mm
\parskip 7.2pt



\title{Frequently recurrent backward shifts}
\keywords{Frequently recurrent operators,  hypercylic operators, weighted shifts}
\subjclass[2010]{
47A16, %Cyclic vectors, hypercyclic and chaotic operators
    37B20  	%Notions of recurrence and recurrent behavior in dynamical systems
	47B37  	%Linear operators on special spaces (weighted shifts, operators on sequence spaces, etc.)
  37A45  %	Relations between ergodic theory and number theory
	47A35   	%Ergodic theory of linear operators
 %37A25  %	Ergodicity, mixing, rates of mixing
		%11B25 % 	Arithmetic progressions		
	%47H60,  Multilinear and polynomial operators
	%37F10, Complex dynamical systems - Polynomials; rational maps; entire and meromorphic functions
	%30D20, %Functions of a complex variable, Entire functions, general theory 
	%30K99  % Functions of a complex variable -Universal holomorphic functions -None of the above, but in this section
	%46E50 %  	Spaces of differentiable or holomorphic functions on infinite-dimensional spaces
	%46T25%    	Functional analysis	Holomorphic maps
	%32H50.  %Several complex variables and analytic spaces	Iteration problems
	%32A19 % 	Holomorphic functions of several complex variables  	Normal families of functions, mappings
%15A69 %Multilinear algebra, tensor products,
%46G25% 	(Spaces of) multilinear mappings, polynomials 
%46G20% 	Infinite-dimensional holomorphy
}



\author{Rodrigo Cardeccia, Santiago Muro}
\address{Instituto Balseiro, Universidad Nacional de Cuyo – C.N.E.A. and CONICET, San Carlos de Bariloche, Rep\'ublica Argentina} \email{rodrigo.cardeccia@ib.edu.ar} 
\address{FCEIA, UNIVERSIDAD NACIONAL DE ROSARIO AND   CIFASIS, CONICET}
\email{muro@cifasis-conicet.gov.ar}

\thanks{Partially supported by ANPCyT PICT 2015-2224, UBACyT 20020130300052BA, PIP 11220130100329CO and CONICET}



\begin{document}
\begin{abstract}
 We study  frequently recurrent backward shift operators on Fréchet sequence spaces. We prove that if a backward shift admits a nonzero frequently recurrent vector, then it supports a dense set of such vectors, and hence the operator is frequently recurrent. As a consequence, we provide two different characterizations for frequently recurrent backward shift operators,
which also imply the dense lineability of the respective set of
frequently recurrent vectors for this class of operators.
%     We study frequently recurrent backward shifts on Fr\'echet sequence spaces. We prove that a backward shift that supports a nonzero frequently recurrent vector is automatically frequently recurrent. As a consequence of the construction we provide a characterization of frequent recurrence for backward shifts.

    For the construction of frequently recurrent vectors it is important to consider some arithmetic structure on the sets of return times along with their positive lower density. To achieve this, we introduce the recurrent arithmetic thickening of a sequence of subsets of $\zN$. We show that they provide an extension, that takes into account lower density, of a well-known result of Furstenberg  determining that sets of return times contain $IP$-sets. 
\end{abstract}

\maketitle






\section*{Introduction}

The central notion in the dynamics of linear operators is that of hypercyclicity. A linear operator \( T \) on a Fr\'echet space $X$ is said to be hypercyclic if it possesses a dense orbit, i.e. the set \( \{T^n x : n \in \mathbb{N}\} \) is dense in the space for some \( x \in X \).
Furthermore, if an orbit visits each nonempty open set with visiting times that grow at most linearly (or equivalently, the sets of visiting times have positive lower density), then the operator is called frequently hypercyclic. Since its introduction in \cite{BayGri06}, frequent hypercyclicity has been a fundamental concept in the development of linear dynamics.

The class of weighted shift operators (unilateral and bilateral) constitutes one of the most fundamental classes of operators and serves as a rich source of examples and counterexamples in linear dynamics; see for example, \cite{BayGri07,BayRuz15,Bes16,BonGro18,ChaErnMen16,charpentier2019chaos}. The characterization of hypercyclic weighted shifts is well-known, and in \cite{BayRuz15}, a characterization of frequently hypercyclic weighted shift operators on \(\ell_p\) and \(c_0\) was presented (see also \cite{BonGro18,charpentier2019chaos} for extensions to more general Fréchet sequence spaces).

While recurrence is a classical research topic in  dynamical systems, only recently in \cite{CosManPar14}, a systematic study on recurrence for linear operators was initiated. This line of research has been intensive in recent years \cite{abakumov2024orbits,BonGroLopPer22JFA,CardeMur22,CardeMur22multipleScand,GriLop23,GriLop25,lopez2024recurrent,LopMen25}. A vector \( x \) is (frequently) recurrent if the set of visiting times to each neighborhood of \( x \) is infinite (and has positive lower density). An operator is (frequently) recurrent if it has a dense set of (frequently) recurrent vectors.

Returning to weighted shifts, a result due to Chan and Seceleanu \cite{ChaSec12} implies that on \(\ell_p(\mathbb{N})\) and \(\ell_p(\mathbb{Z})\) for \(1 \leq p < \infty\), every recurrent weighted shift is hypercyclic and moreover, the existence of a single nonzero recurrent vector is sufficient to establish hypercyclicity. In particular, we have a zero-one type law for the set of nonzero recurrent vectors of a weighted shift: it is either empty or dense in the space. These zero-one laws properties  were extended to general Fr\'echet sequence spaces \cite{bonilla2025zeroEdinburgh,he2018jf-class} and investigated in other contexts, like adjoint of multipliers on function spaces \cite{ChSe10,bonilla2025zeroEdinburgh},   $\Gamma$-supercyclicity \cite{abakumov2024orbits} or chain recurrence for weighted shifts \cite{lopez2025shifts}. It is worth noting that zero-one laws do not hold for weighted shifts on directed trees \cite{abakumov2024several} or composition operators \cite{ChSe10}. 

%In this article, we will focus on frequently recurrent backward shifts inspired by this zero-one law perspective. 

In this article, we prove a zero-one law for frequently recurrent backward shifts and provide several applications. We will need to construct new recurrent vectors from a given one.
A well known property on (non-linear) dynamical systems is that the set of visiting times of a recurrent vector possesses a rich arithmetic structure. Indeed, it was proved by Furstenberg (see \cite[Theorem 2.17]{Fur81}) that they always contain $IP$-sets. Recall that a subset  $A\subset\mathbb N$ is an $IP$-set if there are natural numbers $p_1\le p_2\le\dots$ such that $A$ consists of
the numbers $p_i$ together with all finite sums $p_{i_1}+\dots + p_{i_k}$ with
$i_1<\dots<i_k$. Thus, if we want to construct a recurrent vector by pasting an appropriate sequence of vectors, we must do it carefully so that we respect the arithmetic structure. When we study frequent recurrence, there is a new issue: this arithmetic structure should somehow care about the lower density of the sets involved. Indeed, an inspection of the proof of   \cite[Theorem 2.17]{Fur81}, shows that the $IP$-set constructed there does not need to have positive lower density, no matter the properties of the recurrent vector. 

% It is then natural to seek for an appropriate extension of Furstenberg's result, that takes into account positive lower density.
% So we ask:

% Is there a family of sets $\mathcal F$ that preserves both the arithmetic structure and the positive lower density? that is, is there a family $\mathcal F$ containing $IP$-sets and the positive lower density sets such that an operator is $\mathcal F$-recurrent if and only if it is frequently recurrent? 

To solve this problem  we introduce what we call the \emph{recurrent arithmetic thickening} ($RAT$ from now on) of a sequence $(A_p)_p$ of subsets of $\zN$, denoted by $RAT((A_p)_p)$.  We show that these sets appear naturally in the context of recurrence: the sets of return times of recurrent points (on any dynamical system) always contain a $RAT$-set. The recurrent arithmetic thickening is always an $IP$-set and preserves the positivity of the lower density, so it becomes crucial for our  constructions of frequently recurrent vectors. The introduction of $RAT$-families allows us to prove an extension of \cite[Theorem 2.17]{Fur81} to arbitrary families.


The main results of the article are as follows: if a backward shift supports a non-trivial frequently recurrent vector, then there must be a dense set of such vectors. Equivalently, the operator is frequently recurrent. Moreover, using the idea of the construction of the dense set of frequently recurrent vectors, and by an application of the zero-one law, we can provide two different characterizations for frequent recurrence: the first one is in the same vein as the characterization of frequently hypercyclic shifts obtained in \cite{BonGro18}, while the second one is related to the existence of orbits that frequently approximate ``big'' subsets of the space $X$. We summarize our main results in a single statement.
%, that in the unilateral  case,  reads as follows.

\begin{theorem}\label{main theorem}
Let $X$ be a Fr\'echet sequence space with unconditional basis $(e_n)_{n\in \mathbb N_0}$ (resp. $(e_n)_{n\in \mathbb Z}$), let $B$ be the respective backward shift operator on $X$ and assume that it is continuous. 
 The following are equivalent:
\begin{enumerate}
    \item $B$ supports a nonzero frequently recurrent vector.
       \item $B$ is frequently recurrent.
    \item There is a dense vector subspace in which every nonzero vector is frequently recurrent and cyclic for $B$.
     \item  There are $(\varepsilon_p)_{p\ge 0}\subset \mathbb R_+$ with $\varepsilon_p\to 0$, and sets of positive lower density $A_p$  
     such that: 
\begin{enumerate}
    \item [i)]for every $p\in\mathbb N_0$, $\sum_{n\in A_p}e_{n+p}$ converges; 
    \item [ii)]
    for every $q\in\mathbb N_0$ and every  $m\in A_q$,
    $$
    \quad\Big\|\sum_{\overset{N\in RAT((A_p)_p)}{N>m}} e_{N-m}\Big\|_q<\varepsilon_q, \qquad \Big(\textrm{resp. } \Big\|\sum_{\overset{N\in RAT((A_p)_p)}{N\neq m}} e_{N-m}\Big\|_q<\varepsilon_q\Big).
    $$
    % $$\quad\Big\|\sum_{p=0}^\infty \sum_{n\in A_p, n>m} \sum_{j=0}^p e_{n-m+N_j}\Big\|_q<\varepsilon_q, \qquad \Big(\textrm{resp. }\Big\|\sum_{p=0}^\infty \sum_{n\in A_p, n\ne m} \sum_{j=0}^p e_{n-m+N_j}\Big\|_q<\varepsilon_q\Big).$$
        \end{enumerate}
   \item There are a set $A\subseteq \mathbb N_0$ of positive lower density and $x\in X$ with $supp(x)\subseteq A$, such that $x$ is a frequently hypercyclic vector for $\mathfrak B_1(A):=\{y\in X: supp(y)\subseteq A \text{ and } |y_n|<1 \text { for every } n\geq 0\}$. 
    
\end{enumerate}
\end{theorem}
We refer to Section \ref{section RAT} for the definition $RAT((A_p)_p)$ and to Section \ref{subsection segunda aplicacion} for the definition of \emph{frequent hypercyclicity for $Y$}, where $Y \sub X$ is a given subset of $X$ (the latter was originally introduced in \cite{Gro19bilateral}). 


The results presented in the above statement depend on each other: for the proof of the dense lineability result $(3)$  (which is proved in Theorem \ref{cyclic and frequently recurrent vector}) we have to use the zero-one law ($(1)\Rightarrow (2)$) together with  the  characterization $(5)$ (which in turn uses $(4)$). 
%We also note that $(2)$ implies the existence of a dense vector space in which every nonzero vector is cyclic and frequently recurrent (i.e. they form a  dense lineable set).

The structure of the paper is the following. In the first section we present some preliminaries.
In Section \ref{section unilateral} we prove the main theorem of the paper, Theorem \ref{main theorem}, in the unilateral case. We will first prove the zero-one law for the case $X=c_0$ in Theorem \ref{0-1 law frequent recurrence c0}.  Although this result is in fact a particular case of Theorem \ref{0-1 law frequent recurrence}, we have chosen to present it separately because its more elementary proof captures the main idea behind the construction of frequently recurrent vectors and motivates the introduction of the $RAT$. This construction  follows a ``self-similar'' process which, in the limit, converges to a frequently recurrent vector. 
%$This theorem is more elementary than the general case, and the construction of frequently recurrent vectors follows a ``self-similar'' process, without an explicit mention of the recurrent arithmetic thickening. 
We introduce the $RAT$ and  extend \cite[Theorem 2.17]{Fur81}   to arbitrary families in Section \ref{section RAT}. We also show how the $RAT$ naturally fits in the context of frequently recurrent operators.
This is applied to show the general case of the zero-one law, i.e. $(1)\Leftrightarrow (2)$ of Theorem \ref{main theorem}, in Section \ref{section 0-1 general case}.  A careful look at the construction of frequently recurrent vectors in the zero-one law, together with some properties of the recurrent arithmetic thickening will allow us to characterize frequently recurrent backward shift operators in Theorem \ref{Characterization frequent recurrence} ($ (4)$ in Theorem \ref{main theorem}). Applying this characterization together with the zero-one law 
we show, in Theorem \ref{thm: caract S_1(A)}, a second characterization for frequently recurrent backward shifts ($ (5)$ in Theorem \ref{main theorem}). In Theorem \ref{cyclic and frequently recurrent vector} we apply all the previous results to obtain dense lineability.
Since our characterization $(4)$ is weaker than the characterization of frequent hypercyclicity from \cite{BonGro18}, it is natural to consider two classes of operators that are formally weaker than frequently hypercyclic but stronger than frequently  recurrent.  In Corollary \ref{Teo: caract freq hyp on S1} and Proposition \ref{prop: condicion con el min} we study them and describe the dynamical notions that they determine.

Finally, in Section \ref{section:bilateral} we prove the  main Theorem \ref{main theorem} in the bilateral setting.
%and in Section \ref{} we generalize Theorem \ref{main theorem} for Furstenberg families satisfying the Separation Lemma \ref{seccion separation lemma}.

\section{Preliminaries and notation}\label{preliminaries}
Through the paper $X$ will denote a Fr\'echet sequence space over $\mathbb K=\mathbb R$ or $\mathbb C$, i.e. $X$ is a subspace of $\mathbb K^{\mathbb N_0}$ or of $\mathbb K^{\mathbb Z}$ such that the functionals $X\ni x=(x_n)_n\mapsto x_n$ are continuous. The topology of $X$ is defined by a sequence of seminorms $(\|\cdot\|_p)_{p\in \mathbb N_0}$. We always assume that for each $x\in X$, the mapping $p\mapsto \|x\|_p$ is non-decreasing. Equivalently, the topology of $X$ may be defined via an $F$-norm, $\|\cdot\|=\sum_{p=0}^\infty \frac{1}{2^{p+1}} \min\{1,\|\cdot\|_p\}$. In that case, $(X,\|\cdot\|)$ results a complete metric space with distance $d(x,y)=\|x-y\|.$ 
%We will denote $D(z,\varepsilon)$ the ball centered at $z$ of radius $\varepsilon$, that is $D(z,\varepsilon)=\{w:\|w-z\|<\varepsilon\}.$ 
Sometimes it will convenient to refer to $[x]_n$ as the $n$-th coordinate of $x$, that is $x_n$. 

An unconditional basis in a Fr\'echet sequence space $X$ is a basis which satisfies that $(a_n\cdot x_n)_n\in X$ whenever $a\in \ell_\infty $ and $x=(x_n)_n\in X$. This implies that for every $p\ge 0$ there are $m_p\geq p$ and a constant $C_p\ge1$ such that 
\begin{equation}\label{incondicionalidad}
    \|ax\|_p\le C_p\|a\|_{\infty} \|x\|_{m_p}.
\end{equation} 
We will always assume that the $C_p$'s are increasing.

Recall that an operator $T$ on $X$ is said to be hypercyclic provided that it has a dense orbit. It is called recurrent if it has a dense set of recurrent vectors, that is, of vectors $x\in X$ such that the set of hitting times $N_T(x,U):=\{n\in\mathbb N: T^n(x)\in U\}$ is infinite for each neighbourhood $U$ of $x.$ 

A set $A\sub \mathbb N_0$ is said to have positive lower density provided that
$$\underline {dens}(A):=\liminf_{n\to \infty} \frac{\# A\cap [0,n]}{n+1}>0$$
and positive upper density if 
$$\overline {dens}(A):=\limsup_{n\to \infty} \frac{\# A\cap [0,n]}{n+1}>0.$$

The lower density is finitely invariant:  $\underline {dens}(A)=\underline {dens}(A\cap F)$ for every cofinite set $F$. The same holds for the upper density.

Given an operator $T$ on $X,$
a vector $x\in X$ is said to be \emph{frequently recurrent} provided that for every open set $U$ around $x$ the set of hitting times $N_T(x,U) = \{n \in \mathbb{N} : T^n(x) \in U\}$ has positive lower density. If the set of frequently recurrent vectors is dense in the space, then we will say that the operator is frequently recurrent. A vector $x\in X$ is called \emph{frequently hypercyclic } if 
$\underline {dens}(N_T(x,U))>0$ for every nonempty open set $U\sub X$. If $T$ has a frequently hypercyclic vector then $T$ is said to be frequently hypercyclic. We refer to the books \cite{BayMat09,GroPer11} for basic concepts on linear dynamics.


Although we will mostly work with frequently recurrent backward shift operators, it will be useful to have a characterization of frequent hypercyclicity at hand. The following theorem was proved by Bonilla and Grosse-Erdmann in \cite{BonGro18}.
\begin{theorem}[Bonilla, Grosse-Erdmann]\label{caracterizacion hiperciclicidad frecuente} Let $X$ be a Fr\'echet sequence space with unconditional basis $(e_n)_{n\geq 0}$. Then $B$ is frequently hypercyclic if and only if there are sets with positive lower density $A_p$ and a sequence of positive nubers $\varepsilon_p\to 0 $ such that
\begin{enumerate}
    \item For every $p\in\mathbb N$, $\sum_{n\in A_p} e_{n+p}\in X$ and
    \item for every $p,q$, $m\in A_q$ and $0\leq j\leq p$
    $$\|\sum_{n\in A_p,n>m} e_{n-m+j}\|_q<\min\{\varepsilon_p,\varepsilon_q\}.$$
    
\end{enumerate}
\end{theorem}
We will always assume that the backward shift $B:X\to X$ is well defined (and hence continuous). Although we will mostly work with unweighted backward shifts all the results can be translated to the  weighted setting via an easy conjugation argument, if $B_w:X\to X$ is a weighted backward shift then we consider $v_n=\prod_{l=1}^n w_l^{-1}$ and $X(v)=\{(x_n)_n:(x_nv_n)_n\in X\}$. It follows that $B:X(v)\to X(v)$ is a well defined backward shift. Moreover, $B:X(v)\to X(v)$ and $B_w:X\to X$ are conjugate via the isomorphism $\phi: X(v)\to X$, $(x_n)_n\to (x_nv_n)_n $. See \cite{GroPer11,bonilla2025zeroEdinburgh} for more details on the matter.

\subsection{A separation Lemma for sets with positive lower density}\label{seccion separation lemma}
We now present a Separation Lemma \ref{separation lemma} for sets with positive lower density, which is crucial for our proofs. The idea of considering sets with separation properties is not new in linear dynamics. For example, a key step to prove that the operators satisfying the frequent hypercyclicity criterion are indeed frequently hypercyclic is a separation lemma of the natural numbers \cite[Lemma 2.2]{BayGri06}. In \cite[Proposition 1]{Bes16} the authors extended this lemma to general families of the natural numbers by proving that if there exist an $\F$-hypercyclic operator, then the family $\F$ must satisfy the separation lemma \cite[Lemma 2.2]{BayGri06}. We finally mention that the proof of the characterization of frequently hypercyclic backward shift operators \cite{BonGro18} uses several times that some sets $A_p$ with positive lower density are separated, where the $A_p$ are sets of hitting times of a frequently hypercyclic vector. 
These ideas have in common that  the separation property obtained is a consequence of a dynamical property: if $x\in X$ is a frequently hypercyclic vector for an operator $T$, and $(y_p)_p\subset X$ is a sequence of pairwise different vectors, then there are $\varepsilon_p>0$ such that the sets $A_p:=\{n\in\mathbb N:\|T^n(x)-y_p\|<\varepsilon_p\}$ satisfy some separation property. However, it is not possible to follow the same strategy for frequent recurrence because in this case the orbit is only known to  approximate a single vector. We have to change the strategy and to look at this property as a consequence of a number theoretic result rather than a dynamical property. To achieve this we will use a result proved recently by Martin, Menet and Puig in \cite{MarMenPui22} (see also \cite{BayGriMatMen24}, where the same lemma is applied).
\begin{lemma}[Separation Lemma]\label{separation lemma}
Let $(A_p)_p$ be a sequence of sets with positive lower density and $(k_p)$ a sequence of natural numbers. Then there exist pairwise disjoint sets of positive lower density $(B_p)_p$, such that for every $p$, $B_p\sub A_p$ and such that for every $n\in B_p$, $m\in B_q$, $n\ne m$, we have that $|n-m|>k_p+k_q$.
\end{lemma}

% \begin{proof}
% Without loss of generality we may suppose that $k_p$ is increasing.

% We will inductively construct sets $B_{1,N}\ldots B_{N,N}$, a decreasing sequence of positive real numbers $(\delta_N)_N$ and a sequence of natural numbers $(n_N)_N$ such that
% \begin{enumerate}[label=\roman*)]
%     \item \label{item contencion A_i} $B_{N,N}\sub A_N$;
%     \item for every $1\leq i<j \leq N$, we have that $B_{i,j}\sub B_{i,j-1}$;
%     \item \label{ item interseccion no vacia}for every $1\leq i<N$, $B_{i,N}\cap [0,N]=B_{i,N-1}\cap[0,N];$
%     \item \label{densidad B_i,N}for every $1\leq i\le j\leq N$  and every $n>n_i$
%     $$\frac{\# B_{i,j}\cap [0,n]}{n+1}\geq \delta_i-\sum_{l=1}^j\frac{\delta_i}{2^{j+1}}.$$
%     In particular, for every $1\leq i\le j\leq N$ we have that 
%     $\underline{dens}(B_{i,j})\geq\delta_i-\sum_{l=1}^j\frac{\delta_i}{2^{j+1}}>\frac{\delta_i}{2}$;
%     \item \label{item pairwise disjoint}the sets $B_{1,N}\ldots B_{N,N}$ are pairwise disjoint and
%     \item \label{item separados}for every $1\leq i,i'\leq N$ we have that for every $n\in B_{i,N}$ and $m\in B_{i',N}$ with $n\neq m$ then
%     $$|n-m|> k_i+k_{i'}.$$

% \end{enumerate}

% For the first step we will construct the sets $B_{1,1}$, $B_{1,2}$ and $B_{2,2}$ and $0<\delta_1<\delta_2$ and $n_1<n_2\in\mathbb N$ satisfying i)-vi). 

% Let $B_{1,1}$ be any subset of $A_1$ with positive lower density such that for any $n,m\in A_1$, $n\neq m$, then $|n-m|\geq 2k_1$. Let $\delta_1<\underline{dens}(B_{1,1})$ and $n_1\in\zN$ such that for every $n\geq n_1,$ $\frac{\# B_{1,1}\cap [0,n]}{n+1}\geq \delta_1$. 


% Let $B_{2,2}\sub A_2$ be any subset with positive lower density  such that:
% \begin{itemize}
% \item[a)] for every $n\in \zN$ $\frac{\# B_{2,2}\cap[0,n]}{n+1}\leq \frac{\delta_1}{32k_2}$; 
% \item[b)] $\min \{B_{2,2}\}>\frac{8(2+2k_2)}{\delta_1}+2k_2$ and such that
% \item [c)]for any $n,m\in B_{2,2}$, $n\neq m$, we have that $|n-m|> 2k_2$.
% \end{itemize}
% To justify the existence of this set we may consider for big enough $l$,  $\{a_{ln}:n\in \mathbb N\}\cap[2k_2+n_1+1,\infty)$, where $a_{n}$ is the sequence defining $A_2$. This  set has positive lower density $\frac{\underline{dens}(A_2)}{l}$ and upper density bounded by $\frac{1}{l}$. So, for big enough $l$, the set has upper density less than $\frac{\delta_1}{32k_2}$. By removing finitely many numbers we obtain a).




% Let $0<\delta_2<\underline{dens}(B_{2,2})$ and $n_2\in\zN$ such that for every $n\geq n_2$ we have that $\frac{\# B_{2,2}\cap[0,n]}{n+1}\geq\delta_2$. In particular,
% $$
% \delta_2<\underline{dens}(B_{2,2})\leq\overline{dens}(B_{2,2})\leq \frac{\delta_1}{32k_2}<\delta_1.$$
% If $B_{2,2}=(b_n)_n$ consider $\tilde B_{2}:= \bigcup_{n} [b_n-2k_2,b_n+2k_2]=B_{2,2}+[-2k_2,2k_2]$. It follows by a) that for every $n\ge \min\{B_{2,2}\}-2k_2\ge \frac{8(2+2k_2)}{\delta_1}$ we have,
% $$
% \frac{\# \tilde B_{2}\cap[0,n]}{n+1}\leq \frac{4k_2\#  B_{2,2}\cap[0,n]}{n+1}+\frac{2k_2}{n+1}\leq \frac{\delta_1}{8}+\frac{\delta_1}{8}=\frac{\delta_1}{4},
% $$  
% and for $n< \min\{B_{2,2}\}-2k_2$ we have, $\tilde B_{2}\cap[0,n]=\emptyset$. 
% This implies that  $B_{1,2}:=B_{1,1}\setminus \tilde B_{2}\sub B_{1,1}$ has lower density greater than $\delta_1-\frac{\delta_1}{4}$. Moreover, for every $n\geq n_1$,
% $$\frac{\# B_{1,2}\cap [0,n]}{n+1}\geq \delta_1-\frac{\delta_1}{4}.$$

% Note that by construction we have that if $n\in B_{1,2}$ and $m\in B_{2,2}$ then $|n-m|>2k_2>k_1+k_2$ and that $B_{1,2}$ and $B_{2,2}$ are disjoint. Also b) implies that $B_{1,2}\cap[0,2]=(B_{1,1}\setminus \tilde B_{2}) \cap[0,2]=B_{1,1}\cap[0,2]$, because $\tilde B_{2}$ is contained in $(2,\infty)$.


% Suppose that we have constructed $B_{i,j}$, $1\leq i\leq j\leq N$ and $\delta_{i},n_i$, $1\leq i\leq N$ satisfying properties i)-vi).

%  Let $B_{N+1,N+1}$ be any subset of $A_{N+1}$ with positive lower density
%  such that 

% \begin{itemize}
% \item[a)] for every $n\in \zN$ $\frac{\# B_{N+1,N+1}\cap[0,n]}{n+1}\leq \frac{\delta_N}{2^{N+2}k_{N+1}}$;
% \item[b)] $\min \{B_{N+1,N+1}\}\geq \frac{2^{N+2}(2+2k_{N+1})}{\delta_N}+2k_{N+1}$ and such that
% \item [c)]for any $n,m\in B_{N+1,N+1}$, $n\neq m$, we have that $|n-m|> 2k_{N+1}$.
% \end{itemize}
  
%  Let $0<\delta_{N+1}<\underline{dens}(B_{N+1,N+1})$ and $n_N$ such that for every $n\geq n_N$ we have that
%   $\frac{\# B_{N+1,N+1}\cap[0,n]}{n+1}\geq\delta_{N+1}$.
 
%  If $B_{N+1,N+1}=(b_n)_n$ consider $\tilde B_{N+1}=\bigcup_n [b_n-2k_{N+1},b_n+2k_{N+1}]$. Notice that for every $n\ge \min\{B_{N+1,N+1}\}-2k_{N+1}\ge \frac{2^{N+2}(2+2k_{N+1})}{\delta_N}$ we have,
%  \begin{equation}\label{densidad superior chica}
%      \frac{\# \tilde B_{N+1}\cap[0,n]}{n+1}\leq \frac{4k_{N+1}\#  B_{N+1,N+1}\cap[0,n]}{n+1}+\frac{2k_N}{n+1} \leq\frac{\delta_N}{2^{N+1}}.
%  \end{equation} 
%  On the other hand, if $n< \min\{B_{N+1,N+1}\}-2k_{N+1}$ then $\tilde B_{N+1}\cap[0,n]$. Hence, for any $n\in\mathbb N,$ we have $\# \tilde B_{N+1}\cap[0,n]\le (n+1)\frac{\delta_N}{2^{N+1}}.$

 
%  Consider now, for every $1\leq i\leq N$, the sets $B_{i,{N+1}}=B_{i,N}\setminus \tilde B_{N+1}$.
 
%  Let us show that the sets $B_{1,N+1},\ldots, B_{N+1,N+1}$ satisfy conditions i)-vi). Clearly they satisfy conditions i) and ii).
 
% \ref{ item interseccion no vacia}.
%  Observe that by b), $\tilde B_{N+1}$ is contained in $(N+1,\infty)$ and hence for every $1\leq i<N+1$, $B_{i,N+1}\cap [0,N+1]=(B_{i,N}\setminus \tilde B_{N+1})\cap [0,N+1]=B_{i,N}\cap [0,N+1].$
 
%  \ref{densidad B_i,N}.
%  By \eqref{densidad superior chica} and the inductive hypothesis  \ref{densidad B_i,N} applied to $B_{i,N}$  have that for every $i$ and every $n\geq n_i$
%   $$\frac{\# B_{i,N+1}\cap[0,n]}{n+1}\geq \delta_i-\sum_{l=1}^N \frac{\delta_i}{2^{l+1}}-\frac{\delta_N}{2^{N+1}}\geq \delta_i-\sum_{l=1}^{N+1} \frac{\delta_i}{2^{l+1}}.$$

% \ref{item pairwise disjoint} Since $B_{N+1,N+1}$ is contained in $\tilde B_{N+1}$ and $\tilde B_{N+1}$ is disjoint with every $B_{i,N+1}$ with $1\leq i\leq N$ it follows that $B_{N+1,N+1}$ is disjoint with every $B_{i,N+1}$ with $1\leq i\leq N$. On the other hand, every $B_{i,N+1}$ with $1\leq i\leq N$ is contained in $B_{i,N}$, and the $B_{i,N}$ with $1\leq i\leq N$ are pairwise disjoint by the inductive hypothesis \ref{item pairwise disjoint}. We conclude that the $B_{i,N+1}$ with $1\leq i\leq N+1$ are pairwise disjoint. 

% \ref{item separados}.
%  By construction we have that if $n\in B_{i,{N+1}}$ with $i<N+1$ and $m\in B_{N+1,N+1}$ then $|n-m|> 2k_{N+1}\geq k_i+k_{N+1}$. On the other hand, since each $B_{i,N+1}$ is contained in $B_{i,N}$ we have by the inductive hypothesis \ref{item separados} that for every $1\leq i,i'<N+1$ and every $n\in B_{i,N+1}$ and $m\in B_{i',N+1}$
%  that $|n-m|> k_i+k_{i'}$. 
  
%  This shows the construction of the sets $B_{1,N+1}\ldots B_{N+1,N+1}$.
  
%   Let for each $i$, $B_i=\bigcap_{j\geq i} B_{i,j}$. Note that by ii), the sequence $(B_{i,j})_{j\ge i}$ is decreasing  and by \ref{item contencion A_i} $B_{i}\sub B_{i,i}\sub A_i
% $. 
% Also the $B_i$ are pairwise disjoint. Indeed, for every $i,i'$ we have that $B_i\sub B_{i,i+i'}$ and $B_{i'}\sub B_{i',i+i'}$. Since $B_{i,i+i'}$ is disjoint with   $B_{i',i+i'}$ we get  that $B_i$ is disjoint with $B_{i'}$.

% We show now that each $B_i$ has positive lower density. This implies in particular that each $B_i$ is nonempty. Let $N\geq n_i$, then by \ref{ item interseccion no vacia} $B_i\cap [0,N]=B_{i,N}\cap[0,N]$ and therefore
% $$\frac{\#B_i\cap [0,N]}{N+1}=\frac{\#B_{i,N}\cap [0,N]}{N+1}\geq \delta_i-\sum_{l=1}^N\frac{\delta_i}{2^{l+1}}.$$
%  Thus,
%   $$\underline{dens}(B_i)\geq \delta_i-\sum_{l=1}^\infty\frac{\delta_i}{2^{l+1}}=\frac{\delta_i}{2}.$$ 
%     Finally, 
%   if $n\in B_{i}$ and $m\in B_{i'}$ with $n\neq m$, we have that $n\in B_{i,i+i'}$ and $m\in B_{i',i+i'}$ which yields $|n-m|> k_i+k_{i'}$.  
% \end{proof}


\section{Frequently recurrent backward shifts: the unilateral case}\label{section unilateral}
In this section we prove the unilateral case of the main Theorem \ref{main theorem} which is an immediate consequence of Theorem \ref{0-1 law frequent recurrence} (the zero-one law), Theorem \ref{Characterization frequent recurrence} (First characterization of frequent recurrence) and Theorem \ref{thm: caract S_1(A)} (Second characterization of frequent  recurrence).
As an application we prove in Subsection \ref{subsection segunda aplicacion} that the set of frequently recurrent vectors of a frequently recurrent backward shift is dense lineable.
An important ingredient of the construction of frequently recurrent vectors is the recurrence arithmetic thickening, which is addressed in Subsection \ref{section RAT}.


\subsection{A zero-one law in ${c_0}$}\label{subsection primer caracterizacion}
We begin this section by proving a zero-one law for frequently recurrent backward shift operators: the existence of one nonzero frequently recurrent vector implies that there is a dense set of such vectors. Then, by studying the construction of this dense set of vectors  we can describe the necessary conditions to obtain frequent recurrence, and we show hence a characterization of frequently recurrent shifts in Theorem \ref{Characterization frequent recurrence}.


We begin with the case of $c_0$, which is  technically less demanding and which illustrates the main underlying ideas for treating the general case. Moreover, the space $c_0$ has long served as a rich source of counter-examples in linear dynamics, for instance in \cite{BayGri07} the authors showed that frequent hypercyclicity  does not imply chaos, topological mixing or the existence of invariant Gaussian measures, by providing suitable examples of backward shifts on $c_0$. 

We will consider the backward shift operator on the weighted $c_0(v)$ space. Here $v=(v_n)_{n\ge 0}$ is a sequence of positive weights and for $x\in c_0(v)$ the norm is given by $\|x\|_{c_0(v)}=\sup_n|v_nx_n|$. We will assume, as we do in the whole manuscript, that $B$ is bounded, which in this case is equivalent to the condition of $(v_n/v_{n+1})_n$ being  bounded.

Our goal is to construct plenty of recurrent vectors starting from a single one. The construction of (frequently) recurrent vectors is in some sense fractal: We start by an initial vector $x_0$, and then we will  paste infinitely many copies of $x_0$ to obtain a vector $x_1\sim x_0$ and  such that the orbit of $x_1$ approximates $x_0$ along some set $A_1$ of positive lower density. In the $k$-step we will paste infinitely many copies of the vector $x_{k-1}$ and the orbit of $x_k$ will approximate each $x_{k-1}$ along  $A_k$. The vector $x=\lim_{k\to \infty} x_k$ will be our candidate to be a frequently recurrent vector. 

\begin{theorem}[Zero-one law for frequently recurrent vectors for $c_0$]\label{0-1 law frequent recurrence c0}
Let $X=c_0(v)$. If $B$ supports a nontrivial frequently recurrent vector then $B$ is frequently recurrent.
\end{theorem}
\begin{proof}
It suffices to prove that for every $R\in\mathbb N_0$, and every $\varepsilon>0$ there is a frequently recurrent vector $y\in X$ such that $\|y-e_R\|_{c_0(v)}<\varepsilon$. Indeed, if $U$ is a nonempty open set, let $z=\sum_{j=0}^R z_je_j\in U$. 
If $y\in X$ is a frequently recurrent vector sufficiently close to $e_R$, the continuity of $B$, of the sum and of the product by scalars imply that the vector $\sum_{j=0}^R z_j B^j(y)$ is a frequently recurrent vector in $U$.

So,  let $x\neq 0$ be a frequently recurrent vector,  $R\in\mathbb N_0$ and $\varepsilon>0$. Without loss of generality, we may suppose that $x_R\neq 0$ (otherwise consider $B^n(x)$ which is frequently recurrent). Let $\delta\in(0,1)$ such that $|x_R|>\delta.$


Since $x$ is a recurrent vector,  there is for every $p\in\mathbb N_0$, a  $q_p>p+R$ such that $|x_{q_p}|>\delta$. By the continuity of the coordinate functionals, there are a decreasing sequence of positive numbers $(\varepsilon_p)_p$ and a sequence of natural numbers $(k_p)_p$ such that for each $p\in\mathbb N_0$, $\varepsilon_{p}<\frac{\varepsilon}{2^{p+1}}$,  $k_p>p+1$ and
\begin{align}\label{eq: beta_p (zero one freq) UNIL c_0}
 \begin{split}
  \|w\|_{c_0(v)}<\varepsilon_{p}\Longrightarrow  |w_{j}|<\frac{\delta}{4} ,  \text{ for every } 0\leq j\leq q_p, &\quad \text{and}\\
 \|x-\sum_{0\leq j\le k_p} x_j e_j\|_{{c_0(v)}}<\frac{\varepsilon}{2^{p+1}}\frac{\delta}{8}.&
 \end{split}
 \end{align}

We will now present sets $A_p$  that will allow us to construct a frequently recurrent vector near $e_R$. For each $p\in\mathbb N_0$, $A_p$ will be a set of positive lower  density such that  
\begin{equation}\label{eq: A_p zero one freq c_0}
A_p\sub N_B(x,U_p), \text{ where } U_p=\{y\in X:\|y-x\|_{c_0(v)}<\varepsilon_p\frac{\delta}{8}\}.
\end{equation}
Applying the Separation Lemma \ref{separation lemma}
we may assume that the $A_p$'s are pairwise disjoint and that for every $n\in A_p,m\in A_q$ with $n\neq m$ we have that
\begin{equation}\label{eq separation zero one freq c_0}
|n-m|>k_p+k_q+2R>p+1+q+1+2R.    
\end{equation}

Notice that for every $p\in\zN_0$ and every $0\leq j\leq p$ we have that   $\sum_{n\in A_p} e_{n+R+j}$ converges. Indeed, by \eqref{eq: beta_p (zero one freq) UNIL c_0} and \eqref{eq: A_p zero one freq c_0} we have that $|x_{n+q_p}-x_{q_p}|<\frac{\delta}{4}$ and hence $|x_{n+q_p}|>\frac{\delta}{2}$ whenever $n\in A_p$. By uncondionality this implies that $\sum_{n\in A_p} e_{n+q_p}$ converges and by applying the backward shift $q_p-(R+j)$ times it follows that $\sum_{n\in A_p} e_{n+R+j}$ converges.  By the finite invariance property of  the sets of positive lower density, we may assume that 
\begin{equation}\label{eq A_p empieza en p c_0}
    A_p\sub [p+R+1,\infty)
\end{equation} and that,  for every $0\le j\le p,$
\begin{equation}\label{eq: sumas de e_j norma chica (zero one freq) c_0}
 \|\sum_{n\in A_p} e_{n+R+j}\|_{c_0(v)}<\frac{\varepsilon}{2^{p+1}}.   
\end{equation}
We will now proceed to construct a frequently recurrent vector near $e_R$ by induction. The rough idea is to construct first a vector $y^1$ near $e_R$ whose orbit is frequently near $y^0:=e_R$, and by induction a sequence of vectors $(y^k)_k$, contained in a small neighbourhood of $e_R$, and such that the orbit of $y^k$ is frequently near $y^{k-1}$. Those vectors will be constructed so that they   form a convergent sequence to a vector $y$, which will be a frequently recurrent vector near $e_R.$ Each vector $y^k$ will consist on carefully pasting together shifted copies of a truncation of the vector $y^{k-1}$ (and hence of copies of $e_R$).
 
Set $N_0=M_0=0$ and let $y^1=e_R+\sum_{n\in A_0}
e_{n+R}=e_R+\sum_{n\in A_0}
e_{n+N_0+R}$.  The vector is well defined, because of \eqref{eq: sumas de e_j norma chica (zero one freq) c_0}. Notice also that
$$\|y^1-e_R\|_{c_0(v)}=\|\sum_{n\in A_0}
e_{n+R}\|_{c_0(v)}<\frac{\varepsilon}{2}.$$

For the second step, let $N_1\in\zN$ such that $N_1+R$ is the second nonzero coordinate of $y^1$, which in this step coincides with $\min A_0$. Let $M_1>N_1$  such that $ \frac{\delta}{8}\|B\|^{N_{1}}\varepsilon_{M_{1}}<\frac{\varepsilon_0}{2^{2}}$.

We will add to $y^1$ infinitely many copies of a truncation $P_{N_1+R}(y^1)$ of $y^1$, where $P_{N_1+R}$ denotes the projection into the first $N_1+R$ coordinates:

\begin{align*}
  y^2&=y^1+\sum_{n\in A_{M_1}}\sum_{j=0}^{N_1+R} y^1_j e_{n+j} =y^1+\sum_{n\in A_{M_1}}e_{n+N_0+R}+e_{n+N_1+R}  =e_R+\sum_{p=0}^{1}\sum_{n\in A_{M_p}}\sum_{j=0}^p e_{n+N_j+R}.
  \end{align*}
%where we adopted the notation $A_{M_0}=A_0.$

Observe that, since $|n-m|>k_p+k_q+2R$ whenever $n\in A_p,m\in A_q$ and $k_p>p$, the sums involved have disjoint support. Also, since we imposed in \eqref{eq A_p empieza en p c_0} that $n>p+R$ for every $n\in A_p$ and $M_1>N_1$ we have that the first two nonzero coordinates of $y^2$ are $R$ and $N_1+R$.  It follows by \eqref{eq: sumas de e_j norma chica (zero one freq) c_0} that the sums are well defined and that 
$$\|y^2-y^1\|_{c_0(v)}=\|\sum_{n\in A_{M_1}}\sum_{j=0}^1 e_{n+N_j+R}\|_{c_0(v)}\leq  \frac{\varepsilon}{2^{M_1+1}}<\frac{\varepsilon}{2^2}.$$

Let us briefly sketch the third step. Let $N_2\in\zN$ such that $N_2+R$ is the third nonzero coordinate of $y^2$ (In this case $N_2=\min\{n\in A_{M_0}\cup A_{M_1}\,:\, n>N_1\}.$ Let $M_2>N_2$ such that $ \frac{\delta}{8}\|B\|^{N_{2}}\varepsilon_{M_{2}}<\frac{\varepsilon_0}{2^{3}}$.

% Let i.e. $N_2$ is either the first element of $A_{N_1}$ or the second element of $A_0=A_{N_0}$.
% Note that $N_2+R$ is the index of the third nonzero  coordinate of $y^2$.

Let us define $y^3$, with three copies of $e_R$ starting at each element of $n\in A_{M_2}$:
\begin{align*}
y^3&=y^2+\sum_{n\in A_{M_2}}\sum_{j=0}^{N_2+R} y^2_j e_{n+j}    =y^2+\sum_{n\in A_{M_2}} \sum_{j=0}^2  e_{n+N_j+R} =e_R+\sum_{p=0}^{2}\sum_{n\in A_{M_p}}\sum_{j=0}^p e_{n+N_j+R}.
\end{align*}
Then, proceeding as before, $y^3$ will be well defined and will satisfy, by \eqref{eq: sumas de e_j norma chica (zero one freq) c_0}, that $\|y^3-y^2\|_{c_0(v)}<\frac{\varepsilon}{2^3}$ and that its first three nonzero coordinates are $R,N_1+R$ and $N_2+R$.

We point out that there is a subtle difference from the fourth step onwards. Again $N_3\in\zN$ is defined so that $N_3+R$ is the fourth nonzero coordinate of $y^3$, but in this case  $N_3$ is not necessarily $\min\{n\in A_{M_0}\cup A_{M_1}\cup A_{M_2}\,:\, n>N_2\}.$ This is because it may happen that $N_1+N_2+R$ is the fourth nonzero coordinate of $y^3$, and $N_1+N_2\notin A_{M_0}\cup A_{M_1}\cup A_{M_2}$. In any case, since the nonzero coordinates of $y^3$ are of the form $n+N_j+R$, with $n\in A_{M_p}$, $1\leq p\leq 2$ and $0\leq j\leq p$, we have that 
$$N_3=\min(N_2,\infty)\bigcap\left(\bigcup_{p=0}^2\bigcup _{j=0}^p A_{M_p}+N_j\right).$$
The definition of $y^4$ and $M_3$ is analogous.



Let us now proceed to the inductive step. Suppose that we have defined $N_1,\dots, N_{k-1}$, $M_1,\dots,M_{k-1}$ and $y^1,\dots,y^k$ in the same way and that for every  $1\leq j\le k,$
$$
\|y^{j}-y^{j-1}\|_{c_0(v)}<\frac{\varepsilon}{2^{j}}.
$$
Define $N_k$ such that $N_k+R$ is the $(k+1)$-th nonzero coordinate of $y^k$  and let $M_k>N_k$ such that 
\begin{equation}\label{eq: definicion M_k c_0}
    \frac{\delta}{8}\|B\|^{N_{k}} \varepsilon_{M_{k}}<\frac{\varepsilon_0}{2^{k+1}}.
\end{equation}

% :=\min\{n>N_{k-1}:n\in  \bigcup_{j=0}^{k-1} A_{N_j}\}$. Note that  
% \begin{align}\label{eq: indices of the first copies}
%     N_0+R,N_1+R, \dots,N_{k}+R
% \end{align} are the indices of the first $k+1$ nonzero coordinates of $y^k$. 

% Note also, for later use, that with this definition,
% \begin{align}\label{eq: A es la union de los N_k}
%     \{0\}\cup\bigcup_{j \in \mathbb N_0} A_{N_j}=\bigcup_{j \in \mathbb N_0}\{N_j\}.
% \end{align}

Let
\begin{align}\label{eq: def y^{k+1}}
y^{k+1}=y^k+\sum_{n\in A_{M_{k}}}\sum_{j=0}^{N_{k}+R} y^k_j e_{n+j} =y^k+\sum_{n\in A_{M_{k}}}\sum_{j=0}^{k}  e_{n+N_j+R}=e_R+\sum_{p=0}^{k}\sum_{n\in A_{M_p}}\sum_{j=0}^p e_{n+N_j+R}.    
\end{align}
Note that by \eqref{eq: sumas de e_j norma chica (zero one freq) c_0},
\begin{align*}
\|y^{k+1}-y^k\|_{c_0(v)} &=\|\sum_{n\in A_{M_{k}}}\sum_{j=0}^{k}e_{n+N_j+R}\|_{c_0(v)}
\leq \frac{\varepsilon}{2^{M_{k}+1}}<\frac{\varepsilon}{2^{k+1}}.    
\end{align*}
This shows that $(y^k)_k$ is a Cauchy sequence so that the limit vector exists and that

\begin{align*}
y:=\lim_{k\to \infty} y^k &=e_R+\sum_{p=0}^{\infty}\sum_{n\in A_{M_p}}\sum_{j=0}^pe_{n+N_j+R}=e_R+\sum_{p=0}^\infty\sum_{n\in A_{M_p}}\sum_{j=0}^{N_p+R}y_je_{n+j}
\end{align*}
is well defined and that $\|y-e_R\|_{c_0(v)}<\varepsilon.$
Note that by the construction of $y^k$ we have that
% by \eqref{eq: indices of the first copies c_0}, we have
\begin{align*}
y=\sum_{j=0}^\infty e_{N_j+R}.
\end{align*}


It remains to prove that $y$ is a frequently recurrent vector. From now and until the end of the proof, we fix $q\in\zN$ and $m\in A_{M_{q}}$. Recall that by \eqref{eq separation zero one freq c_0}, since $k_p>p$ for every $p,$ we have that for every $p\in\mathbb N$ and $n\in A_{M_p}$, $|n-m|>M_p+M_{q}+2R>N_p+N_{q}+2R$. Thus, if $n<m$ and $n\in A_{M_{p}}$, then $n-m+N_j+N_q+R<0$ for every $0\le j\le p$. Hence,
$$B^m\left(\sum_{n\in A_{M_p},n<m}\sum_{j=0}^pe_{n+N_j+R}\right)=\sum_{n\in A_{M_p},n<m}\sum_{j=0}^pe_{n-m+N_j+R}=0.$$
Noting that $m$ belongs only to $A_{M_{q}}$, because the $A_{M_p}$'s are pairwise disjoint, we get that
\begin{align*}
B^m(y)-\sum_{j=0}^{N_{q}+R}y_je_j &=B^m(y)-\sum_{j=0}^{q}e_{N_j+R}\\
&=\sum_{p\neq q}\sum_{n\in A_{M_p},n>m}\sum_{j=0}^{p}e_{n-m+N_j+R}+\sum_{n\in A_{M_{q}},n\geq m}\sum_{j=0}^{q}e_{n-m+N_j+R}-\sum_{j=0}^qe_{N_j+R}\\
&=\sum_{p=0}^\infty \sum_{n\in A_{M_p},n>m}\sum_{j=0}^pe_{n-m+N_j+R}\\
\end{align*}
Thus,
\begin{equation}\label{eq norma frec zero one c_0}
 \|B^m(y)-\sum_{j=0}^{N_{q}+R}y_je_j \|_{c_0(v)}= \left\|\sum_{p=0}^\infty \sum_{n\in A_{M_p},n>m}\sum_{j=0}^pe_{n-m+N_j+R}\right\|_{c_0(v)}.
\end{equation}


We claim that for every $N_k$ we have that $\|B^{N_k}x-x\|_{c_0(v)}<\varepsilon_0-\frac{\varepsilon_0}{2^k}\leq\varepsilon_0$.
Indeed, for $N_0=0$ it is plain and for $N_1=\min A_0$ it follows because $A_0\subseteq \{n\in \zN: \|B^n(x)-x\|_{c_0(v)}< \frac{\varepsilon_0 \delta}{8}\}$ and $\delta<1$.

Let $k>1$. By the construction of $N_k$, we have that $N_k+R$ is the  $(k+1)$-th non zero coordinate of $$y^{k}=e_R+\sum_{p=0}^{k-1}\sum_{n\in A_{M_p}} \sum_{j=0}^p e_{n+N_j+R}.$$ 
By a quick inspection we have that $N_k$ must be a number of the form $n+N_j$, where $n\in A_{M_p}$ for some $0\leq p\leq {k-1}$ and some $0\leq j\leq p$. Hence by \eqref{eq: definicion M_k c_0} and the inductive hypothesis,
\begin{align}
\label{eq: A contenido en N(x,epsilon)} \|B^{N_k}x-x\|_{c_0(v)} & =\|B^{n+N_j}x-x\|_{c_0(v)}\leq \|B^n B^{N_j}x-B^{N_j}x\|_{c_0(v)}+\|B^{N_j}x-x\|_{c_0(v)}    \\
\nonumber &\le \|B\|^{N_j}\,\|B^{n}x-x\|_{c_0(v)}+ \varepsilon_0(1-\frac1{2^j}) < \|B\|^{N_j-N_p}\,\frac{\varepsilon_0}{2^{p+1}}+  \varepsilon_0(1-\frac1{2^{j}}) < \varepsilon_0(1-\frac1{2^k}).
\end{align}



This implies that for every $j\in \zN_0$, $\|B^{N_j}x-x\|_{{c_0(v)}}<\varepsilon_0$ and since $|x_R|>\delta$, we have by \eqref{eq: beta_p (zero one freq) UNIL c_0} that $|x_{N_j+R}|>\frac{\delta}{2}$. On the other hand, we have that for every $n\in A_{M_p},$ $\|B^{n}x-x\|_{{c_0(v)}}<\varepsilon_{M_p}$ and hence by \eqref{eq: beta_p (zero one freq) UNIL c_0} and the fact that $q_{M_p}>M_p+R>N_j+R,$ we obtain for $0\leq j\leq p$ and $n\in A_{M_p}$ that
$|x_{n+N_j+R}-x_{N_j+R}|<\frac{\delta}{4}$. This implies that for every $n\in A_{M_p}$ and $0\leq j\leq p$, $|x_{n+N_j+R}|>\frac{\delta}{4}.$
 
 Since $m\in A_{M_q}$ and by \eqref{eq separation zero one freq c_0}, $n-m>M_p+k_{M_{q}}+R>k_{N_{q}}$ whenever $n\in A_{M_p}$ and $n>m$,  we have  $[\sum_{l=0}^{k_{N_{q}}}x_le_l]_{n-m+N_j+R}=0$ for each $0\le j\le p$. Hence it follows that  
 % \begin{align}\label{cota B^my, n>m}
 % \|B^m(y)-\sum_{j=1}^{N_{s_0}+R}y_je_j \|&\leq CR\|B\|^R\|z\|_\infty\left\|\sum_{s=0}^\infty \sum_{n\in A_{N_s},n>m}\sum_{l=0}^se_{n-m+N_l+R}\right\|\\
 % &\leq \frac{4C^2}{\delta}R\|B\|^R\|z\|_\infty\left\|\sum_{s=1}^{\infty}\sum_{n\in A_{N_s},n>m}\sum_{l=0}^sx_{n+N_l+R}e_{n-m+N_l+R}\right\|\\
 % &=\frac{4C^2}{\delta}R\|B\|^R\|z\|_\infty\left\|\sum_{s=0}^{\infty}\sum_{n\in A_{N_s},n>m}\sum_{l=0}^s[B^m(x)-\sum_{j=1}^{k_{N_{s_0}}}x_j]_{n-m+N_l+R}e_{n-m+N_l+R}\right\|\\
 % &\leq \frac{4C^3}{\delta}R\|B\|^R\|z\|_\infty\left(\|B^m(x)-x\|+\|x-\sum_{j=1}^{k_{N_{s_0}}}x_j\|\right)< \varepsilon_{N_{s_0}}.
 % \end{align}
  \begin{align*}%\label{cota B^my, n>m c_0}
 \begin{split}
 \Bigg\|\sum_{p=0}^\infty \sum_{n\in A_{M_p},n>m}\sum_{j=0}^p & 
  e_{n-m+N_j+R} \Bigg\|_{c_0(v)}
 =\Bigg\|\sum_{p=0}^{\infty}\sum_{n\in A_{M_p},n>m}\sum_{j=0}^p\frac{1}{x_{n+N_j+R}}\cdot{x_{n+N_j+R}}e_{n-m+N_j+R}\Bigg\|_{c_0(v)}\\
 &=\Bigg\|\sum_{p=0}^{\infty}\sum_{n\in A_{M_p},n>m}\sum_{j=0}^p\frac{1}{x_{n+N_j+R}}\cdot[B^m(x)-\sum_{l=0}^{k_{N_{q}}}x_le_l]_{n-m+N_j+R}e_{n-m+N_j+R}\Bigg\|_{c_0(v)}\\
 &\leq \frac{4}{\delta}\left(\|B^m(x)-x\|_{{c_0(v)}}+\|x-\sum_{l=0}^{k_{N_{q}}}x_le_l\|_{{c_0(v)}}\right).
  \end{split}
 \end{align*}

Therefore, using  \eqref{eq norma frec zero one c_0}  together with \eqref{eq: beta_p (zero one freq) UNIL c_0} and the definition of $A_{M_p}$ in \eqref{eq: A_p zero one freq c_0}, we conclude that for any $m\in A_{M_{q}},$
\begin{align}\label{eq: B^my final c_0}
\|B^m(y)-\sum_{j=0}^{N_{q}+R}y_je_j \|_{c_0(v)}\le \frac{4}{\delta}\left(\|B^m(x)-x\|_{{c_0(v)}}+\|x-\sum_{l=0}^{k_{N_{q}}}x_le_l\|_{{c_0(v)}}\right)< \frac{\varepsilon}{2^{N_{q}+1}}.
    \end{align}

 
Finally let $V$ be an open set around $y$. Then since  ${N_q}\to\infty$ as $q\to\infty,$    $\sum_{j=0}^{N_q+R} y_je_j\to y$ and hence,  for big enough  $q\in\mathbb N$, it follows by \eqref{eq: B^my final c_0} that $A_{M_{q}}\sub N_B(y,V).$
This shows that $y$ is a frequently recurrent vector.
\end{proof}


\subsection{The recurrent arithmetic thickening\label{section RAT}}
A careful look at the proof of Theorem \ref{0-1 law frequent recurrence c0} shows that the sequences $(N_k)_k$ and $(A_{M_p})_p$ constructed there have an arithmetical relation, specifically, 
$$
N_{k}=\min\{ n>N_{k-1}\,:\, n\in\bigcup_{p=0}^{k-1}\bigcup_{j=0}^p A_{M_p}+N_j\}.
$$
This can be seen from equation \eqref{eq: def y^{k+1}} and recalling that $N_{k}$ is defined as the $(k+1)$-th nonzero coordinate of $y^k$. 
Moreover, if $A:=\{N_k:k\in\zN\}=\bigcup_{p=0}^{\infty}\bigcup_{j=0}^p A_{M_p}+N_j$, $A$ is contained in the set of return times of a small neighborhood of the recurrent vector, see \eqref{eq: A contenido en N(x,epsilon)}.

This at first sight strange condition should not be entirely surprising, since it is well-known that the sets of return times of recurrent vectors always contain some arithmetical structure. Indeed, it is shown in \cite[Theorem 2.17]{Fur81} that the sets of return times of recurrent points always contain $IP$-sets. A key point to achieve our goals is not only to obtain arithmetic structure, but to keep the  positive density of the set. 

Inspired by these comments we will define the recurrent arithmetic thickening of a sequence of subsets of $\zN$. 
\begin{definition}
    Given a sequence of subsets $A_p\sub \zN$ we define the recurrent arithmetic thickening of $(A_p)_p$ as the set
    \begin{equation}\label{RAT(A_p)}
    RAT((A_p)_p):=\{0\}\cup\bigcup_{p=0}^\infty\bigcup_{j=0}^p A_{p}+N_j ,\, \text{ where, }N_k=\min\{ n>N_{k-1}\,:\, n\in\bigcup_{p=0}^{k-1}\bigcup_{j=0}^p A_{p}+N_j\}
       \end{equation}
     for $k\ge1$ and $N_0=0$.
     \end{definition}
It is clear that, for any sequence $(A_p)_p$, it holds $(N_j)_j\sub RAT((A_p)_p)$. We will say that the sequence $(A_p)_p$ is \emph{compatible} if $RAT((A_p)_p)=(N_j)_j$.


% To this aim, we need first to study arithmetic properties of sets $A_p$ being contained in the sets of hitting times of a frequently recurrent vector. 

% Let $A_p$ have positive lower density, what arithmetic condition must they satisfy to imply the existence of a frequently recurrent vector $x$ and an operator $T$ such that $A_p\subseteq N(x,U_p)$, where $U_p$ is a decreasing chain of open sets around $x$?



\begin{remark}\label{rem:A_0 contenido en RAT}
    For any sequence of subsets $(A_p)_p$, let $N_k$ be the numbers appearing in the definition of $RAT((A_p)_p)$, i.e. $N_{k+1}=\min\{ n>N_{k}\,:\, n\in\bigcup_{p=0}^{k}\bigcup_{j=0}^p A_{p}+N_j\}$. Then $A_0\sub (N_k)_k$ and moreover, for any $0\le l\le p$,
    $$
    A_p\cap [N_p,\infty)+N_l\sub (N_k)_k.
    $$
    In particular, $(A_p)_p$ is compatible if $A_k\cap [1,N_k)=\emptyset$ for every $k\ge1.$
\end{remark}
\begin{proof}
Let $n\in A_p$, $n\ge N_p$. If $n=N_p$ the conclusion is obvius. If $n>N_p$ then $n+N_l\in A_p+N_l\subseteq  \bigcup_{m=0}^{k}\bigcup_{j=0}^m A_{m}+N_j$ for every $k\ge p$.  Then, since $N_k$ is the increasing sequence of minimums $N_{k+1}=\min\{ n>N_{k}\,:\, n\in\bigcup_{m=0}^{k}\bigcup_{j=0}^m A_{m}+N_j\}$ and $n+N_l> N_p$, $n+N_l$ will eventually coincide with some $N_{k}$, $k>p$.  Hence $n+N_l\in \{N_{p+1},N_{p+2},\dots\}$. 
    % Note that $A_p+N_l\sub \bigcup_{m=0}^{k}\bigcup_{j=0}^m A_{m}+N_j$ for every $k\ge p$.   Let $n\in A_p$ with  $n\ge N_p$, then since $N_{k+1}=\min\{ n>N_{k}\,:\, n\in\bigcup_{m=0}^{k}\bigcup_{j=0}^m A_{m}+N_j\}$
    % and , then 
\end{proof}



% The following proposition gives a necessary condition for this question  and will be crucial for our construction.
We will next show that the recurrent arithmetic thickening appears naturally for any recurrent point of any dynamical system and we will then apply this fact to show the general case of the Zero-one law.
\begin{proposition}\label{prop: condic necesaria vectores recurrentes}
    Let $(X,T)$ be a dynamical system (i.e. $X$ a first-countable topological space and $T:X\to X$  continuous). Let $x\in X$ be a recurrent vector, $(U_p)_p$ be a decreasing base of neighbourhoods of $x$ and $A_p\subset N_T(x,U_p)$ be non-empty sets. Then for each $q\in \zN$ there exists a subsequence $(A_{M_p})_p$ such that 
    % the set defined by the relations
    % \begin{equation}\label{definicion A}
    %     A=\bigcup_{p=0}^\infty\bigcup_{j=0}^p A_{M_p}+N_j ,\, \text{ and, }\,\\
    %  N_k=\min\{ n>N_{k-1}\,:\, n\in\bigcup_{p=0}^{k-1}\bigcup_{j=0}^p A_{M_p}+N_j\}
    %    \end{equation}
       $RAT((A_{M_p})_p)\subseteq N_T(x,U_q)$. 
\end{proposition}

\begin{proof}
Let $q\in \zN$. Let $M_0=q$. For $N_1=\min A_q$ we have $T^{N_1}x\in U_q$. Thus there is some $M_1$ such that $T^{N_1}(U_{M_1})\sub U_q$. Observe that this implies that for every $n\in A_{M_1}$ and every $0\leq j\leq 1$ we have that $T^{n+N_j}(x)\in U_q$. Indeed, for $j=0$ it is plain and for $j=1$, $T^{n+N_1}(x)=T^{N_1} T^n(x)\in T^{N_1}(U_{M_1})\subseteq U_q.$ 

We now define $M_{k+1}$ inductively. So we assume that we have defined $M_0,\dots,M_k$ such that for every $0\le j\le p\le k,$
$$
T^{N_j}(U_{M_p})\sub U_q,
$$
where the $N_j$ are defined by the relations \eqref{RAT(A_p)}.
Then since $N_{k+1}=n+N_j$ for some $n\in A_{M_p}$ with $0\le j\le p\le k,$ we have by the inductive hypothesis that $T^{N_{k+1}}x\in U_q$. So by continuity there is $M_{k+1}$ such that $T^{N_{j}}(U_{M_{k+1}})\in U_q$ for every $0\le j\le k+1$. 

It remains to show that $RAT((A_{M_p})_p)\subseteq N_T(x,U_q)$. Let $N\in RAT((A_{M_p})_p)$. Then  $N=n+N_j$ for some $n\in A_{M_k}$ with $0\le j\le k.$ Therefore $T^Nx=T^{n+N_j}x\in T^{N_j}(U_{M_k})\sub U_q $ by the definition of $M_k$.
\end{proof}
\begin{remark}\label{rem:M_k>N_k}
    Note that in the above proof $M_k$ is chosen after defining $N_k$, so we may assume without loss of generality that $M_k>N_k$.
\end{remark}


    We note next that for any sequence $(A_p)_p$ of infinite sets, the recurrent arithmetic thickening $RAT((A_{p})_p)$ is an $IP$-set. In particular we obtain as a consequence of \cite[Theorem 2.17]{Fur81} the following converse of Proposition \ref{prop: condic necesaria vectores recurrentes}.
\begin{proposition}\label{prop: los RAT son IP}
    Let $(A_p)_p$ any sequence of infinite sets, then the recurrent arithmetic thickening $RAT((A_{p})_p)$ contains an $IP$-set.  Moreover, there are a recurrent vector $x$  of some dynamical system $(X,T)$ and a neighborhood $U$ of $x$ such that $N_T(x,U)\sub RAT((A_{p})_p)$.

    Conversely, any $IP$-set contains $RAT((A_{p})_p)$ for some sequence  $(A_p)_p$ of infinite sets.  
\end{proposition}
\begin{proof}
    Take $m_1=N_1=\min A_0$ and $j_1=1$. Take $m_2\in A_{N_{j_1}}\cap (N_{j_1},\infty)$, where $(N_k)_k$ is the sequence in \eqref{RAT(A_p)}.
    By Remark \ref{rem:A_0 contenido en RAT}, we have that $\{m_1,m_{2}, m_1+m_2\}\sub \{N_k:k\in \zN\}$, thus     $m_1+m_2=N_{j_{2}}$ for some $j_{2}> j_1$. Then if $m_3\in A_{N_{j_2}}\cap (N_{j_2},\infty)$ we obtain again by Remark \ref{rem:A_0 contenido en RAT} that $\{m_1,m_{2},m_3, m_1+m_2,m_2+m_3,m_1+m_3,m_1+m_2+m_3\}\sub \{N_k:k\in \zN\}$. Proceeding inductively we obtain the first claim. 

    The second assertion now follows immediately from \cite[Theorem 2.17]{Fur81}.


    For the converse, if $A$ is the $IP$-set generated by the finite sums of an infinite set $A_0$. Consider the constant sequence given by $A_p=A_0$ for every $p$. Then $RAT((A_{p})_p)\sub A.$
\end{proof}


As mentioned above one advantage of the $RAT$ construction (over the $IP$-set construction) is that it allows to preserve properties like density of the original sets. This invites us to consider the notion of $RAT$-set associated to a given family in $\zN$.
\begin{definition}
    Let $\mathcal F$ be any hereditarily upward family of subsets of $\zN$ (i.e. $B_1\in\mathcal F$ and $B_1 \sub B_2$ imply $B_2\in\mathcal F$). We will say that a set $A\sub\zN$ is in $RAT_{\mathcal F}$ if there is a sequence $(A_p)_p$ of sets in $\mathcal F$ such that $ RAT((A_{p})_p)\sub A$. 
\end{definition}
Recall that, given an hereditary upward family $\mathcal F$, a vector $x$ is said to be \textit{$\mathcal F$-recurrent} for $T$ provided that for every open set $U$ around $x$ we have that $N_T(x,U)\in \mathcal F$.

Note that \cite[Theorem 2.17]{Fur81} implies that recurrence is equivalent to $IP$-recurrence.
As an immediate consequence of Proposition \ref{prop: condic necesaria vectores recurrentes} we obtain the following generalization of Furstenberg's Theorem to arbitrary families.
\begin{proposition}\label{prop: RAT_F}
Let $\mathcal F$ be an hereditary upward family and $(X,T)$ a dynamical system. Then $x\in X$ is a $\mathcal F$-recurrent point if and only if it is a $RAT_\mathcal F$-recurrent point.
\end{proposition}

We will not pursue any further on the investigation of $RAT_\mathcal F$-recurrence but we mention that  the notion of $RAT_\mathcal F$-recurrence, for $\mathcal F$ the family of sets with positive lower density, is of key importance for the proof of our main results. In the following construction we will check that
the vectors obtained are $RAT$-frequently recurrent, and hence frequently recurrent.


\subsection {Zero-one law. The general case}\label{section 0-1 general case}
We will prove now the zero-one law for arbitrary Fr\'echet spaces. The same construction of Theorem \ref{0-1 law frequent recurrence c0} works, however some technical details appear when dealing with unconditional basis on non-normable Fr\'echet spaces.
We will present a slightly different proof which relies on the recurrence arithmetic thickening: we first apply the separation Lemma to obtain separated subsets $A_p\subseteq N_B(x,U_p)$. Then we apply Proposition \ref{prop: condic necesaria vectores recurrentes} to obtain a compatible subsequence $A_{M_p}$ so that  $RAT((A_p)_p)=\{N_j\}$. The vector $1_{RAT((A_p)_p)}=\sum_{n\in {RAT((A_p)_p))}} e_{N_j}$ will be our candidate to be frequently recurrent. 

% our first zero-one law Theorem for frequently recurrent backward shifts. The strategy to construct frequently recurrent vectors (near $e_0$) will be the following: starting with a frequently recurrent vector $x\in X$ and a decreasing chain of open sets $U_p$ around $x$, we will use the Separation Lemma \ref{separation lemma} to obtain well separated sets $A_p$ contained in $N_T(x,U_p)$. Then we will fractaly construct our vector $y$ near $e_0$ as follows: in the first step we will paste infinitely many copies of $e_0$ along $A_0$. In the $k$-step, we will paste infinitely many copies of $y^{k-1}$ along some adequate set $A_{M_{k-1}}$. The choice of the sets $A_{M_{k}}$ comes from the fact that the candidate of frequently hypercyclic vector will be of the form $1_A=\sum_{n\in A}e_n$, where $A$ is the set defined by the relations \eqref{definicion A}. In order for $1_A$ to be frequently recurrent we will need $A$ to satisfy the conclusion of Proposition \ref{prop: condic necesaria vectores recurrentes}. 


\begin{theorem}[Zero-one law for frequently recurrent vectors]\label{0-1 law frequent recurrence}
Let $X$ be a Fr\'echet sequence space with unconditional basis $(e_n)_{n\geq 0}$. If $B$ supports a nontrivial frequently recurrent vector then $B$ is frequently recurrent.
\end{theorem}

\begin{proof}
Let $x\neq 0$ be a frequently recurrent vector. It suffices to prove that for each $R\in\mathbb N_0$, and each $\varepsilon>0$ there is a frequently recurrent vector $y\in X$ such that $\|y-e_R\|<\varepsilon$, where $\|\cdot\|$ denotes an $F$-norm on $X$ (see Section \ref{preliminaries}). 

We may suppose that   $|x_R|>\delta$ for some $\delta>0$.   
Since $x$ is a recurrent vector,  there is for every $p\in\mathbb N_0$, a  $q_p>p+R$ such that $|x_{q_p}|>\delta$. By the continuity of the coordinate functionals,  there are a decreasing sequence of positive numbers $(\varepsilon_p)_p$ and a sequence of natural numbers $(k_p)_p$ such that for each $p\in\mathbb N_0$, $\varepsilon_{p}<\frac{\varepsilon}{2^{p+1}}$, $k_p>p+1$ and   
\begin{align}\label{eq: beta_p (zero one freq) UNIL}
 \begin{split}
  \|w\|_{k_p}<\varepsilon_{p}\Longrightarrow  |w_{j}|<\frac{\delta}{4} ,  \text{ for every } 0\leq j\leq q_p, &\quad \text{and}\\
 \|x-\sum_{0\leq j\le k_p} x_j e_j\|_{m_p}<\frac{\varepsilon}{2^{p+1}}\frac{\delta}{8C_p},&
 \end{split}
 \end{align}
 where $C_p$ and $\|\cdot\|_{m_p}$ are the unconditional constant basis and seminorm associated to $\|\cdot\|_p$ (see \eqref{incondicionalidad}). 



We will now present sets $A_p\sub \zN$  that will allow us to construct a frequently recurrent vector near $e_R$. For each $p\in\mathbb N_0$, $A_p$ will be a set of positive lower  density such that  
\begin{equation}\label{eq: A_p zero one freq}
A_p\sub N_B(x,U_p),\quad\text{where }\, U_p=\left\{y\,:\, \|y-x\|_{k_p+m_p}<\varepsilon_p\frac{\delta}{8C_p}\right\}.
\end{equation}

Applying the Separation Lemma \ref{separation lemma}
we may assume that the $A_p$'s are pairwise disjoint and that for every $n\in A_p,m\in A_q$ with $n\neq m$ we have that
\begin{equation}\label{eq separation zero one freq}
|n-m|>k_p+k_q+2R>p+1+q+1+2R.    
\end{equation}

Notice that for each  $p\in \zN_0$ and each $0\leq j\leq p$ we have that, $\sum_{n\in A_p} e_{n+R+j}$ converges. Indeed, by \eqref{eq: beta_p (zero one freq) UNIL} and \eqref{eq: A_p zero one freq} we have that $|x_{n+q_p}-x_{q_p}|<\frac{\delta}{4}$ and hence $|x_{n+q_p}|>\frac{\delta}{2},
$ for all $n\in A_p.$ By unconditionality this implies that $\sum_{n\in A_p} e_{n+q_p}$ converges and by applying the backward shift $q_p-(R+j)$ times it follows that $\sum_{n\in A_p} e_{n+j+R}$ converges for every $0\leq j\leq p$. By the finite invariance property of  the sets of positive lower density, 
we may assume that 
\begin{equation}\label{eq A_p empieza en p}
    A_p\sub [p+R+1,\infty)
\end{equation}
and that,  for every $0\le j\le p,$
\begin{equation}\label{eq: sumas de e_j norma chica (zero one freq)}
 \|\sum_{n\in A_p} e_{n+R+j}\|<\frac{\varepsilon}{(p+1)2^{p+1}}.   
\end{equation}

Applying Proposition \ref{prop: condic necesaria vectores recurrentes} and Remark \ref{rem:M_k>N_k} there is a sequence $(M_p)_p$ for which $RAT((A_{M_p})_p)\subseteq N_B(x,U_0)$ and such that $M_p>N_p$ for every $p\geq 1$ and $M_0=N_0=0$. The fact that $M_p>N_p$ together with \eqref{eq A_p empieza en p}, implies that $(A_{M_p})_p$ is compatible, i.e. 
 $RAT((A_{M_p})_p)=\{N_k\}$.
 % Let  $(M_p)_p$ be the sequence given by Proposition \ref{prop: condic necesaria vectores recurrentes} (for $q=0$).   Consider the recurrent arithmetic thickening $RAT((A_{M_p})_p)=\bigcup_{p=0}^{\infty}\bigcup _{j=0}^p (A_{M_p}+N_j)$, where $N_{k}=\min_{n>N_{k-1}}\{n\in\bigcup_{p=0}^{k-1}\bigcup _{j=0}^p (A_{M_p}+N_j) \}$. By Remark \ref{rem:M_k>N_k} We may assume that $N_k<M_k$ for every $k\ge1$ which, together with \eqref{eq A_p empieza en p}, implies that $(A_{M_p})_p$ is compatible, i.e. 
 % $RAT((A_{M_p})_p)=\{N_k\}$. 

We will show that the vector
\begin{align*}
y:=\hspace{-.3cm}\sum_{n\in RAT((A_{M_p})_p) }\hspace{-.3cm}e_{n+R}=e_R+\sum_{p=0}^{\infty}\sum_{n\in A_{M_p}}\sum_{j=0}^pe_{n+N_j+R}=e_R+\sum_{p=0}^\infty\sum_{n\in A_{M_p}}\sum_{j=0}^{N_p+R}y_je_{n+j}=\sum_{j=0}^\infty e_{N_j+R}.
\end{align*}
is a (well-defined) frequently recurrent vector such that $\|y-e_R\|<\varepsilon.$ The well-definition and the inequality $\|y-e_R\|<\varepsilon$ follow from \eqref{eq: sumas de e_j norma chica (zero one freq)}. 

We prove that $y$ is frequently recurrent. 
 Fix $q\in\zN$ and $m\in A_{M_{q}}$. We will show that $B^m(y)$ is close to $\sum_{j=0}^{N_{q}+R}y_je_j$.
 
 
 Recall that by \eqref{eq separation zero one freq}, since $k_p>p$ for every $p,$ we have that for every $p\in\mathbb N$ and $n\in A_{M_p}$, $|n-m|>M_p+M_{q}+2R>N_p+N_{q}+2R$. Thus, if $n<m$ and $n\in A_{M_{p}}$, then $n-m+N_j+N_q+R<0$ for every $0\le j\le p$. Hence,
$$B^m\left(\sum_{n\in A_{M_p},n<m}\sum_{j=0}^pe_{n+N_j+R}\right)=\sum_{n\in A_{M_p},n<m}\sum_{j=0}^pe_{n-m+N_j+R}=0.$$
Thus, noting that $m$ belongs only to $A_{M_{q}}$, because the $A_{M_p}$'s are pairwise disjoint, we get that
\begin{align*}
B^m(y)-\sum_{j=0}^{N_{q}+R}y_je_j &=B^m(y)-\sum_{j=0}^{q}e_{N_j+R}\\
&=\sum_{p\neq q}\sum_{n\in A_{M_p},n>m}\sum_{j=0}^{p}e_{n-m+N_j+R}+\sum_{n\in A_{M_{q}},n\geq m}\sum_{j=0}^{q}e_{n-m+N_j+R}-\sum_{j=0}^qe_{N_j+R}\\
&=\sum_{p=0}^\infty \sum_{n\in A_{M_p},n>m}\sum_{j=0}^pe_{n-m+N_j+R}.
\end{align*}
By Proposition \ref{prop: condic necesaria vectores recurrentes},  for $0\le j\le p,$ $n\in A_{M_p}$,
$$
B^{n+N_j}x\in U_0,\quad\text{and in particular,}\quad \| B^{n+N_j}x -x\|_{k_0+m_0}<\varepsilon_0.
$$
Then since $|x_R|>\delta,$ we have by \eqref{eq: beta_p (zero one freq) UNIL} that 
\begin{equation*}
    |x_{n+N_j+R}|\ge |x_R|-|x_{n+N_j+R}-x_R| \ge \delta - \frac{\delta}{4}.
\end{equation*}

Since by \eqref{eq separation zero one freq}, $n-m>k_{M_p}+k_{M_{q}}+R>k_{M_q}>k_{N_q}$ whenever $n\in A_{M_p}$ and $n>m$,  we have  $[\sum_{l=0}^{k_{N_{q}}}x_le_l]_{n-m+N_j+R}=0$ for each $0\le j\le p$. Hence it follows that  
 % \begin{align}\label{cota B^my, n>m}
 % \|B^m(y)-\sum_{j=1}^{N_{s_0}+R}y_je_j \|&\leq CR\|B\|^R\|z\|_\infty\left\|\sum_{s=0}^\infty \sum_{n\in A_{M_s},n>m}\sum_{l=0}^se_{n-m+N_l+R}\right\|\\
 % &\leq \frac{4C^2}{\delta}R\|B\|^R\|z\|_\infty\left\|\sum_{s=1}^{\infty}\sum_{n\in A_{M_s},n>m}\sum_{l=0}^sx_{n+N_l+R}e_{n-m+N_l+R}\right\|\\
 % &=\frac{4C^2}{\delta}R\|B\|^R\|z\|_\infty\left\|\sum_{s=0}^{\infty}\sum_{n\in A_{M_s},n>m}\sum_{l=0}^s[B^m(x)-\sum_{j=1}^{k_{N_{s_0}}}x_j]_{n-m+N_l+R}e_{n-m+N_l+R}\right\|\\
 % &\leq \frac{4C^3}{\delta}R\|B\|^R\|z\|_\infty\left(\|B^m(x)-x\|+\|x-\sum_{j=1}^{k_{N_{s_0}}}x_j\|\right)< \varepsilon_{N_{s_0}}.
 % \end{align}
  \begin{align}\label{eq: B^my final}
 \begin{split}
 \left\| B^m(y)-\sum_{j=0}^{N_{q}+R}y_je_j\right\|_q&=
 \left\|\sum_{p=0}^\infty \sum_{n\in A_{M_p},n>m}\sum_{j=0}^pe_{n-m+N_j+R}\right\|_q\\
 &=\left\|\sum_{p=0}^{\infty}\sum_{n\in A_{M_p},n>m}\sum_{j=0}^p\frac{1}{x_{n+N_j+R}}\cdot{x_{n+N_j+R}}e_{n-m+N_j+R}\right\|_q\\
 &=\left\|\sum_{p=0}^{\infty}\sum_{n\in A_{M_p},n>m}\sum_{j=0}^p\frac{1}{x_{n+N_j+R}}\cdot[B^m(x)-\sum_{l=0}^{k_{N_{q}}}x_le_l]_{n-m+N_j+R}e_{n-m+N_j+R}\right\|_q\\
 &\leq \frac{4C_q}{\delta}\left(\|B^m(x)-x\|_{m_q}+\|x-\sum_{l=0}^{k_{N_{q}}}x_le_l\|_{m_q}\right)\\
 & 
\le \frac{4C_q}{\delta}\left(\|B^m(x)-x\|_{k_{M_q}+m_{M_q}}+\|x-\sum_{l=0}^{k_{N_{q}}}x_le_l\|_{m_{N_q}}\right)< \frac{\varepsilon}{2^{N_{q}+1}},
  \end{split}
 \end{align}
where in the last inequality we have used \eqref{eq: beta_p (zero one freq) UNIL} and the definition of $A_{M_p}$ in \eqref{eq: A_p zero one freq}.


 
Finally let $V$ be an open set around $y$. Then since  ${N_q}\to\infty$ as $q\to\infty,$    $\sum_{j=0}^{N_q+R} y_je_j\to y$ and hence,  for big enough  $q\in\mathbb N$, it follows by \eqref{eq: B^my final} that $A_{M_{q}}\sub N_B(y,V).$
This shows that $y$ is a frequently recurrent vector.


\end{proof}


% \begin{proof}
% It suffices to prove that for every $R\geq 0$, and every $\varepsilon>0$ there is a frequently recurrent vector $y\in X$ such that $\|y-e_R\|<\varepsilon$, where $\|\cdot\|$ denotes an $F$-norm on $X$ (see Section \ref{preliminaries}). Indeed, if $U$ is a nonempty open set, let $z=\sum_{j=0}^R z_je_j\in U$. 
% If $y\in X$ be a frequently recurrent vector sufficiently close to $e_R$, the continuity of $B$ and of the sum imply that $\sum_{j=0}^R z_j B^j(y)$ is a frequently recurrent vector in $U$.

% So,  let $x\neq 0$ be a frequently recurrent vector,  $R\geq 0$ and $\varepsilon>0$. Without loss of generality, we may suppose that $x_R\neq 0$ (otherwise consider $B^n(x)$ which is frequently recurrent). Let $\delta>0$ such that $|x_R|>\delta.$


% Since $x$ is a recurrent vector,  there is for every $p\in\mathbb N_0$, a  $q_p>p+R$ such that $|x_{q_p}|>\delta$. By the continuity of the coordinate functionals, there are for each $p\in\mathbb N_0$, $0<\varepsilon_{p}<\frac{\varepsilon}{2^{p+1}}$ and $k_p>p+1$ such that  
% \begin{align}\label{eq: beta_p (zero one freq) UNIL}
%  \begin{split}
%   \|w\|_{k_p}<\varepsilon_{p}\Longrightarrow  |w_{j}|<\frac{\delta}{4} ,  \text{ for every } 0\leq j\leq q_p, &\quad \text{and}\\
%  \|x-\sum_{0\leq j\le k_p} x_j e_j\|_{m_p}<\frac{\varepsilon}{2^{p+1}}\frac{\delta}{8C_p},&
%  \end{split}
%  \end{align}
%  where $C_p$ and $\|\cdot\|_{m_p}$ are the unconditional constant basis and seminorm associated to $\|\cdot\|_p$ (see \eqref{incondicionalidad}). \textcolor{red}{We may assume, without loss of generality, that $m_{p}>k_{p-1}$.}



% We will now present sets $A_p$  that will allow us to construct a frequently recurrent vector near $e_R$. For each $p$, $A_p$ will be a set of positive lower  density such that  
% {\color{red}
% \begin{equation}\label{eq: A_p zero one freq}
% A_p\sub N(x,U_p),\quad\text{where }\, U_p=\left\{y\,:\, \|y-x\|_{k_p+m_p}<\varepsilon_p\frac{\delta}{8C_p}\right\}.
% \end{equation}
% }
% Applying the Separation Lemma \ref{separation lemma}
% we may assume that the $A_p$'s are pairwise disjoint and that for every $n\in A_p,m\in A_q$ with $n\neq m$ we have that
% \begin{equation}\label{eq separation zero one freq}
% n-m>k_p+k_q+2R>p+1+q+1+2R.    
% \end{equation}

% Notice that for each  $p$, $\sum_{n\in A_p} e_{n+p+R}$ converges. Indeed, by \eqref{eq: beta_p (zero one freq) UNIL} and \eqref{eq: A_p zero one freq} we have that $|x_{n+q_p}-x_{q_p}|<\frac{\delta}{4}$ and hence $|x_{n+q_p}|>\frac{\delta}{2}.
% $ By unconditionality this implies that $\sum_{n\in A_p} e_{n+q_p}$ converges and by applying the backward shift $q_p-p+R$ times it follows that $\sum_{n\in A_p} e_{n+p+R}$ converges. By the finite invariance property of  the sets of positive lower density, 
% we may assume that 
% \begin{equation}\label{eq A_p empieza en p}
%     A_p\sub [p+R,\infty)
% \end{equation}
% and that,  for every $0\le j\le p+1,$
% \begin{equation}\label{eq: sumas de e_j norma chica (zero one freq)}
%  \|\sum_{n\in A_p} e_{n+R+j}\|<\frac{\varepsilon}{(p+1)2^{p+1}}.   
% \end{equation}

% {\color{red}
% Let $(N_j)_j$ and $(M_p)_p$ be defined as in Proposition \ref{prop: condic necesaria vectores recurrentes} (for $q=0$). We thus have that for $0\le j\le p,$ $m\in A_{M_p}$,
% $$
% B^{m+N_j}x\in U_0,\quad\text{and in particular, } \| B^{m+N_j}x -x\|_{k_0+m_0}<\varepsilon_0.
% $$
% Note that since $|x_R|>\delta,$ we have by \eqref{eq: beta_p (zero one freq) UNIL} that 
% \begin{equation}\label{eq: coord de x grandes}
%     |x_{m+N_j+R}|\ge |x_R|-|x_{m+N_j+R}-x_R| \ge \delta - \frac{\delta}{4}.
% \end{equation}

% Let now $q\in \mathbb N_0$.
%  Since, by \eqref{eq separation zero one freq}, $|n-m|>N_p+k_{N_{q}}+R$ whenever $n\in A_{M_p}$,  we have  $[\sum_{l=0}^{k_{N_{q}}}x_le_l]_{n-m+N_j+R}=0$ for each $0\le j\le p$ and hence it follows that  
%  % \begin{align}\label{cota B^my, n>m}
%  % \|B^m(y)-\sum_{j=1}^{N_{s_0}+R}y_je_j \|&\leq CR\|B\|^R\|z\|_\infty\left\|\sum_{s=0}^\infty \sum_{n\in A_{M_s},n>m}\sum_{l=0}^se_{n-m+N_l+R}\right\|\\
%  % &\leq \frac{4C^2}{\delta}R\|B\|^R\|z\|_\infty\left\|\sum_{s=1}^{\infty}\sum_{n\in A_{M_s},n>m}\sum_{l=0}^sx_{n+N_l+R}e_{n-m+N_l+R}\right\|\\
%  % &=\frac{4C^2}{\delta}R\|B\|^R\|z\|_\infty\left\|\sum_{s=0}^{\infty}\sum_{n\in A_{M_s},n>m}\sum_{l=0}^s[B^m(x)-\sum_{j=1}^{k_{N_{s_0}}}x_j]_{n-m+N_l+R}e_{n-m+N_l+R}\right\|\\
%  % &\leq \frac{4C^3}{\delta}R\|B\|^R\|z\|_\infty\left(\|B^m(x)-x\|+\|x-\sum_{j=1}^{k_{N_{s_0}}}x_j\|\right)< \varepsilon_{N_{s_0}}.
%  % \end{align}
%   \begin{align}\label{cota B^my, n>m}
%  \begin{split}
%  \left\|\sum_{p=0}^\infty \sum_{n\in A_{M_p},n>m}\sum_{j=0}^pe_{n-m+N_j+R}\right\|_q&=\left\|\sum_{p=0}^{\infty}\sum_{n\in A_{M_p},n>m}\sum_{j=0}^p\frac{1}{x_{n+N_j+R}}\cdot{x_{n+N_j+R}}e_{n-m+N_j+R}\right\|_q\\
%  &=\left\|\sum_{p=0}^{\infty}\sum_{n\in A_{M_p},n>m}\sum_{j=0}^p\frac{1}{x_{n+N_j+R}}\cdot[B^m(x)-\sum_{l=0}^{k_{N_{q}}}x_le_l]_{n-m+N_j+R}e_{n-m+N_j+R}\right\|_q\\
%  &\leq \frac{4C_q}{\delta}\left(\|B^m(x)-x\|_{m_q}+\|x-\sum_{l=0}^{k_{N_{q}}}x_le_l\|_{m_q}\right).
%   \end{split}
%  \end{align}


% }

% We will now proceed to construct a frequently recurrent vector near $e_R$ by induction. \textcolor{red}{ojo que esto ya lo decimos muchas veces}\textcolor{blue}{Si, yo lo volvi a decir 2 veces. Esta prueba en el caso Frechet es la misma que en $c_0$, hay que justificar mucho menos.} The rough idea is to construct first a vector $y^1$ near $e_R$ whose orbit is frequently near $y^0:=e_R$, and by induction a sequence of vectors $(y^k)_k$, contained in a small neighbourhood of $e_R$, and such that the orbit of $y^k$ is frequently near $y^{k-1}$. Those vectors will be constructed so that they   form a convergent sequence to a vector $y$, which will be a frequently recurrent vector near $e_R.$ Each vector $y^k$ will consist on carefully pasting together shifted copies of a truncation of the vector $y^{k-1}$ (and hence of copies of $e_R$).

% {\color{red}
%  {Recall that we are using the sequences $(N_j)_j$ and $(M_p)_p$ given by Proposition \ref{prop: condic necesaria vectores recurrentes}, so that $N_0=M_0=0,$   and $N_{k}=\min_{n>N_{k-1}}\{n\in\bigcup_{p=0}^{k-1}\bigcup _{j=0}^p (A_{M_p}+N_j) \}.$ We will assume that $N_k<M_k$ for every $k\ge1.$}

% Let $y^1=e_R+\sum_{n\in A_0}
% e_{n+R}$. Note that $N_1+R$ is the second nonzero coordinate of $y^1$ and that, in this step, $N_1=\min A_0$. The vector is well defined, because of \eqref{eq: sumas de e_j norma chica (zero one freq)}.
% Notice also that
% $$\|y^1-e_R\|=\|\sum_{n\in A_0}
% e_{n+R}\|<\frac{\varepsilon}{2}.$$




% For the second step, we will add to $y^1$ infinitely many copies of a truncation $P_{N_1+R}(y^1)$ of $y^1$, each of these copies will start at the elements $n$ of $A_{M_1}$. 
% Thus,
% \begin{align*}
%   y^2&=y^1+\sum_{n\in A_{M_1}}\sum_{j=0}^{N_1+R} y^1_j e_{n+j} =y^1+\sum_{n\in A_{M_1}}e_{n+R}+e_{n+N_1+R}  =e_R+\sum_{p=0}^{1}\sum_{j=0}^p\sum_{n\in A_{M_p}} e_{n+N_j+R}. 
% \end{align*}
% Note that since $N_2=\min\{n>N_1\,:\,n\in\bigcup_{p=0}^1\bigcup _{j=0}^p (A_{M_p}+N_j) \},$
% then $N_2+R$ is the third nonzero coordinate of $y^2$. In this case $N_2=\min\{n\in A_{M_0}\cup A_{M_1}\,:\, n>N_1\},$ it may either be the first element of $A_{M_1}$ or the second of $A_{M_0}=A_0$.


% Observe also that, since $|n-m|>k_p+k_q+2R$ whenever $n\in A_p,m\in A_q$ and $k_p>p$, the sums involved have disjoint support. It follows by \eqref{eq: sumas de e_j norma chica (zero one freq)} that the sums are well defined and that 
% $$\|y^2-y^1\|=\|\sum_{n\in A_{M_1}}\sum_{j=0}^1 e_{n+N_j+R}\|\leq  2\frac{\varepsilon}{M_12^{M_1}}<\frac{\varepsilon}{2^2}.$$

% Let us briefly sketch the third step:  
% we define $y^3$, adding to $y^2$ infinitely many copies of $P_{N_2+R}(y^2)$ (which consists of three copies of $e_R$), starting at each element of $n\in A_{M_2}$:
% \begin{align*}
% y^3&=y^2+\sum_{n\in A_{M_2}}\sum_{j=0}^{N_2+R} y^2_j e_{n+j}    =y^2+\sum_{n\in A_{M_2}} \sum_{j=0}^2  e_{n+N_j+R} =e_R+\sum_{p=0}^{2}\sum_{j=0}^p\sum_{n\in A_{M_p}} e_{n+N_j+R}.
% \end{align*}
% Then, proceeding as before by \eqref{eq: sumas de e_j norma chica (zero one freq)}, $y^3$ will be well defined and  $\|y^3-y^2\|<\frac{\varepsilon}{2^3}$. It also satisfies  that its first four nonzero coordinates are $R,\, N_1+R,N_2+R$ and $N_3+R$. \textcolor{red}{habria que explicar mejor por que $N_3+R$ es la 4ta no nula??}


% Let us now proceed to the inductive step. Suppose that we have defined  $y^1,\dots,y^k$ in the same way so that for every  $1\leq j\le k,$
% $$
% \|y^{j}-y^{j-1}\|<\frac{\varepsilon}{2^{j}},
% $$
% and  $R,\, N_1+R,\dots,N_k+R$ are the first $(k+1)$ nonzero coordinates of $y^k$.
% % Define $N_k:=\min\{n>N_{k-1}:n\in  \bigcup_{j=0}^{k-1} A_{N_j}\}$. Note that  
% % \begin{align}\label{eq: indices of the first copies}
% %     N_0+R,N_1+R, \dots,N_{k}+R
% % \end{align} are the indices of the first $k+1$ nonzero coordinates of $y^k$. Note also, for later use, that with this definition,
% % \begin{align}\label{eq: A es la union de los N_k}
% %     \{0\}\cup\bigcup_{j \in \mathbb N_0} A_{N_j}=\bigcup_{j \in \mathbb N_0}\{N_j\}.
% % \end{align}
% We define
% \begin{align}\label{eq: def y^{k+1}}
% y^{k+1}=y^k+\sum_{n\in A_{M_{k}}}\sum_{j=0}^{N_{k}+R} y^k_j e_{n+j} =y^k+\sum_{n\in A_{M_{k}}}\sum_{j=0}^{k}  e_{n+N_j+R}=e_R+\sum_{p=0}^{k}\sum_{j=0}^p\sum_{n\in A_{M_p}} e_{n+N_j+R}.    
% \end{align}

% Note that by \eqref{eq: sumas de e_j norma chica (zero one freq)},
% \begin{align*}
% \|y^{k+1}-y^k\| &=\|\sum_{n\in A_{M_{k}}}\sum_{j=0}^{k}e_{n+N_j+R}\|
% \leq \frac{(k+1)\varepsilon}{(M_{k}+1)2^{M_{k}+1}}<\frac{\varepsilon}{2^{k+1}}.    
% \end{align*}
% This shows that $(y^k)_k$ is a Cauchy sequence so that the limit vector exists and that
% \begin{align*}
% y:=\lim_{k\to \infty} y^k &=e_R+\sum_{p=0}^{\infty}\sum_{n\in A_{M_p}}\sum_{j=0}^pe_{n+N_j+R}=e_R+\sum_{p=0}^\infty\sum_{n\in A_{M_p}}\sum_{j=0}^{N_p+R}y_je_{n+j}
% \end{align*}
% is well defined and that $\|y-e_R\|<\varepsilon.$
% Note that by construction we have
% \begin{align*}
% y=\sum_{j=0}^\infty e_{N_j+R}.
% \end{align*}

% Let us now proceed to the inductive step. Fix $q\in\zN$ and $m\in A_{M_{q}}$. Recall that by \eqref{eq separation zero one freq}, since $k_p>p$ for every $p,$ we have that for every $p\in\mathbb N$ and $n\in A_{M_p}$, $|n-m|>M_p+M_{q}+2R>N_p+N_{q}+2R$. Thus, if $n<m$ and $n\in A_{M_{p}}$, then $n-m+N_j+N_q+R<0$ for every $0\le j\le p$. Hence,
% $$B^m\left(\sum_{n\in A_{M_p},n<m}\sum_{j=0}^pe_{n+N_j+R}\right)=\sum_{n\in A_{M_p},n<m}\sum_{j=0}^pe_{n-m+N_j+R}=0.$$
% Noting that $m$ belongs only to $A_{M_{q}}$, because the $A_{M_p}$'s are pairwise disjoint, we get that
% \begin{align*}
% B^m(y)-\sum_{j=0}^{N_{q}+R}y_je_j &=B^m(y)-\sum_{j=0}^{q}e_{N_j+R}\\
% &=\sum_{p\neq q}\sum_{n\in A_{M_p},n>m}\sum_{j=0}^{p}e_{n-m+N_j+R}+\sum_{n\in A_{M_{q}},n\geq m}\sum_{j=0}^{q}e_{n-m+N_j+R}-\sum_{j=0}^qe_{N_j+R}\\
% &=\sum_{p=0}^\infty \sum_{n\in A_{M_p},n>m}\sum_{j=0}^pe_{n-m+N_j+R}\\
% \end{align*}
% Thus from \eqref{cota B^my, n>m}, together with \eqref{eq: beta_p (zero one freq) UNIL} and the definition of $A_{M_p}$ in \eqref{eq: A_p zero one freq}, we conclude that for any $m\in A_{M_{q}},$
% \begin{align}%\label{eq norma frec zero one}
%  \|B^m(y)-\sum_{j=0}^{N_{q}+R}y_je_j \|_q &= \left\|\sum_{p=0}^\infty \sum_{n\in A_{M_p},n>m}\sum_{j=0}^pe_{n-m+N_j+R}\right\|_q \leq \frac{4C_q}{\delta}\left(\|B^m(x)-x\|_{m_q}+\|x-\sum_{l=0}^{k_{N_{q}}}x_le_l\|_{m_q}\right)\\
% & \label{eq: B^my final}
% \le \frac{4C_q}{\delta}\left(\|B^m(x)-x\|_{k_{M_q}+m_{M_q}}+\|x-\sum_{l=0}^{k_{N_{q}}}x_le_l\|_{m_{N_q}}\right)< \frac{\varepsilon}{2^{N_{q}+1}}.
% \end{align}

 
% Finally let $V$ be an open set around $y$. Then since  ${N_q}\to\infty$ as $q\to\infty,$    $\sum_{j=0}^{N_q+R} y_je_j\to y$ and hence,  for big enough  $q\in\mathbb N$, it follows by \eqref{eq: B^my final} that $A_{M_{q}}\sub N(y,V).$
% This shows that $y$ is a frequently recurrent vector.

% }

% % For the second step, we define $N_1\in\zN$ such that $N_1+R$ is the second nonzero coordinate of $y^1$ which, in this step, coincides with $\min A_0$. Let $M_1>N_1$  such that \textcolor{red}{$ \frac{\delta}{8 C_{M_1}}\|B\|^{N_{1}}\varepsilon_{M_{1}}<\frac{\varepsilon_0}{2^{2}}$. hay que pedirles a los $\varepsilon_p$ que se vayan a 0 mas rapido que $C_p$?} \textcolor{blue}{Hay que pedirselo al $\epsilon$ (asi creo que funciona), o al $M_p$}


% %  We will add to $y^1$ infinitely many copies of a truncation $P_{N_1+R}(y^1)$ of $y^1$:
% % %We construct a vector $y^2$ that will have two copies of $z$ (separated by $N_1-R$ zeros) starting at each $n$ that belongs to $A_{N_1}$, that is,
% % \begin{align*}
% %   y^2&=y^1+\sum_{n\in A_{M_1}}\sum_{j=0}^{N_1+R} y^1_j e_{n+j} =y^1+\sum_{n\in A_{M_1}}e_{n+R}+e_{n+N_1+R}  =e_R+\sum_{p=0}^{1}\sum_{n\in A_{M_p}}\sum_{j=0}^p e_{n+N_j+R}, 
% % \end{align*}
% % where we have set $N_0=M_0=0.$


% % Observe that, since $|n-m|>k_p+k_q+2R$ whenever $n\in A_p,m\in A_q$ and $k_p>p$, the sums involved have disjoint support. Adopting the notation $N_0=0$, it follows by \eqref{eq: sumas de e_j norma chica (zero one freq)} that the sums are well defined and that 
% % $$\|y^2-y^1\|=\|\sum_{n\in A_{M_1}}\sum_{j=0}^1 e_{n+N_j+R}\|\leq  2\frac{\varepsilon}{M_12^{M_1}}<\frac{\varepsilon}{2^2}.$$

% % Let us briefly sketch the third step. Let $N_2\in\zN$ such that $N_2+R$ is the third nonzero coordinate of $y^2$. In this case $N_2=\min\{n\in A_{M_0}\cup A_{M_1}\,:\, n>N_1\},$ it may either be the first element of $A_{M_1}$ or the second of $A_{M_0}=A_0$. \textcolor{red}{Let $M_2>N_2$ such that $ \frac{\delta}{8 C_{M_2}}\|B\|^{N_{2}}\varepsilon_{M_{2}}<\frac{\varepsilon_0}{2^{3}}$.}

% % Let us define $y^3$, adding to $y^2$ infinitely many copies of $P_{N_2+R}(y^2)$ (which consists of three copies of $e_R$), starting at each element of $n\in A_{M_2}$:
% % \begin{align*}
% % y^3&=y^2+\sum_{n\in A_{M_2}}\sum_{j=0}^{N_2+R} y^2_j e_{n+j}    =y^2+\sum_{n\in A_{M_2}} \sum_{j=0}^2  e_{n+N_j+R} =e_R+\sum_{p=0}^{2}\sum_{n\in A_{M_p}}\sum_{j=0}^p e_{n+N_j+R}.
% % \end{align*}
% % Then, proceeding as before, $y^3$ will be well defined and will satisfy, by \eqref{eq: sumas de e_j norma chica (zero one freq)}, that $\|y^3-y^2\|<\frac{\varepsilon}{2^3}$ and that its first three nonzero coordinates are $R,\, N_1+R$ and $N_2+R$.

% % We point out that there is a subtle difference from the fourth step onwards. Again $N_3\in\zN$ is defined so that $N_3+R$ is the fourth nonzero coordinate of $y^3$, but in this case  $N_3$ is not necessarily $\min\{n\in A_{M_0}\cup A_{M_1}\cup A_{M_2}\,:\, n>N_2\}.$ This is because it may happen that $N_1+N_2+R$ is the fourth nonzero coordinate of $y^3$, and $N_1+N_2\notin A_{M_0}\cup A_{M_1}\cup A_{M_2}$. \textcolor{blue}{ In any case, since the nonzero coordinates of $y^3$ are of the form $n+N_j+R$, with $n\in A_{M_p}$, $1\leq p\leq 2$ and $0\leq j\leq p$, we have that 
% % $$N_3=\min_{n>N_2}\{n\in\bigcup_{p=0}^2\bigcup _{j=0}^p (A_{M_p}+N_j) \}.$$} 
% % The definition of $y^4$ and $M_3$ is analogous.



% % Let us now proceed to the inductive step. Suppose that we have defined $N_1,\dots, N_{k-1},M_1,\dots, M_{k-1}$ and $y^1,\dots,y^k$ in the same way and that for every  $1\leq j\le k,$
% % $$
% % \|y^{j}-y^{j-1}\|<\frac{\varepsilon}{2^{j}}.
% % $$
% % \textcolor{red}{Define $N_k$ such that $N_k+R$ is the $(k+1)$-th nonzero coordinate of $y^k$  and let $M_k>N_k$ such that 
% % \begin{equation}\label{eq: definicion M_k}
% %     \frac{\delta}{8 C_{M_k}}\|B\|^{N_{k}} \varepsilon_{M_{k}}<\frac{\varepsilon_0}{2^{k+1}}.
% % \end{equation}}

% % % Define $N_k:=\min\{n>N_{k-1}:n\in  \bigcup_{j=0}^{k-1} A_{N_j}\}$. Note that  
% % % \begin{align}\label{eq: indices of the first copies}
% % %     N_0+R,N_1+R, \dots,N_{k}+R
% % % \end{align} are the indices of the first $k+1$ nonzero coordinates of $y^k$. Note also, for later use, that with this definition,
% % % \begin{align}\label{eq: A es la union de los N_k}
% % %     \{0\}\cup\bigcup_{j \in \mathbb N_0} A_{N_j}=\bigcup_{j \in \mathbb N_0}\{N_j\}.
% % % \end{align}

% % Let
% % \begin{align}\label{eq: def y^{k+1}}
% % y^{k+1}=y^k+\sum_{n\in A_{M_{k}}}\sum_{j=0}^{N_{k}+R} y^k_j e_{n+j} =y^k+\sum_{n\in A_{M_{k}}}\sum_{j=0}^{k}  e_{n+N_j+R}=e_R+\sum_{p=0}^{k}\sum_{n\in A_{M_p}}\sum_{j=0}^p e_{n+N_j+R}.    
% % \end{align}
% % Note that by \eqref{eq: sumas de e_j norma chica (zero one freq)},
% % \begin{align*}
% % \|y^{k+1}-y^k\| &=\|\sum_{n\in A_{M_{k}}}\sum_{j=0}^{k}e_{n+N_j+R}\|
% % \leq \frac{(k+1)\varepsilon}{(M_{k}+1)2^{M_{k}+1}}<\frac{\varepsilon}{2^{k+1}}.    
% % \end{align*}
% % This shows that $(y^k)_k$ is a Cauchy sequence so that the limit vector exists and that
% % \begin{align*}
% % y:=\lim_{k\to \infty} y^k &=e_R+\sum_{p=0}^{\infty}\sum_{n\in A_{M_p}}\sum_{j=0}^pe_{n+N_j+R}=e_R+\sum_{p=0}^\infty\sum_{n\in A_{M_p}}\sum_{j=0}^{N_p+R}y_je_{n+j}
% % \end{align*}
% % is well defined and that $\|y-e_R\|<\varepsilon.$
% % Note that by construction we have
% % \begin{align*}
% % y=\sum_{j=0}^\infty e_{N_j+R}.
% % \end{align*}


% % \textcolor{red}{hasta aca cambie}

% % It remains to prove that $y$ is a frequently recurrent vector. Fix $q\in\zN$ and $m\in A_{M_{q}}$. Recall that by \eqref{eq separation zero one freq}, since $k_p>p$ for every $p,$ we have that for every $p\in\mathbb N$ and $n\in A_{M_p}$, $|n-m|>M_p+M_{q}+2R>N_p+N_{q}+2R$. Thus, if $n<m$ and $n\in A_{M_{p}}$, then $n-m+N_j+N_q+R<0$ for every $0\le j\le p$. Hence,
% % $$B^m\left(\sum_{n\in A_{M_p},n<m}\sum_{j=0}^pe_{n+N_j+R}\right)=\sum_{n\in A_{M_p},n<m}\sum_{j=0}^pe_{n-m+N_j+R}=0.$$
% % Noting that $m$ belongs only to $A_{M_{q}}$, because the $A_{M_p}$'s are pairwise disjoint, we get that
% % \begin{align*}
% % B^m(y)-\sum_{j=0}^{N_{q}+R}y_je_j &=B^m(y)-\sum_{j=0}^{q}e_{N_j+R}\\
% % &=\sum_{p\neq q}\sum_{n\in A_{M_p},n>m}\sum_{j=0}^{p}e_{n-m+N_j+R}+\sum_{n\in A_{M_{q}},n\geq m}\sum_{j=0}^{q}e_{n-m+N_j+R}-\sum_{j=0}^qe_{N_j+R}\\
% % &=\sum_{p=0}^\infty \sum_{n\in A_{M_p},n>m}\sum_{j=0}^pe_{n-m+N_j+R}\\
% % \end{align*}
% % Thus,
% % \begin{equation}\label{eq norma frec zero one}
% %  \|B^m(y)-\sum_{j=0}^{N_{q}+R}y_je_j \|_q= \left\|\sum_{p=0}^\infty \sum_{n\in A_{M_p},n>m}\sum_{j=0}^pe_{n-m+N_j+R}\right\|_q.
% % \end{equation}

% % %%%%%%%%%%%
% % We will show that the expression in \eqref{eq norma frec zero one} converges to 0 as $q$ increases to $\infty$. To achieve this, we claim that for every $N_j$ we have that $\|B^{N_j}x-x\|_{k_j}<\varepsilon_0-\frac{\varepsilon_0}{2^j}<\varepsilon_0$.
% % Indeed, for $N_0=0$ it is plain and for $N_1=\min A_0$ it follows because $A_0\subseteq \{n\in \zN: \|B^n(x)-x\|_{k_1}\le \|B^n(x)-x\|_{k_0+m_0}< \frac{\varepsilon_0 \delta}{8 C_0}\}.$

% % Let $k>1$. By the construction of the  $N_k$'s, we have that $N_k+R$ is the  $(k+1)$-th nonzero coordinate of $$y^{k}=\sum_{p=0}^{k-1}\sum_{n\in A_{M_p}} \sum_{j=0}^p e_{n+N_j+R}.$$ 
% % By a quick inspection we have that $N_k$ must be a number of the form $n+N_j$, where $n\in A_{M_p}$, $0\leq p\leq {k-1}$ and $0\leq j\leq p$. Hence by \eqref{eq: definicion M_k} and the inductive hypothesis,
% % \begin{align*}
% % \|B^{N_k}x-x\|_\infty & =\|B^{n+N_j}x-x\|_\infty\leq \|B^n B^{N_j}(x)-B^{N_j}x\|_\infty+\|B^{N_j}(x)-x\|_\infty    \\
% % &\le \|B\|^{N_j}\,\|B^{n}x-x\|_\infty+ \varepsilon_0(1-\frac1{2^j}) < \frac{\varepsilon_0}{2^{k}}+  \varepsilon_0(1-\frac1{2^{k-1}}) = \varepsilon_0(1-\frac1{2^k}).
% % \end{align*}



% % This implies that for every $j\in \zN_0$, $\|B^{N_j}x-x\|_{\infty}<\varepsilon_0$ and since $|x_R|>\delta$, we have by \eqref{eq: beta_p (zero one freq) UNIL c_0} that $|x_{N_j+R}|>\frac{\delta}{2}$. On the other hand, we have that for every $n\in A_{M_p},$ $\|B^{n}x-x\|_{\infty}<\varepsilon_{M_p}$ and hence by \eqref{eq: beta_p (zero one freq) UNIL c_0} and the fact that $q_{M_p}>M_p+R>N_j+R,$ we obtain for $0\leq j\leq p$ and $n\in A_{M_p}$ that
% % $|x_{n+N_j+R}-x_{N_j+R}|<\frac{\delta}{4}$. This implies that for every $n\in A_{M_p}$ and $0\leq j\leq p$, $|x_{n+N_j+R}|>\frac{\delta}{4}.$


% % %%%%%%%%%%


% % Note that for every $j\ge 0$, $N_j\in\{0\}\cup\bigcup_{p=0}^\infty A_p$. This implies that for every $j\in \zN_0$, $\|B^{N_j}x-x\|_{k_0}<\varepsilon_0$ and since $|x_R|>\delta$, we have by \eqref{eq: beta_p (zero one freq) UNIL} that $|x_{N_j+R}|>\frac{\delta}{2}$. On the other hand, we have that for every $n\in A_{M_p},$ $\|B^{n}x-x\|_{k_p}<\varepsilon_{N_p}$ and hence by \eqref{eq: beta_p (zero one freq) UNIL} we obtain for $0\leq j\leq p$ and $n\in A_{M_p}$ that
% % $|x_{n+N_j+R}-x_{N_j+R}|<\frac{\delta}{4}$. This implies that for every $n\in A_{M_p}$ and $0\leq j\leq p$, $|x_{n+N_j+R}|>\frac{\delta}{4}.$
 
% %  Since, by \eqref{eq separation zero one freq}, $|n-m|>N_p+k_{N_{q}}+R$ whenever $n\in A_{M_p}$,  we have  $[\sum_{l=0}^{k_{N_{q}}}x_le_l]_{n-m+N_j+R}=0$ for each $0\le j\le p$ and hence it follows that  
% %  % \begin{align}\label{cota B^my, n>m}
% %  % \|B^m(y)-\sum_{j=1}^{N_{s_0}+R}y_je_j \|&\leq CR\|B\|^R\|z\|_\infty\left\|\sum_{s=0}^\infty \sum_{n\in A_{M_s},n>m}\sum_{l=0}^se_{n-m+N_l+R}\right\|\\
% %  % &\leq \frac{4C^2}{\delta}R\|B\|^R\|z\|_\infty\left\|\sum_{s=1}^{\infty}\sum_{n\in A_{M_s},n>m}\sum_{l=0}^sx_{n+N_l+R}e_{n-m+N_l+R}\right\|\\
% %  % &=\frac{4C^2}{\delta}R\|B\|^R\|z\|_\infty\left\|\sum_{s=0}^{\infty}\sum_{n\in A_{M_s},n>m}\sum_{l=0}^s[B^m(x)-\sum_{j=1}^{k_{N_{s_0}}}x_j]_{n-m+N_l+R}e_{n-m+N_l+R}\right\|\\
% %  % &\leq \frac{4C^3}{\delta}R\|B\|^R\|z\|_\infty\left(\|B^m(x)-x\|+\|x-\sum_{j=1}^{k_{N_{s_0}}}x_j\|\right)< \varepsilon_{N_{s_0}}.
% %  % \end{align}
% %   \begin{align}\label{cota B^my, n>m}
% %  \begin{split}
% %  \left\|\sum_{p=0}^\infty \sum_{n\in A_{M_p},n>m}\sum_{j=0}^pe_{n-m+N_j+R}\right\|_q&=\left\|\sum_{p=0}^{\infty}\sum_{n\in A_{M_p},n>m}\sum_{j=0}^p\frac{1}{x_{n+N_j+R}}\cdot{x_{n+N_j+R}}e_{n-m+N_j+R}\right\|_q\\
% %  &=\left\|\sum_{p=0}^{\infty}\sum_{n\in A_{M_p},n>m}\sum_{j=0}^p\frac{1}{x_{n+N_j+R}}\cdot[B^m(x)-\sum_{l=0}^{k_{N_{q}}}x_le_l]_{n-m+N_j+R}e_{n-m+N_j+R}\right\|_q\\
% %  &\leq \frac{4C_q}{\delta}\left(\|B^m(x)-x\|_{m_q}+\|x-\sum_{l=0}^{k_{N_{q}}}x_le_l\|_{m_q}\right).
% %   \end{split}
% %  \end{align}

% % Therefore, using  \eqref{eq norma frec zero one}  together with \eqref{eq: beta_p (zero one freq) UNIL} and the definition of $A_{M_p}$ in \eqref{eq: A_p zero one freq}, we conclude that for any $m\in A_{M_{q}},$
% % \begin{align}\label{eq: B^my final}
% % \|B^m(y)-\sum_{j=0}^{N_{q}+R}y_je_j \|_q\le \frac{4C_q}{\delta}\left(\|B^m(x)-x\|_{k_{N_q}+m_{N_q}}+\|x-\sum_{l=0}^{k_{N_{q}}}x_le_l\|_{m_{N_q}}\right)< \frac{\varepsilon}{2^{N_{q}+1}}.
% %     \end{align}

 
% % Finally let $V$ be an open set around $y$. Then since  ${N_q}\to\infty$ as $q\to\infty,$    $\sum_{j=0}^{N_q+R} y_je_j\to y$ and hence,  for big enough  $q\in\mathbb N$, it follows by \eqref{eq: B^my final} that $A_{M_{q}}\sub N(y,V).$
% % This shows that $y$ is a frequently recurrent vector.
% \end{proof}


% \begin{proof}
% It suffices to prove that for every $R\geq 0$, and every $\varepsilon>0$ there is a frequently recurrent vector $y\in X$ such that $\|y-e_R\|<\varepsilon$, where $\|\cdot\|$ denotes an $F$-norm on $X$ (see Section \ref{preliminaries}). Indeed, if $U$ is a nonempty open set, let $z=\sum_{j=0}^R z_je_j\in U$. 
% If $y\in X$ be a frequently recurrent vector sufficiently close to $e_R$, the continuity of $B$ and of the sum imply that $\sum_{j=0}^R z_j B^j(y)$ is a frequently recurrent vector in $U$.

% So,  let $x\neq 0$ be a frequently recurrent vector,  $R\geq 0$ and $\varepsilon>0$. Without loss of generality, we may suppose that $x_R\neq 0$ (otherwise consider $B^n(x)$ which is frequently recurrent). Let $\delta>0$ such that $|x_R|>\delta.$


% Since $x$ is a recurrent vector,  there is for every $p\in\mathbb N_0$, a  $q_p>p+R$ such that $|x_{q_p}|>\delta$. By the continuity of the coordinate functionals, there are for each $p\in\mathbb N_0$, $0<\varepsilon_{p}<\frac{\varepsilon}{2^{p+1}}$ and $k_p>p+1$ such that  
% \begin{align}\label{eq: beta_p (zero one freq) UNIL}
%  \begin{split}
%   \|w\|_{k_p}<\varepsilon_{p}\Longrightarrow  |w_{j}|<\frac{\delta}{4} ,  \text{ for every } 0\leq j\leq q_p, &\quad \text{and}\\
%  \|x-\sum_{0\leq j\le k_p} x_j e_j\|_{m_p}<\frac{\varepsilon}{2^{p+1}}\frac{\delta}{8C_p},&
%  \end{split}
%  \end{align}
%  where $C_p$ and $\|\cdot\|_{m_p}$ are the unconditional constant basis and seminorm associated to $\|\cdot\|_p$ (see \eqref{incondicionalidad}). 



% We will now present sets $A_p$  that will allow us to construct a frequent recurrent vector near $e_R$. For each $p$, $A_p$ will be a set of positive lower  density such that  
% \begin{equation}\label{eq: A_p zero one freq}
% A_p\sub\{n\in\zN:\|B^n(x)-x\|_{k_p+m_p}<\varepsilon_p\frac{\delta}{8C_p}\}.
% \end{equation}
% Applying the Separation Lemma \ref{separation lemma}
% we may assume that the $A_p$'s are pairwise disjoint and that for every $n\in A_p,m\in A_q$ with $n\neq m$ we have that
% \begin{equation}\label{eq separation zero one freq}
% n-m>k_p+k_q+2R>p+1+q+1+2R.    
% \end{equation}

% Notice that if $n\in A_p$ then $\sum_{n\in A_p} e_{n+p+R}$ converges. Indeed, by \eqref{eq: beta_p (zero one freq) UNIL} and \eqref{eq: A_p zero one freq} we have that $|x_{n+q_p}-x_{q_p}|<\frac{\delta}{4}$ and hence $|x_{n+q_p}|>\frac{\delta}{2}.
% $ By uncondionality this implies that $\sum_{n\in A_p} e_{n+q_p}$ converges and by applying the backward shift $q_p-p+R$ times it follows that $\sum_{n\in A_p} e_{n+p+R}$ converges. By the finite invariance property of  the sets of positive lower density, we may assume that $A_p\sub [p+R,\infty)$ and that,  for every $0\le j\le p+1,$
% \begin{equation}\label{eq: sumas de e_j norma chica (zero one freq)}
%  \|\sum_{n\in A_p} e_{n+R+j}\|<\frac{\varepsilon}{(p+1)2^{p+1}}.   
% \end{equation}

% We will now proceed to construct a frequently recurrent vector near $e_R$ by induction. The rough idea is to construct first a vector $y^1$ near $e_R$ whose orbit is frequently near $y^0:=e_R$, and by induction a sequence of vectors $(y^k)_k$, contained in a small neighbourhood of $e_R$, and such that the orbit of $y^k$ is frequently near $y^{k-1}$. Those vectors will be constructed so that they   form a convergent sequence to a vector $y$, which will be a frequently recurrent vector near $e_R.$ Each vector $y^k$ will consist on carefully pasting together shifted copies of a truncation of the vector $y^{k-1}$ (and hence of copies of $e_R$).
 
% Let $y^1=e_R+\sum_{n\in A_0}
% e_{n+R}$.  The vector is well defined, because of \eqref{eq: sumas de e_j norma chica (zero one freq)}. Notice also that
% $$\|y^1-e_R\|=\|\sum_{n\in A_0}
% e_{n+R}\|<\frac{\varepsilon}{2}.$$

% For the second step let $N_1=\min A_0$. Note that $N_1+R$ is the index of the second nonzero  coordinate of $y^1$. We will add to $y^1$ infinitely many copies of a truncation $P_{N_1+R}(y^1)$ of $y^1$:
% %We construct a vector $y^2$ that will have two copies of $z$ (separated by $N_1-R$ zeros) starting at each $n$ that belongs to $A_{N_1}$, that is,
% \begin{align*}
%   y^2&=y^1+\sum_{n\in A_{N_1}}\sum_{j=0}^{N_1+R} y^1_j e_{n+j} =y^1+\sum_{n\in A_{N_1}}e_{n+R}+e_{n+N_1+R}  =e_R+\sum_{p=0}^{1}\sum_{n\in A_{N_p}}\sum_{j=0}^p e_{n+N_j+R}. 
% \end{align*}

% Observe that, since $|n-m|>k_p+k_q+2R$ whenever $n\in A_p,m\in A_q$ and $k_p>p$, the sums involved have disjoint support. Adopting the notation $N_0=0$, it follows by \eqref{eq: sumas de e_j norma chica (zero one freq)} that the sums are well defined and that 
% $$\|y^2-y^1\|=\|\sum_{n\in A_{N_1}}\sum_{j=0}^1 e_{n+N_j+R}\|\leq  2\frac{\varepsilon}{N_12^{N_1}}<\frac{\varepsilon}{2^2}.$$

% Let us briefly sketch the third step. 
% Let $N_2=\min\{n\in A_{N_0}\cup A_{N_1}\,:\, n>N_1\}$, i.e. $N_2$ is either the first element of $A_{N_1}$ or the second element of $A_0=A_{N_0}$.
% Note that $N_2+R$ is the index of the third nonzero  coordinate of $y^2$.
% Let us define $y^3$, with three copies of $e_R$ starting at each element of $n\in A_{N_2}$:
% \begin{align*}
% y^3&=y^2+\sum_{n\in A_{N_2}}\sum_{j=0}^{N_2+R} y^2_j e_{n+j}    =y^2+\sum_{n\in A_{N_2}} \sum_{j=0}^2  e_{n+N_j+R} =e_R+\sum_{p=0}^{2}\sum_{n\in A_{N_p}}\sum_{j=0}^p e_{n+N_j+R}.
% \end{align*}
% Then, proceeding as before, $y^3$ will be well defined and will satisfy, by \eqref{eq: sumas de e_j norma chica (zero one freq)}, that $\|y^3-y^2\|<\frac{\varepsilon}{2^3}.$

% Let us now proceed to the inductive step. Suppose that we have defined $N_1,\dots, N_{k-1}$ and $y^1,\dots,y^k$ in the same way and that for every  $1\leq j\le k,$
% $$
% \|y^{j}-y^{j-1}\|<\frac{\varepsilon}{2^{j}}.
% $$
% Define $N_k:=\min\{n>N_{k-1}:n\in  \bigcup_{j=0}^{k-1} A_{N_j}\}$. Note that  
% \begin{align}\label{eq: indices of the first copies}
%     N_0+R,N_1+R, \dots,N_{k}+R
% \end{align} are the indices of the first $k+1$ nonzero coordinates of $y^k$. Note also, for later use, that with this definition,
% \begin{align}\label{eq: A es la union de los N_k}
%     \{0\}\cup\bigcup_{j \in \mathbb N_0} A_{N_j}=\bigcup_{j \in \mathbb N_0}\{N_j\}.
% \end{align}

% Let
% \begin{align}\label{eq: def y^{k+1}}
% y^{k+1}=y^k+\sum_{n\in A_{N_{k}}}\sum_{j=0}^{N_{k}+R} y^k_j e_{n+j} =y^k+\sum_{n\in A_{N_{k}}}\sum_{j=0}^{k}  e_{n+N_j+R}=e_R+\sum_{p=0}^{k}\sum_{n\in A_{N_p}}\sum_{j=0}^p e_{n+N_j+R}.    
% \end{align}
% Note that by \eqref{eq: sumas de e_j norma chica (zero one freq)},
% \begin{align*}
% \|y^{k+1}-y^k\| &=\|\sum_{n\in A_{N_{k}}}\sum_{j=0}^{k}e_{n+N_j+R}\|
% \leq \frac{(k+1)\varepsilon}{(N_{k}+1)2^{N_{k}+1}}<\frac{\varepsilon}{2^{k+1}}.    
% \end{align*}
% This shows that $(y^k)_k$ is a Cauchy sequence so that the limit vector exists and that

% \begin{align*}
% y:=\lim_{k\to \infty} y^k &=e_R+\sum_{p=0}^{\infty}\sum_{n\in A_{N_p}}\sum_{j=0}^pe_{n+N_j+R}=e_R+\sum_{p=0}^\infty\sum_{n\in A_{N_p}}\sum_{j=0}^{N_p+R}y_je_{n+j}
% \end{align*}
% is well defined and that $\|y-e_R\|<\varepsilon.$
% Note that by \eqref{eq: indices of the first copies}, we have
% \begin{align*}
% y=\sum_{j=0}^\infty e_{N_j+R}.
% \end{align*}


% It remains to prove that $y$ is a frequently recurrent vector. Fix $q\in\zN$ and $m\in A_{N_{q}}$. Recall that by \eqref{eq separation zero one freq}, since $k_p>p$ for every $p,$ we have that for every $p\in\mathbb N$ and $n\in A_{N_p}$, $|n-m|>N_p+N_{q}+2R$. Thus, if $n<m$ and $n\in A_{N_{p}}$, then $n-m+N_j+N_q+R<0$ for every $0\le j\le p$. Hence,
% $$B^m\left(\sum_{n\in A_{N_p},n<m}\sum_{j=0}^pe_{n+N_j+R}\right)=\sum_{n\in A_{N_p},n<m}\sum_{j=0}^pe_{n-m+N_j+R}=0.$$
% Noting that $m$ belongs only to $A_{N_{q}}$, because the $A_{N_p}$'s are pairwise disjoint, we get that
% \begin{align*}
% B^m(y)-\sum_{j=0}^{N_{q}+R}y_je_j &=B^m(y)-\sum_{j=0}^{q}e_{N_j+R}\\
% &=\sum_{p\neq q}\sum_{n\in A_{N_p},n>m}\sum_{j=0}^{p}e_{n-m+N_j+R}+\sum_{n\in A_{N_{q}},n\geq m}\sum_{j=0}^{q}e_{n-m+N_j+R}-\sum_{j=0}^qe_{N_j+R}\\
% &=\sum_{p=0}^\infty \sum_{n\in A_{N_p},n>m}\sum_{j=0}^pe_{n-m+N_j+R}\\
% \end{align*}
% Thus,
% \begin{equation}\label{eq norma frec zero one}
%  \|B^m(y)-\sum_{j=0}^{N_{q}+R}y_je_j \|_q= \left\|\sum_{p=0}^\infty \sum_{n\in A_{N_p},n>m}\sum_{j=0}^pe_{n-m+N_j+R}\right\|_q.
% \end{equation}

% Note that for every $j\ge 0$, $N_j\in\{0\}\cup\bigcup_{p=0}^\infty A_p$. This implies that for every $j\in \zN_0$, $\|B^{N_j}x-x\|_{k_0}<\varepsilon_0$ and since $|x_R|>\delta$, we have by \eqref{eq: beta_p (zero one freq) UNIL} that $|x_{N_j+R}|>\frac{\delta}{2}$. On the other hand, we have that for every $n\in A_{N_p},$ $\|B^{n}x-x\|_{k_p}<\varepsilon_{N_p}$ and hence by \eqref{eq: beta_p (zero one freq) UNIL} we obtain for $0\leq j\leq p$ and $n\in A_{N_p}$ that
% $|x_{n+N_j+R}-x_{N_j+R}|<\frac{\delta}{4}$. This implies that for every $n\in A_{N_p}$ and $0\leq j\leq p$, $|x_{n+N_j+R}|>\frac{\delta}{4}.$
 
%  Since, by \eqref{eq separation zero one freq}, $|n-m|>N_p+k_{N_{q}}+R$ whenever $n\in A_{N_p}$,  we have  $[\sum_{l=0}^{k_{N_{q}}}x_le_l]_{n-m+N_j+R}=0$ for each $0\le j\le p$ and hence it follows that  
%  % \begin{align}\label{cota B^my, n>m}
%  % \|B^m(y)-\sum_{j=1}^{N_{s_0}+R}y_je_j \|&\leq CR\|B\|^R\|z\|_\infty\left\|\sum_{s=0}^\infty \sum_{n\in A_{N_s},n>m}\sum_{l=0}^se_{n-m+N_l+R}\right\|\\
%  % &\leq \frac{4C^2}{\delta}R\|B\|^R\|z\|_\infty\left\|\sum_{s=1}^{\infty}\sum_{n\in A_{N_s},n>m}\sum_{l=0}^sx_{n+N_l+R}e_{n-m+N_l+R}\right\|\\
%  % &=\frac{4C^2}{\delta}R\|B\|^R\|z\|_\infty\left\|\sum_{s=0}^{\infty}\sum_{n\in A_{N_s},n>m}\sum_{l=0}^s[B^m(x)-\sum_{j=1}^{k_{N_{s_0}}}x_j]_{n-m+N_l+R}e_{n-m+N_l+R}\right\|\\
%  % &\leq \frac{4C^3}{\delta}R\|B\|^R\|z\|_\infty\left(\|B^m(x)-x\|+\|x-\sum_{j=1}^{k_{N_{s_0}}}x_j\|\right)< \varepsilon_{N_{s_0}}.
%  % \end{align}
%   \begin{align}\label{cota B^my, n>m}
%  \begin{split}
%  \left\|\sum_{p=0}^\infty \sum_{n\in A_{N_p},n>m}\sum_{j=0}^pe_{n-m+N_j+R}\right\|_q&=\left\|\sum_{p=0}^{\infty}\sum_{n\in A_{N_p},n>m}\sum_{j=0}^p\frac{1}{x_{n+N_j+R}}\cdot{x_{n+N_j+R}}e_{n-m+N_j+R}\right\|_q\\
%  &=\left\|\sum_{p=0}^{\infty}\sum_{n\in A_{N_p},n>m}\sum_{j=0}^p\frac{1}{x_{n+N_j+R}}\cdot[B^m(x)-\sum_{l=0}^{k_{N_{q}}}x_le_l]_{n-m+N_j+R}e_{n-m+N_j+R}\right\|_q\\
%  &\leq \frac{4C_q}{\delta}\left(\|B^m(x)-x\|_{m_q}+\|x-\sum_{l=0}^{k_{N_{q}}}x_le_l\|_{m_q}\right).
%   \end{split}
%  \end{align}

% Therefore, using  \eqref{eq norma frec zero one}  together with \eqref{eq: beta_p (zero one freq) UNIL} and the definition of $A_{N_p}$ in \eqref{eq: A_p zero one freq}, we conclude that for any $m\in A_{N_{q}},$
% \begin{align}\label{eq: B^my final}
% \|B^m(y)-\sum_{j=0}^{N_{q}+R}y_je_j \|_q\le \frac{4C_q}{\delta}\left(\|B^m(x)-x\|_{k_{N_q}+m_{N_q}}+\|x-\sum_{l=0}^{k_{N_{q}}}x_le_l\|_{m_{N_q}}\right)< \frac{\varepsilon}{2^{N_{q}+1}}.
%     \end{align}

 
% Finally let $V$ be an open set around $y$. Then since  ${N_q}\to\infty$ as $q\to\infty,$    $\sum_{j=0}^{N_q+R} y_je_j\to y$ and hence,  for big enough  $q\in\mathbb N$, it follows by \eqref{eq: B^my final} that $A_{N_{q}}\sub N(y,V).$
% This shows that $y$ is a frequently recurrent vector.
% \end{proof}

% \textcolor{blue}{Funcionar\'a algo as\'i?
% \begin{remark}\label{remark caract unilat}
% We extract from the proof of Theorem \ref{0-1 law frequent recurrence} the following facts which will be of vital importance for the proof of the next Theorem. Taking  $R=0$ and considering $A=\{0\}\cup\bigcup_{p \geq 0} \bigcup_{j=0}^pA_{M_p}+N_j=\bigcup_{k\geq 0}\{N_k\}$ (see \eqref{eq: A es la union de los N_k}), there is $\varepsilon_p\to 0$ such that
% \begin{itemize}
%     \item for every $p\in \zN$, $\sum_{n\in A_{N_p}} e_{n+N_p}$ converges (see \eqref{eq: sumas de e_j norma chica (zero one freq)}), and
%     \item for infinitely many $q\in \mathbb N$ we have that for every $m\in A_{M_{q}}$ (see \eqref{cota B^my, n>m} and \eqref{eq: B^my final}),
%     $$\|\sum_{p=0}^\infty \sum_{n\in A_{M_p}, n>m} \sum_{j=0}^p e_{n-m+N_j}\|_q<\varepsilon_q.$$
% \end{itemize}
% \end{remark}
% }


\subsection{First characterization of frequently recurrent backward shifts\label{section 0-1 general case, primera caracterizacion}}
A careful look to the proof of Theorem \ref{0-1 law frequent recurrence} allows us to isolate  the conditions that provide our first characterization for frequently recurrent backward shifts.



\begin{theorem}[First characterization of frequent recurrence]\label{Characterization frequent recurrence}
{Let $X$ be a Fr\'echet sequence space with unconditional basis $(e_n)_{n\geq 0}$. The following are equivalent: }
\begin{enumerate}
\item $B$ is frequently recurrent;
\item
 there are subsets $A_p$ with positive lower density 
 and a sequence $\varepsilon_p\to 0$,  such that
\begin{enumerate}
    \item [i)]for every $p\in\mathbb N_0$, $\sum_{n\in A_p}e_{n+p}$ converges; 
    \item [ii)]
    for any $q\in\mathbb N_0$ and every  $m\in A_q$,
    $$
    \left\|\sum_{\overset{N\in RAT((A_p)_p)}{N>m}} e_{N-m}\right\|_q<\varepsilon_q.
    $$
\end{enumerate}

\end{enumerate}
\end{theorem}  

\begin{proof}[Proof of $(1)\Rightarrow (2)$.] 
     This is a consequence of the proof of Theorem \ref{0-1 law frequent recurrence}. Indeed, we have seen there that if $B$ is frequently recurrent then taking $R=0$,  the sequences $(N_j)_j$, $(A_{M_q})_q$ and $(\varepsilon_p)_p=(\frac{\varepsilon}{2^{p+1}})_p$ satisfy condition $i)$ by \eqref{eq: sumas de e_j norma chica (zero one freq)} and condition $ii)$ by \eqref{eq: B^my final}.  
\end{proof}    

For the converse, we wish to assume that the sets $A_p$ are well separated. This is not longer an immediate application of separation Lemma \ref{separation lemma} because we don't know whether condition $ii)$ remains to hold when we take subsets of the $A_p's$.


To avoid these difficulties we will use the following lemma which after an application of the separation Lemma considers an appropriate subsequence $A_{M_p}'\sub A_{M_p}$ which form a compatible sequence. We will then show that the estimations $i)$ and  $ii)$ also hold for $RAT((A_{M_p}')_p)$.




{ 
\begin{lemma}\label{lem: RAT subconjuntos}
    Let $(A_p)_p$ be a sequence of subsets of $\zN$  and let  $A_p'\sub A_p\cap [N_p,\infty)$ be  a sequence of subsets. Then  there is a subsequence  $(M_p)_p$ (which can be chosen to grow as fast as desired)  such that 
    $(A_{M_p}')_p$ is a compatible sequence and $RAT((A_{M_p}')_p)\sub RAT((A_p)_p)$.
\end{lemma}

\begin{proof}
Let $N_1'=\min A_0'$. Since by Remark \ref{rem:A_0 contenido en RAT}, $A_0'\sub A_0\sub (N_k)_k$, there is $k_0\in \zN$ such that $N_1'=N_{k_0}$. 
Then $(A_{p_0}'+N_1')\cup A_{p_0}'\sub  (A_{p_0}+N_{k_0})\cup A_{p_0}\subseteq RAT((A_p)_p)$ for every $p_0\ge k_0$.
Thus if $M_1>\max\{k_0,N_1'\}$ then 
$$
\bigcup_{p=0}^{1} \bigcup_{j=0}^p A_{M_p}'+N_j'= A_0'\cup A_{M_1}'\cup (A_{M_1}'+N_1')\subseteq RAT((A_p)_p).
$$
Suppose now that we have defined $M_1,\ldots,M_{k-1}$, such that, 
\begin{enumerate}
    \item[a)] $N_{l}'=\min_{n>N_{l-1}'} \bigcup_{p=0}^{l-1} \bigcup_{j=0}^p A_{M_p}'+N_j' \in \{N_0,N_1,\dots,N_{M_l}\}$ for every $1\leq l\leq k-1$,
   \item[b)] $M_{l}>N_l'$ for every $1\leq l\leq k-1$,
    \item [c)]$\bigcup_{p=0}^{k-1} \bigcup_{j=0}^p A_{M_p}'+N_j'\subseteq RAT((A_p)_p)$.
\end{enumerate}

 
Let $N_k'=\min_{n>N_{k-1}'}\bigcup_{p=0}^{k-1} \bigcup_{j=0}^p A_{M_p}'+N_j'$. By Remark \ref{rem:A_0 contenido en RAT} and the fact that $A_p'\sub A_p\cap [N_p,\infty)$ it follows for $j\le p$ that since $N_j'<M_p<N_{M_p}$,  $A_{M_p}'+N_j' \in (N_i)_i$ for every $0\le j\le p\le k-1$. This implies that $N_k'\in (N_i)_i$. 


Thus we have for large enough $p$, 
\begin{equation*}
    A_{p}'+N_j'\subseteq A_p+ N_j'\subseteq RAT((A_p)_p)
\end{equation*} 
for every $0\leq j\leq k$. So, we choose  $M_k$ large enough so that $A_{M_{k}}'+N_j'\subseteq RAT((A_p)_p)$ for every $0\leq j\leq k$, $M_k>N_k'$. Notice that, by c) we have  $\bigcup_{p=0}^{k} \bigcup_{j=0}^p A_{M_p}'+N_j'\subseteq RAT((A_p)_p)$.



By c) we obtain that $RAT((A_{M_p}')_p)\subseteq RAT((A_{p})_p)$.

Moreover, since $A_p'\sub [N_p,\infty)$ and  $M_p>N_p'$, Remark \ref{rem:A_0 contenido en RAT} implies that $(A_{M_p}')_p$ is compatible.
\end{proof}
\begin{remark}
    \label{rem: cosas del lemma que usamos}
We note for later use that the sequence $(M_p)_p$ in the above lemma can be chosen so that 
$$
N_{l}'\in \{N_0,N_1,\dots,N_{M_l}\}  \quad \text{ and, }\quad M_{l}>N_l' \text{ for every }l\ge1,
$$
\end{remark}
}


\begin{proof}[Proof of  $(2)\Rightarrow (1)$ in Theorem \ref{Characterization frequent recurrence}]    
 by Theorem \ref{0-1 law frequent recurrence} it suffices to construct a nonzero frequently recurrent vector near $e_0$. Let $(N_k)_k$ be the sequence associated to $RAT((A_p)_p)$, i.e. $N_k=\min \{n>N_{k-1}: n\in \bigcup_{p=0}^{k-1} \bigcup_{j=0}^p A_p+N_j\}$ and $N_0=0.$

We begin applying the separation Lemma \ref{separation lemma} to obtain sets of positive lower density $A_p'\subseteq A_p$  such that 
\begin{equation}\label{eq: caracterizacion 1 separation de los A_p}
    |n-m|>p+q
\end{equation} whenever $n\in A_p',m\in A_q'$ and $n\neq m$. By $i)$ and unconditionality,  $\sum_{n\in A_p'} e_{n+j}\in X$; and by finite invariance we may suppose that  $A_p'\sub [N_{p}+1,\infty)$ and that 
\begin{equation}\label{eq: caracterizacion 1 finitely invariance}
    \|\sum_{n\in A_p'} e_{n+j}\|\leq \frac{1}{2^p}
\end{equation} for every $0\le j\le p$.

By Lemma \ref{lem: RAT subconjuntos}, there is a subsequence $(M_p)_p$ such that $(A_{M_p}')_p$ is compatible and $RAT((A_{M_p}')_p)\sub RAT((A_p)_p)$.
 Moreover, if $RAT((A_{M_p}')_p)=(N_j')_j$, $N_{l}'=\min_{n>N_{l-1}'} \bigcup_{p=0}^{l-1} \bigcup_{j=0}^p A_{M_p}'+N_j'$ for every $l$ and we may assume that 
\begin{enumerate}
    %\item[a)] $N_{l}'=\min_{n>N_{l-1}'} \bigcup_{p=0}^{l-1} \bigcup_{j=0}^p A_{M_p}'+N_j'$ for every $1\leq l\leq k-1$,
   \item[a)] \label{eq: caracterizacion 1 Los M_l son grandes} $M_{l}>\max\{m_{l},N_l'\}$ for every $1\leq l$, (see Remark \ref{rem: cosas del lemma que usamos})
   \item [b)] $\varepsilon_{M_l}<\varepsilon_{M_{l-1}}$ for $1\leq l$, and,
    \item [c)] $C_l\cdot\varepsilon_{M_l}<\frac{1}{l}$ for every $1\leq l$,
\end{enumerate}
where $C_l$ and $m_l$ are the constants arising     from the unconditionality of the basis.

%Let $N_1'=\min A_0'$. Since $A_0'\sub A=\{N_k\}$, there is $k_0$ such that $N_1'=N_{k_0}$. By the definition of $A$ we have that $A_p'+N_{k_0}\sub A_p+N_{k_0}\subseteq A$, whenever $p>k_0$. If $m_1$ and $C_1$ are the numbers given in \eqref{incondicionalidad}, we choose $M_1>\max\{k_0,m_1,N_1'\}$ large enough so that $C_1\cdot \varepsilon_{M_1}<1$.
% Suppose now that we have defined $M_1,\ldots,M_{k-1}$, such that, 
% \begin{enumerate}
%     \item[a)] $N_{l}'=\min_{n>N_{l-1}'} \bigcup_{p=0}^{l-1} \bigcup_{j=0}^p A_{M_p}'+N_j'$ for every $1\leq l\leq k-1$,
%    \item[b)] \label{eq: caracterizacion 1 Los M_l son grandes} $M_{l}>\max\{m_{l},N_l'\}$ for every $1\leq l\leq k-1$,
%    \item [c)] $\varepsilon_{M_l}<\varepsilon_{M_{l-1}}$ for $1\leq l\leq k-1$,
%     \item [d)] $C_l\cdot\varepsilon_{M_l}<\frac{1}{l}$ for every $1\leq l\leq k-1$,
%      and,
%     \item [e)]$\bigcup_{p=0}^{k-1} \bigcup_{j=0}^p A_{M_p}'+N_j'\subseteq A$.
% \end{enumerate}
% Let $N_k'=\min_{n>N_{k-1}'}\bigcup_{p=0}^{k-1} \bigcup_{j=0}^p A_{M_p}'+N_j'$.
% Since $N_1',\ldots ,N_k'\in A$, we have that for large enough $p$, 
% \begin{equation}\label{eq: caracterizacion 1 Ap'}
%     A_{p}'+N_j'\subseteq A_p+ N_j'\subseteq A
% \end{equation} 
% for every $0\leq j\leq k$. So, we choose large enough $M_k$ so that $A_{M_{k}}'+N_j'\subseteq A$ for every $1\leq j\leq k$, $M_k>\max\{m_k,N_k'\}$, $C_k\cdot \varepsilon_{M_k}<\frac{1}{k}$ and $\varepsilon_{M_l}<\varepsilon_{M_{l-1}}.$ 

% Notice that, by e) and \eqref{eq: caracterizacion 1 Ap'}, we have  $\bigcup_{p=0}^{k} \bigcup_{j=0}^p A_{M_p}'+N_j'\subseteq A$. 

We have so far obtained sets of positive lower density $A_{M_p}'$  satisfying the separation property \eqref{eq: caracterizacion 1 separation de los A_p} and the estimation \eqref{eq: caracterizacion 1 finitely invariance}. Moreover, since $RAT((A_{M_p}')_p)\subseteq RAT((A_{p})_p)$, we have by unconditionality  that $RAT((A_{M_p}')_p)$ satisfy the inequality in $ii)$ for  $\varepsilon_q'=\frac1{q}$
%but only  for the subsequence $(M_{m_q})_q$
, i.e. for every $q\in \zN$ and $m\in A_{M_{q}}'$ we have that
\begin{equation}\label{eq: norma de la suma RAT(A') <1/q}
\left\|\sum_{\overset{N\in RAT((A_{M_p}')_p)}{N>m}} e_{N-m}\right\|_q<\frac{1}{q}.    
\end{equation}



Indeed, let $m\in A_{M_q}'$. We have by the separation property  \eqref{eq: caracterizacion 1 separation de los A_p}    together with b) that for $n\in A_{M_p}'$,
\begin{equation}\label{eq: separation Ap' first characterization}
|m-n|>M_{q}+M_p> M_{q}+N_p'.  
\end{equation}
Thus $\displaystyle \sum_{\overset{N\in RAT((A_{M_p}')_p)}{N>m}} e_{N-m} =    \sum_{p=0}^\infty\sum_{j=0}^p\sum_{\overset{n\in A_{M_{p}}'}{n+N_j'>m}} e_{n-m+N_j'}=    \sum_{p=0}^\infty\sum_{j=0}^p\sum_{\overset{n\in A_{M_{p}}'}{ n>m}} e_{n-m+N_j'}$. 

Moreover, by Remark \ref{rem: cosas del lemma que usamos} we have that for every $p\in \zN$, $\{N_1',\ldots ,N_{p}'\}\subseteq \{N_1,\ldots , N_{M_p}\}$  and hence, using unconditionality,  a) and c), we obtain that
\begin{align}
\nonumber \left\|\sum_{\overset{N\in RAT((A_{M_p}')_p)}{N>m}} e_{N-m}\right\|_q & =    \left\|\sum_{p=0}^\infty\sum_{\overset{n\in A_{M_{p}}'}{ n>m}}\sum_{j=0}^p e_{n-m+N_j'}\right\|_q \leq C_q\cdot \left\|\sum_{p=0}^\infty\sum_{\overset{n\in A_{M_p}}{n>m}}\sum_{j=0}^{M_p} e_{n-m+N_j}\right\|_{m_q}\\
 \label{eq: norma de la suma n>m <1/q}   & \leq C_q\cdot\left\|\sum_{p=0}^\infty\sum_{\overset{n\in A_{M_p}}{n>m}}\sum_{j=0}^{M_p} e_{n-m+N_j}\right\|_{M_q}
     <C_q\cdot \varepsilon_{M_{q}}%<C_q\cdot \varepsilon_{M_q}
    <\frac{1}{q}. 
\end{align}

The hard work has already been done. Denoting $A'=RAT((A_{M_p}')_p)$, we have that $1_{A'}=\sum_{n\in A'} e_n=\sum_{j=0}^\infty e_{N_j'}=\sum_{p=0}^\infty\sum_{n\in A_{M_{p}}'}\sum_{j=0}^p e_{n+N_j'} $ converges because, by \eqref{eq: caracterizacion 1 finitely invariance}, $\sum_{p=0}^\infty \|\sum_{n\in A_{M_p}'} \sum_{0\leq j\leq p} e_{n+N_j'}\|\leq \sum_{p=0}^\infty \frac{p+1}{2^p}<\infty$. We show next that $1_{A'}$ is a frequently recurrent vector.

Let $m\in A_{M_{q}}'$. Then for $N\in A_{M_p}'$, $N<m$ we have  that $ B^m(e_{N})=0$. Hence, by \eqref{eq: norma de la suma n>m <1/q},
\begin{align*}
\left\|B^{m}(1_{A'})-\sum_{j=0}^{q} e_{N_j'}
\right\|_q & = \left\|\sum_{\overset{N\in A'}{N\ge m}} e_{N-m} -\sum_{j=0}^{q} e_{N_j'}\right\|_q \\
&= \left\|\left(\sum_{p=0}^\infty\sum_{\overset{n\in A_{M_{p}}'}{n>m}}\sum_{j=0}^p e_{n-m+N_j'}+\sum_{j=0}^{q} e_{N_j'}\right)-\sum_{j=0}^{q} e_{N_j'}\right\|_q
<\frac{1}{q},    
\end{align*}
which implies that $1_{A'}$ is frequently recurrent.
\end{proof}

% {\color{violet} PRUEBA VIEJA

% \begin{proof}
    
% For the converse, by Theorem \ref{0-1 law frequent recurrence} it suffices to construct a nonzero frequently recurrent vector near $e_0$. 


% We wish to assume that the sets $A_p$ are well separated. This is not longer an immediate application of separation Lemma \ref{separation lemma} because we don't want to loose the arithmetic properties of $A=\{N_j\}=\bigcup_{p=0}^\infty \bigcup_{j=0}^p A_p+N_j$. At the same time, there is another technical obstacle related to the growth of the unconditional constants on Fr\'echet spaces, if we just apply unconditionality in ii), we would obtain an estimation of the form $C_q\varepsilon_{m_q}$ which could not decay to 0. This will be avoided by taking more subsequences.

% To avoid these difficulties we will first apply the separation Lemma and then  consider an appropriate subsequence $A_{M_p}'$ for which  for which $A'=\{N_j'\}=\bigcup_{p=0}^\infty\bigcup_{j=0}^p A'_{M_k}+N_j'\subseteq A$, where for each $k$ we will have $A'_{M_k}\subseteq A_{M_k}$ and the estimations $i)$ and  $ii)$ hold for $A_{M_p}'$ and $A'$.

% We begin applying the separation Lemma \ref{separation lemma} to obtain sets of positive lower density $A_p'\subseteq A_p$  such that 
% \begin{equation}\label{eq: caracterizacion 1 separation de los A_p}
%     |n-m|>p+q
% \end{equation} whenever $n\in A_p',m\in A_q'$ and $n\neq m$. By $i)$ and unconditionality,  $\sum_{n\in A_p'} e_{n+j}\in X$; and by finite invariance we may suppose that  $A_p'\sub [p+1,\infty)$ and that 
% \begin{equation}\label{eq: caracterizacion 1 finitely invariance}
%     \|\sum_{n\in A_p'} e_{n+j}\|\leq \frac{1}{2^p}
% \end{equation} for every $0\le j\le p$.

% Let $N_1'=\min A_0'$. Since $A_0'\sub A=\{N_k\}$, there is $k_0$ such that $N_1'=N_{k_0}$. By the definition of $A$ we have that $A_p'+N_{k_0}\sub A_p+N_{k_0}\subseteq A$, whenever $p>k_0$. If $m_1$ and $C_1$ are the numbers given in \eqref{incondicionalidad}, we choose $M_1>\max\{k_0,m_1,N_1'\}$ large enough so that $C_1\cdot \varepsilon_{M_1}<1$.

% Suppose now that we have defined $M_1,\ldots,M_{k-1}$, such that, 
% \begin{enumerate}
%     \item[a)] $N_{l}'=\min_{n>N_{l-1}'} \bigcup_{p=0}^{l-1} \bigcup_{j=0}^p A_{M_p}'+N_j'$ for every $1\leq l\leq k-1$,
%    \item[b)] \label{eq: caracterizacion 1 Los M_l son grandes} $M_{l}>\max\{m_{l},N_l'\}$ for every $1\leq l\leq k-1$,
%    \item [c)] $\varepsilon_{M_l}<\varepsilon_{M_{l-1}}$ for $1\leq l\leq k-1$,
%     \item [d)] $C_l\cdot\varepsilon_{M_l}<\frac{1}{l}$ for every $1\leq l\leq k-1$,
%      and,
%     \item [e)]$\bigcup_{p=0}^{k-1} \bigcup_{j=0}^p A_{M_p}'+N_j'\subseteq A$.
% \end{enumerate}

 
% Let $N_k'=\min_{n>N_{k-1}'}\bigcup_{p=0}^{k-1} \bigcup_{j=0}^p A_{M_p}'+N_j'$.
% Since $N_1',\ldots ,N_k'\in A$, we have that for large enough $p$, 
% \begin{equation}\label{eq: caracterizacion 1 Ap'}
%     A_{p}'+N_j'\subseteq A_p+ N_j'\subseteq A
% \end{equation} 
% for every $0\leq j\leq k$. So, we choose large enough $M_k$ so that $A_{M_{k}}'+N_j'\subseteq A$ for every $1\leq j\leq k$, $M_k>\max\{m_k,N_k'\}$, $C_k\cdot \varepsilon_{M_k}<\frac{1}{k}$ and $\varepsilon_{M_l}<\varepsilon_{M_{l-1}}.$ 

% Notice that, by e) and \eqref{eq: caracterizacion 1 Ap'}, we have  $\bigcup_{p=0}^{k} \bigcup_{j=0}^p A_{M_p}'+N_j'\subseteq A$. 

% We have sofar obtained sets of positive lower density $A_{M_p'}$  satisfying the separation property \eqref{eq: caracterizacion 1 separation de los A_p} and the estimation \eqref{eq: caracterizacion 1 finitely invariance}. Moreover, since $A'=\{N_j'\}=\bigcup_{p=0}^\infty \bigcup_{j=0}^p A_{M_p}'+ N_j'\subseteq A$, we have by unconditionality  that $(A_{M_p}')$ and $A'$ satisfy the inequality in $ii)$ 
% %but only  for the subsequence $(M_{m_q})_q$
% , i.e. for every $q\in \zN$ and $m\in A_{M_{q}}'$ we have that
% \begin{equation}\label{eq: norma de la suma n>m <1/q}
% \|\sum_{p=0}^\infty\sum_{n\in A_{M_p}', n>m}\sum_{j=0}^{M_p} e_{n-m+N_j'}\|_q=\|\sum_{n\in A', n>m} e_{n-m}\|_q<\frac{1}{q}.    
% \end{equation}



% Indeed, if $m\in A_{M_q}'$ we have by b) that $\{N_0',\dots,N_p'\}\sub \{N_0,\dots,N_{M_p}\}$ and hence, 
% \begin{align*}
%     \|\sum_{p=0}^\infty\sum_{n\in A_{M_{p}}', n>m}\sum_{j=0}^p e_{n-m+N_j'}\|_q&\leq C_q\cdot \|\sum_{p=0}^\infty\sum_{n\in A_{M_p}, n>m}\sum_{j=0}^{M_p} e_{n-m+N_j}\|_{m_q}\\
%     & \leq C_q\cdot\|\sum_{p=0}^\infty\sum_{n\in A_{M_p}, n>m}\sum_{j=0}^{M_p} e_{n-m+N_j}\|_{M_q}
%      <C_q\cdot \varepsilon_{M_{q}}%<C_q\cdot \varepsilon_{M_q}
%     <\frac{1}{q}. 
% \end{align*}

% The hard work has already been done. We have that $1_{A'}=\sum_{n\in A'} e_n=\sum_{j=0}^\infty e_{N_j'}=\sum_{p=0}^\infty\sum_{n\in A_{M_{p}}'}\sum_{j=0}^p e_{n+N_j'} $ converges because, by \eqref{eq: caracterizacion 1 finitely invariance}, $\sum_{p=0}^\infty \|\sum_{n\in A_{M_p}'} \sum_{0\leq j\leq p} e_{n+N_j'}\|\leq \sum_{p=0}^\infty \frac{p+1}{2^p}<\infty$. We show next that $1_{A'}$ is a frequently recurrent vector.

% Let $m\in A_{M_{m_q}}'$. Then for $n\in A_{M_p}'$, $n<m$ we have by the separation property \eqref{eq: caracterizacion 1 separation de los A_p}  together with b) that
% $$
% m-n>M_{q}+M_p> M_{q}+N_p',
% $$
% and thus $B^m\left(\sum_{p=0}^\infty\sum_{n\in A_{M_{p}}', n<m}\sum_{j=0}^p e_{n-m+N_j'}\right)=0$. Hence, by \eqref{eq: norma de la suma n>m <1/q},
% $$
% \left\|B^{m}(1_{A'})-\sum_{j=0}^{q} e_{N_j'}
% \right\|_q= \left\|\left(\sum_{p=0}^\infty\sum_{n\in A_{M_{p}}', n>m}\sum_{j=0}^p e_{n-m+N_j'}+\sum_{j=0}^{q} e_{N_j'}\right)-\sum_{j=0}^{q} e_{N_j'}\right\|_q
% <\frac{1}{q}
% $$
% which implies that $1_{A'}$ is frequently recurrent.
% \end{proof}
% }

With the same proof we can present another characterization which follows the spirit of the characterization of frequent hypercyclicity from Theorem \ref{caracterizacion hiperciclicidad frecuente} obtained by Bonilla and Grosse-Erdmann in \cite{BonGro18}.
If the sets $(A_p)_p$ are sufficiently separated then the equivalence $(2)$ in the following just rephrases Theorem \ref{Characterization frequent recurrence}. The equivalence $(3)$ is a consequence of the proof of Theorem \ref{Characterization frequent recurrence}. 
\begin{theorem}\label{coro: first Characterization frequent recurrence}
{Let $X$ be a Fr\'echet sequence space with unconditional basis $(e_n)_{n\geq 0}$. The following are equivalent: }
\begin{enumerate}
\item $B$ is frequently recurrent;

\item  there  are subsets $A_p\subseteq \zN$ with positive lower density, $(N_k)_k$ defined by the relation
 $N_k=\min_{n>N_{k-1}} \{n\in\bigcup_{p=0}^{k-1}\bigcup_{j=0}^p A_p+N_j\}$  
  and a sequence $\varepsilon_p\to 0$,  such that
\begin{enumerate}
    \item [i)]for every $p\in\mathbb N_0$, $\sum_{n\in A_p}e_{n+p}$ converges; 
    \item [ii)]
    for any $q\in\mathbb N_0$ and every  $m\in A_q$,
    $$
    \left\|\sum_{p=0}^\infty \sum_{n\in A_p, n>m} \sum_{j=0}^p e_{n-m+N_j}\right\|_q<\varepsilon_q.
    $$

\end{enumerate}



\item  there is a compatible sequence  $(A_p)_p$ formed by subsets $A_p\subseteq \zN$ with positive lower density such that
$$
\sum_{n\in RAT((A_p)_p)}\hspace{-.5cm}e_{n}\quad\text{is a frequently recurrent vector.}
$$
\end{enumerate}

 Moreover in (2) we may assume that the sequence $(A_p)_p$ satisfies:
\begin{itemize}
        \item it is well separated: for every $p,q\in\mathbb N_0$ and every $n\in A_p$, $m\in A_q$ with $n\neq m$ we have that $|n-m|>p+q.$
    \item for every $p\in \zN_0$, $\|\sum_{n\in A_p} e_{n+N_p}\|$ is arbitrarily small, and
    \item it is compatible: $\{N_j\}=RAT((A_p)_p).$
\end{itemize}


\end{theorem}  



This characterization of frequent recurrence is similar to the characterization of frequent hypercyclicity from Theorem \ref{caracterizacion hiperciclicidad frecuente} obtained by Bonilla and Grosse-Erdmann in \cite{BonGro18}, but it is weaker in two aspects: while condition i) is the same on both characterizations, we require in ii) an upper bound which depends only on $\varepsilon_q$ (and not on the minimum between  $\varepsilon_p$ and $\varepsilon_q$). Also, the bound of ii) holds only for a subset of the natural numbers $A=\{N_k\}$, while in \cite{BonGro18} the authors need $A=\mathbb N_0$. Of course, the case $A=\mathbb N_0$ is a sufficient condition for frequent recurrence.
\begin{corollary}\label{Corolario zN}
Let $X$ be a Fr\'echet sequence space with unconditional basis $(e_n)_{n\in \mathbb N}$. Suppose that there are sets $A_p$ with positive lower density such that
\begin{enumerate}
\item [i)]for every $p\in\mathbb N_0$, $\sum_{n\in A_p}e_{n+p}$ converges 
    \item [ii)]
    For every $q\in\mathbb N_0$ and every  $m\in A_q$,
    $$\|\sum_{p=0}^\infty \sum_{n\in A_p, n>m} \sum_{j=0}^p e_{n-m+j}\|_q<\varepsilon_q.$$
    \end{enumerate}Then $B$ is frequently recurrent.
\end{corollary}
\begin{proof}
The proof is similar: we construct inductively a sequence $(M_p)_p$ such that $M_0=0$,
    %\item[a)] $N_{l}'=\min_{n>N_{l-1}'} \bigcup_{p=0}^{l-1} \bigcup_{j=0}^p A_{M_p}'+N_j'$ for every $1\leq l\leq k-1$,
  $M_{l}>\max\{m_{l},N_l'\}$,
 $\varepsilon_{M_l}<\varepsilon_{M_{l-1}}$ and,
  $C_l\cdot\varepsilon_{M_l}<\frac{1}{l}$ for every $1\leq l$,
where $C_l,m_l$ are the constants from the unconditionality, $N_0'=0$ and $N_k' =\min_{n>N_{k-1}'} \{n\in\bigcup_{p=0}^{k-1}\bigcup_{j=0}^p A_{M_p}+N_j'\}$.

Hence we have,
\begin{align*}
\|\sum_{p=0}^\infty \sum_{n\in A_{M_p}, n>m}\sum_{j=0}^p e_{n-m+N_j}\|_q &\le C_q\|\sum_{p=0}^\infty \sum_{n\in A_{p}, n>m}\sum_{j=0}^p e_{n-m+j}\|_{m_q}\\
&\le C_q\|\sum_{p=0}^\infty \sum_{n\in A_{p}, n>m}\sum_{j=0}^p e_{n-m+j}\|_{M_q}<C_q\varepsilon_{M_q}<\frac1{q}.   
\end{align*}
Thus, the sequence $(A_{M_p})_p$ satisfies the conditions of Theorem \ref{coro: first Characterization frequent recurrence} (2) and hence $B$ is frequently recurrent.
\end{proof}



Weighted backwards shifts in $c_0$ provide a great source of counter-examples in linear dynamics \cite{BayRuz15,BonGro18,BayGri07,Bes16}. It is therefore meaningful to enunciate Theorem \ref{Characterization frequent recurrence} and Corollary \ref{Corolario zN} for the space $c_0$.
\begin{theorem}
Let $X=c_0$ and $B_w$ be a weighted backward shift. Then $B_w$ is frequently recurrent if and only if there are a sets of positive lower density $A_p$  such that 
\begin{enumerate}
    \item [i)] for each $p\in \zN_0$, $\lim_{n\in A_{p}} \prod_{l=1}^{n+p} w_{l}=0$; 
\end{enumerate}    
    and such that one of the following conditions holds
    \begin{enumerate}
        
    \item [ii)]for every $q\in \zN_0$, every $n\in RAT((A_p)_p)$ and every $m\in A_q$ with $n>m$ we have that
$$\prod_{l=1}^{n-m} w_{l}\leq \varepsilon_q$$
for every $0\leq j\leq p$;
    \item [ii')] If $N_k=\min_{n>N_{k-1}} \{\bigcup_{p=0}^{k-1}\bigcup_{j=0}^k A_p+N_j\}$, then for every $p,q\in \zN_0$ and every $n\in A_p$, $m\in A_q$ with $n>m$ we have that
    
$$\prod_{l=1}^{n-m+N_j} w_{l}\leq \varepsilon_q.$$
    \end{enumerate}
\end{theorem}
\begin{corollary}
Let $X=c_0$ and $B_w$ be a weighted backward shift. Suppose that there are sets of positive lower density $A_p$ and a sequence of positive numbers $\varepsilon_p\to 0$ such that 
\begin{itemize}
    \item[(i)] for every $p\in\zN_0$, $\lim_{n\in A_{p}} \prod_{l=1}^{n+p} w_{l}=0$
    \item[(ii)] for every $p,q\in \zN_0$ and every $n\in A_p$, $m\in A_q$ with $n>m$ we have that
$$\prod_{l=1}^{n-m+j} w_{l}\leq \varepsilon_q$$
for every $0\leq j\leq p$.
\end{itemize} Then $B$ is frequently recurrent.
\end{corollary}


\subsection{Second characterization of frequent recurrence and an application to dense lineability. }\label{subsection segunda aplicacion}
We will now show another equivalence for frequent recurrence in terms of the existence of  orbits with wild behavior. As an application we will show that the set of frequently recurrent vectors is dense lineable. Given a set $A\subseteq \zN_0$ we consider $\mathfrak B_M(A)$ the set of vectors supported in $A$ that have coordinates bounded by $M>0$, i.e.
$$\mathfrak B_M(A):=\{x\in X: supp(x)\subseteq A \text{ and } |x_n|<M \text { for every } n\geq 0\}.$$

We recall the following definition from \cite {Gro19bilateral}:  given $T:X\to X$ and a subset $Y\subseteq X$ we say that $x\in X$ is a \emph{frequently hypercyclic vector for $Y$} if for every open set $U$ such that $U\cap Y\neq \emptyset$ we have that $N_T(x,U)$ has positive lower density.

\begin{theorem}[Second characterization of frequent recurrence] \label{thm: caract S_1(A)}
Let $X$ be a Fr\'echet sequence space with unconditional basis $(e_n)_{n\ge 0}$. Then $B$ is frequently recurrent if and only if there are $A\subseteq \mathbb N_0$ with $\underline {dens}(A)>0$ and $x\in X$, supported in $A$, such that $x$ is a frequently hypercyclic vector for $\mathfrak B_1(A)$. 
\end{theorem}

\begin{proof}
If there are $A$ and $x\in X$ supported in $A$ such that $x$ is frequently hypercyclic for $\mathfrak B_1(A)$ then we have by unconditionality that $P_{\mathfrak B_1(A)} x:= \sum_{n\in A} \min\{|x_n|,1\}sgn(x_n) e_n$ is frequently hypercyclic for $\mathfrak B_1(A)$, where $sgn$ denotes either the real or complex sign of a number depending the field of $X$. Indeed, for any $y\in \mathfrak B_1(A)$, $q\in \mathbb N_0$, we have,
$$
\|B^m(P_{\mathfrak B_1(A)} x)-y\|_q\le C_q \|B^m( x)-y\|_{m_q}.
$$


Since $P_{\mathfrak B_1(A)} x$ belongs itself to $\mathfrak B_1(A)$ then it is a frequently recurrent vector. By the zero-one law Theorem \ref{0-1 law frequent recurrence} we deduce that the operator is frequently recurrent.  

Suppose now that $B$ is frequently recurrent. 
By Theorem \ref{coro: first Characterization frequent recurrence} there are $A_p$'s and $\varepsilon_p$'s, 
\begin{itemize}
    \item for every $p,q\in\mathbb N_0$ and every $n\in A_p$, $m\in A_q$ with $n\neq m$ we have that $|n-m|>p+q.$
    \item for every $p\in \zN_0$, $\|\sum_{n\in A_p} e_{n+N_p}\|\leq\frac{1}{(p+1)2^{p+1}}$ and
    \item $(A_p)_p$ is compatible.
\end{itemize}

Let $A=RAT((A_p)_p)=\{N_j\}$.
The above conditions  imply that the vector $z=\sum_{p=0}^\infty\sum_{n\in A_p} \sum_{j=0}^{p} e_{n+N_j}=\sum_{n\in A} e_n$ is in $X$. 

Let $(y^p)_p$ be a dense sequence in $\mathfrak B_1(A)$ such that $supp(y^p)=\{0,N_1,\ldots, N_p\}.$


Consider now the vector $y=\sum_{p=0}^\infty \sum_{n\in A_p} \sum_{j=0}^p y^p_j e_{n+N_j}.$ This vector is well defined, because 
\begin{align*}
\sum_{p=0}^\infty \|\sum_{n\in A_p} \sum_{j=0}^p y^p_j e_{n+N_j}\|&\leq \sum_{p=0}^\infty \frac{1}{ 2^{p+1}}<\infty.    
\end{align*}

We show now that $y$ is a frequently hypercyclic vector for $\mathfrak B_1(A)$. Let $q\in \zN_0$ and $m\in A_{m_q}$, then by condition $(2)$ $ii)$ of Theorem \ref{coro: first Characterization frequent recurrence} and unconditionality,
$$\|B^m(y)-y^{q}\|_q=\|\sum_{p=0}^\infty \sum_{n\in A_p,n>m}\sum_{j=0}^p y^p_j e_{n-m+N_j}\|_q\leq C_q\varepsilon_{m_q}.$$

\end{proof}
We finish the subsection with an application of Theorem \ref{thm: caract S_1(A)}. Recall that a set $F\subset X$ is said to be dense lineable provided that $F\cup\{0\}$ contains a subspace which is dense in $X$. In the context of linear dynamics, it is immediate that the set of (frequently) hypercyclic vectors is dense lineable, because if $x$ is a (frequently) hypercyclic vector then $\{P(T)(x):P \text{ is a polynomial}\}$ is a dense space of (frequently) hypercyclic vectors. However, the same argument breaks for recurrence. In \cite{LopMen25} the authors found a recurrent operator without a dense lineable set of recurrent vectors. It is still an open question whether the set of frequently recurrent vectors is dense lineable, see \cite[Problem 5.1 and Remark 5.3]{GriLopPerQuestions2} and \cite[Problem 4.2]{LopMen25}. We provide a positive answer for backward shift     operators (both for the unilateral and bilateral case). We note also that in \cite[Corollary 4.8]{GriLopPerQuestions2} it was shown that if $T$ supports a dense set of vectors such that $T^nx\to 0$ (in particular if $T$ is a unilateral backward shift) and if the family $\mathcal F$ is an upper right-invariant Furstenberg family $\mathcal F$, then the $\mathcal F$-recurrence of $T$ implies by \cite{BonGroLopPer22JFA} the $\mathcal F$-hypercyclicity of $T$ and in particular the dense lineability for the set of $\mathcal F$-recurrent vectors of $T$. 

\begin{theorem}\label{cyclic and frequently recurrent vector}
Let $X$ be a Fr\'echet sequence space with unconditional basis $(e_n)_{n\in\mathbb N_0}$. Then  $B$ is frequently recurrent if and only if $B$ has a cyclic and frequently recurrent vector.

In particular, if $B$ is frequently recurrent then
\begin{enumerate}
    \item the set of frequently recurrent vectors is dense lineable and moreover,
    %\item $B\times\cdots \times B:X^n\to X^n$ is frequently recurrent for every $n\in\mathbb N.$
    \item the set of frequently recurrent vectors of $B\times\cdots \times B:X^n\to X^n$ is dense lineable for every $n\in\mathbb N.$
\end{enumerate}
\end{theorem}
\begin{proof}
If $B$ supports a cyclic and frequently recurrent vector, then the operator is frequently recurrent by the zero-one law Theorem \ref{0-1 law frequent recurrence}.

Conversely suppose that $B$ is frequently recurrent. Then there are, by  Theorem \ref{thm: caract S_1(A)}, a set $A$, with $\underline{dens}(A)>0$, and $x\in X$ supported in $A$ such that $x$ is a frequently hypercyclic vector for $\mathfrak B_1(A)$. By the proof of Theorem \ref{thm: caract S_1(A)} we can suppose that $x\in \mathfrak{B}_1(A)$. In particular, $x$ is frequently recurrent. We claim that $x$ is also a cyclic vector.

Let $y=\sum_{j=0}^N y_j e_j$ and $\varepsilon>0$. For each $n\in \mathbb N$ consider $A_n:=A-n$ and let $n_1,\ldots, n_r$ such that $[0,N]\subseteq \bigcup_{l=1}^r A_{n_l}$. 
Thus, $y$ be can written as $\sum_{l=1}^r y^l$, where each $y^l$ is supported in $A_{n_l}$. Let {$M=\max_{1\le l\le r}\|y^l\|_\infty$}. We have then that for every $1\leq l\leq r$, $M B^{n_l}(x)$ is a frequently hypercyclic vector for $\mathfrak B_M(A_{n_l})$, the set of vectors supported in $A_{n_l}$ with coordinate less than $M$. Thus, there are $k_1,\ldots, k_r$ such that for every $1\leq l\leq r$, $\|B^{k_l}(MB^{n_l}(x))-y^l\|<\frac{\varepsilon}{r}.$
It follows that $\|M\sum_{l=1}^r B^{n_l+k_l}(x)-y\|<\varepsilon.$

To prove $(1)$ just observe that $\{P(B)(x):P\in \zK[x]\}$ is a dense subspace formed by frequently recurrent vectors.
On the other hand it is well known that if an operator admits a cyclic and frequently recurrent vector then the $n$-fold product of the operator is frequently recurrent and with a dense lineable set of frequently recurrent vectors, see  \cite[Proposition 3.8]   {CardeMurBlock} or \cite[Proposition 4.2]{GriLopPerQuestions2}.
 \end{proof}

% {\color{red}
% Related to lineability there are two concepts that one may wonder if the set of frequently recurrent vectors satisfy: spaceability (i.e. the existence of a closed infinite dimensional subspace) and algebrability (i.e. the existence of a non-trivial algebra) of the set of frequent recurrent vectors. We refer to \cite{bayart2019hypercyclic,bayart2021baire,bes2020algebrable,falco2020algebrability,falco2020algebras,lopez2024recurrent,menet2014hypercyclic,menet2015existence,montes1996banach} for some related literature. 

% Concerning spaceability, by the results in \cite{lopez2024recurrent} even the set or recurrent vectors is not spaceable in general: it was shown in \cite{lopez2024recurrent} that the set of recurrent vectors of a weakly-mixing operator acting on a Banach space is spaceable if and only if the set of hypercyclic vectors is spaceable. Thus, we may consider the Rolewicz operators $\lambda B$ on $c_0$ or $\ell_p$ with $|\lambda|>1$ for which the set of hypercyclic vectors was shown not to be spaceable in \cite{montes1996banach}.

% Concerning algebrability of Rolewicz operators on $\ell_p$ or $c_0$ (considering the coordinatewise product) the set of hypercyclic vectors (and hence also of recurrent vectors) is algebrable \cite[Corollary 2.8]{falco2020algebrability}, but the set of frequently hypercyclic vectors for $\lambda B$ is not algebrable \cite[Corollary 1.3]{falco2020algebras}. We show next, based on the ideas of \cite{falco2020algebras} that the set of frequently recurrent vectors for $\lambda B$ is also not algebrable. 
% \begin{proposition}
%     Let $X=c_0$ or $X=\ell_p$ considered as an algebra with the coordinatewise product and let $|\lambda|>1$. Then there is no nonzero algebra contained on the set of frequently recurrent vectors of $\lambda B$.
% \end{proposition}



% \begin{proof}
%     Let us assume that there is a nonzero frequently recurrent  vector $x$  such that all the vectors in the algebra generated by $x$ are frequently recurrent.

% Let $0<\varepsilon<\frac1{2}$.
% Then the vector $z:=2\dfrac{x^p}{\|x\|_\infty^p}$ is frequently recurrent for each $p\in\mathbb N$ and, for $p$ big enough we have,
% \begin{align}\label{eq: algeb z~2}
% \left\|z-w\right\|<\frac{\varepsilon}{2},    
% \end{align}
% where $w=2e_{l_1}+\dots+2e_{l_N}$, and the $l_k$'s are the coordinates at which $x$ attains its maximum.   

% We will show that, for big enough $m,$ $z^m$ cannot be a recurrent vector. This is a contradiction because $z^m$ belongs to the algebra generated by $x.$

% Since $z$ is frequently recurrent, the set 
% \begin{align}\label{eq: algeb n_k}
%     (n_k)_k=\{n\in\mathbb N\,:\, \|(\lambda B)^{n}z-w\|<\varepsilon \}\supset \{n\in\mathbb N\,:\, \|(\lambda B)^{n}z-z\|<\frac{\varepsilon}{2} \} 
% \end{align} 
% has positive lower density and thus the sequence $(\frac{n_{k+1}}{k})_k$ is bounded. This implies that there is some $M>0$ such that $n_{k+1}< Mk\le Mn_k$ for every $k$. Let $m\ge M$, and assume that $z^m$ is recurrent. Then there is some $n>n_1$ such that 
% \begin{align*}
%     \frac{\varepsilon}{2}>\|(\lambda B)^{n}z^m-z^m\|\ge |\lambda^nz_{n+l_N}^m-z_{l_N}^m|\ge (2-\frac{\varepsilon}{2})^m-|\lambda^nz_{n+l_N}^m|,
% \end{align*}
% where the last inequality follows from the fact that by \eqref{eq: algeb z~2}, we have $2-\frac{\varepsilon}{2}<|z_{l_N}|$. Thus,
% \begin{align*}
%     |z_{n+l_N}|>\frac{\Big((2-\frac{\varepsilon}{2})^m-\frac{\varepsilon}{2}\Big)^{1/m}}{|\lambda|^{n/m}}> |\lambda|^{-n/m}.
% \end{align*}
% Let $k$ be such that $n_k<n\le n_{k+1}< Mn_k \le m n_k$. Then by \eqref{eq: algeb n_k}, 
% \begin{align*}
%     \varepsilon > &\|(\lambda B)^{n_k}z-w\| \ge \left|\Big[(\lambda B)^{n_k}z-w\Big]_{n+l_N-n_k}\right|= \left|\lambda^{n_k} z_{n+l_N}-w_{n+l_N-n_k}\right| \\ 
%     = & \left|\lambda^{n_k} z_{n+l_N}\right|  > |\lambda|^{n_k-\frac{n}{m}}> 1
% \end{align*}
% \end{proof}
% }

\subsection{Two intermediate conditions\label{intermediate}}
There are two natural cases in which the backward shift operator satisfies conditions  that are formally weaker than frequent hypercyclicity but stronger than frequent recurrence. The first  one is to consider the operators for which the equivalence (2) of Theorem \ref{coro: first Characterization frequent recurrence} holds for $A=\mathbb N_0$.
 Theorem \ref{thm: caract S_1(A)} implies that, in this case, the backward shift operator is frequently hypercyclic for $\mathfrak B_1=\mathfrak B_1(\mathbb N_0)$, the set of vectors with coordinates smaller than one.
\begin{corollary}\label{Teo: caract freq hyp on S1}
    Let $X$ be a Fr\'echet sequence space with  unconditional basis $(e_n)_{n\ge 0}$. The following are equivalent:
\begin{enumerate}
    \item There are sets $A_p$, $\underline {dens}(A_p)>0$ and a sequence $\varepsilon_p\to 0$ such that
    \begin{enumerate}
        \item [i)] for every $p\ge 0,$ $\sum_{n\in A_p} e_{n+p}\in X$ and
        \item [ii)] for every $q\ge 0$ and $m\in A_q$, $\|\sum_{p=0}^\infty \sum_{n\in A_p,n>m}\sum_{j=0}^p e_{n-m+j}\|_q<\varepsilon_q.
    $    
        \end{enumerate}
\item $B$ is frequently hypercyclic for $\mathfrak B_1$.


\end{enumerate}
\end{corollary}

In a similar way to Theorem \ref{cyclic and frequently recurrent vector} we have the following. 
\begin{proposition}
     Let $X$ be a Fr\'echet sequence space. If  $B$ is frequently hypercyclic for $\mathfrak B_1$ then $B$ supports a supercyclic and frequently recurrent vector.
 \end{proposition}
\begin{proof}
If $B$ is frequently hypercyclic for $\mathfrak B_1$, then the proof of Theorem \ref{thm: caract S_1(A)} shows that there is a vector $x\in \mathfrak B_1$, which is frequently hypercyclic for $\mathfrak B_1$. This vector is by definition, frequently recurrent and supercyclic.
\end{proof}


The other case that is worth to consider is when
condition ii) of equivalence (2) in Theorem \ref{Characterization frequent recurrence} is smaller than $\min\{\varepsilon_p,\varepsilon_q\}$. Assuming that $X$ is Banach we obtain an equivalent notion to frequent hypercyclicity.

\begin{proposition}\label{prop: condicion con el min}
    Let $X$ be a Banach space with unconditional basis $(e_n)_{n\ge 0}$. Suppose that there are a subsets $A_p\subseteq \zN$ of positive lower density, a sequence $(N_k)_k$ defined by the relation $N_k=\min_{n>N_{k-1}} \{n\in \bigcup_{p=0}^{k-1}\bigcup_{j=0}^p A_p+N_j\}$ and $\varepsilon_p\to 0$, such that
\begin{enumerate}
    \item [i)]for every $p\in\mathbb N_0$, $\sum_{n\in A_p}e_{n+p}$ converges 
    \item [ii)]
    For every $q\in\mathbb N_0$ and every  $m\in A_q$,
    $$\|\sum_{p=1}^\infty \sum_{n\in A_p, n>m} \sum_{j=0}^p e_{n-m+N_j}\|<\min\{\varepsilon_p,\varepsilon_q\}.$$
   \end{enumerate}
   Then $B$ is frequently hypercyclic.
\end{proposition}
\begin{proof}
By Theorem \ref{caracterizacion hiperciclicidad frecuente} it suffices to prove that there are $A_p$ of positive lower density satisfying i) and the stronger condition
$$\|\sum_{n\in A_p,n>m} e_{n-m+j}\|<\min\{\varepsilon_p,\varepsilon_q\}.$$
We will prove that this holds for a subsequence of $A_p$.

Indeed, considering a subsequence of $A_p$ and using that the basis is unconditional, we may suppose $C\|B\|^{N_p}\varepsilon_p<\frac{1}{2^p}$. Hence, if $p,q\in\zN$ and $m\in A_q$ then
\begin{align*}
  \|\sum_{n\in A_p, n>m} e_{n-m+j}\|&=\|B^{N_p-j}\left(\sum_{n\in A_{p},n>m} e_{n-m+N_p}\right)\|\\
  &\leq C\|B\|^{N_p}\|\sum_{p=1}^\infty \sum_{n\in A_p,n>m}\sum_{j=0}^p e_{n-m+N_j}\|\\
  &\leq C\|B\|^{N_p}\min\{\varepsilon_p,\varepsilon_q\}\leq \min \{\frac{1}{2^p},\frac{1}{2^q}\}.
\end{align*}
\end{proof}





\section{Bilateral backward shifts}\label{section:bilateral}
We now focus  our attention to bilateral backward shifts and extend the results of the preceding section to the bilateral setting. Although all the results in Section \ref{section unilateral} have an extension to the bilateral case, we choose to mention only the main ones.
\begin{theorem}\label{0-1 law frequent recurrence bil}
Let $X$ be a Fr\'echet sequence space with uncondtional basis $(e_n)_{n\in\mathbb Z}$. If $B$ supports a non-trivial frequently recurrent vector then $B$ is frequently recurrent.
\end{theorem}
%The proof is similar to the proof of Theorem \ref{0-1 law frequent recurrence} with some modifications. We only give a sketch of it. 


\begin{proof}
 Let $\varepsilon>0$, $R\in\mathbb Z$, we show that there is a frequently recurrent vector $y\in X$ such that $\|y-e_R\|<\varepsilon$.

Let $x\neq 0$ be a frequently recurrent vector such that $x_R\neq 0$. Let $\delta\in(0,1)$ such that $|x_R|>\delta.$

Since $x$ is a recurrent vector,  there is for every $p\in\mathbb N_0$ a  $q_p>p+|R|$ such that $|x_{q_p}|>\delta$. By the continuity of the coordinate functionals, there are a decreasing sequence of positive numbers $(\varepsilon_{p})_p$ and a sequence of natural numbers $(k_p)_p$  satisfying that for each $p\in\mathbb N_0$, $\varepsilon_{p}<\frac{\varepsilon}{2^{p+1}}$, $k_p>p+1$, and such that  
\begin{align}\label{eq: beta_p (zero one freq) BILAT}
 \begin{split}
  \|w\|_{k_p}<\varepsilon_{p}\Longrightarrow  |w_{j}|<\frac{\delta}{4} ,  \text{ for every } |j|\leq q_p, &\quad \text{and}\\
 \|x-\sum_{|j|\le k_p} x_j e_j\|_{m_p}<\frac{\varepsilon}{2^{p+1}}\frac{\delta}{10C_p},&
 \end{split}
 \end{align}
 where $C_p$ and $\|\cdot\|_{m_p}$ are the unconditional constant basis and seminorm associated to $\|\cdot\|_p$. 



 
 

For each $p\ge 0$ we consider 
\begin{equation}\label{bilateral eq: A_p zero one freq}
A_p\subseteq N_B(x,U_p) \text{ where } U_p=\{y\in X:\|y-x\|_{k_p+m_p}<\frac{\delta\varepsilon_p}{10C_p}\},
\end{equation}
and the $A_p$ have positive lower density.
Applying the Separation Lemma \ref{separation lemma}
we may assume that the sets $A_p$ are pairwise disjoint and that for every $n\in A_p,\, m\in A_q$ with $n\neq m$ we have that
\begin{equation}\label{bilateral eq separation zero one freq}
|n-m|>k_p+k_q+2|R|>p+1+q+1+2|R|.    
\end{equation}

 By \eqref{eq: beta_p (zero one freq) BILAT} and \eqref{bilateral eq: A_p zero one freq} we have that $|x_{n+q_p}-x_{q_p}|<\frac{\delta}{4}$ and hence $|x_{n+q_p}|>\frac{\delta}{2}$ for each $n\in A_p.$ By unconditionality this implies that $\sum_{n\in A_p} e_{n+q_p}$ converges and by applying the backward shift $q_p-p$ times it follows that $\sum_{n\in A_p} e_{n+p+1}$ converges. By the finite invariance property of  the sets of positive lower density, we may assume that $A_p\sub [k_p+|R|+1,\infty)$ and that,  for every $|j|\le p+1,$
\begin{equation}\label{bilateral eq: sumas de e_j norma chica (zero one freq)}
 \|\sum_{n\in A_p} e_{n+j}\|<\frac{\varepsilon}{(p+1)2^{p+1}}.   
\end{equation}

Applying Proposition \ref{prop: condic necesaria vectores recurrentes} and Remark \ref{rem:M_k>N_k} there is a sequence $(M_p)_p$ for which 
\begin{equation}\label{eq: bilateral rat en U_0}
    RAT((A_{M_p})_p)\subseteq N_B(x,U_0)
\end{equation} and such that $M_p>N_p$ for every $p\geq 1$. The fact that $M_p>N_p$ together with the contention $A_p\sub [k_p+|R|+1,\infty)$, implies that $(A_{M_p})_p$ is compatible, i.e. 
 $RAT((A_{M_p})_p)=\{N_k\}$.

Let $$y=\sum_{n\in RAT((A_{M_p})_p)} e_{n+R}=e_R+\sum_{p=0}^{\infty}\sum_{n\in A_{M_p}}\sum_{j=0}^pe_{n+N_j+R},$$
which we claim to be frequently recurrent.

Using \eqref{bilateral eq: sumas de e_j norma chica (zero one freq)} and applying the same arguments as in \eqref{eq: def y^{k+1}} it can be shown that $y\in X$ and that $\|y-e_R\|<\varepsilon$.

We will show now that $y$ is frequently recurrent. 
 Fix $q\in\zN$ and $m\in A_{M_{q}}$.
 
Noting that $m$ belongs only to $A_{M_{q}}$, because the sets $A_{M_p}$ are pairwise disjoint, we get that
\begin{align*}
B^m(y)-\sum_{j=-|R|}^{N_{q}+|R|}y_je_j &=B^m(y)-\sum_{j=0}^{q}e_{N_j+R}\\
&=e_{R-m}+\sum_{p\neq q}\sum_{n\in A_{M_p}}\sum_{j=0}^{p}e_{n-m+N_j+R}+\sum_{n\in A_{M_{q}}}\sum_{j=0}^{q}e_{n-m+N_j+R}-\sum_{j=0}^{q}e_{N_j+R}\\
&=e_{R-m}+\sum_{p=0}^\infty \sum_{n\in A_{M_p},n\neq m}\sum_{j=0}^pe_{n-m+N_j+R}.
\end{align*}

% Thus,
% \begin{equation}\label{eq norma frec zero one bilateral}
%  \|B^m(y)-\sum_{j=-|R|}^{N_{s_0}+|R|}y_je_j \|\le C(2|R|+1)\|B\|^{2|R|}\|z\|_\infty\left\|\sum_{s=0}^\infty \sum_{n\in A_{N_s},n\neq m}\sum_{l=0}^se_{n-m+N_l+R}\right\|.
% \end{equation}
By \eqref{eq: bilateral rat en U_0} we have that $RAT((A_{M_p})_p)\subseteq N_B(x,U_0)$. Thus, we have that for every $n\in A_{M_p}$ and every $0\le j\le p,$ that
$$
B^{n+N_j}x\in U_0,\quad\text{and in particular,}\quad \| B^{n+N_j}x -x\|_{k_0+m_0}<\varepsilon_0.
$$
Then since $|x_R|>\delta,$ we have by \eqref{eq: beta_p (zero one freq) BILAT} that 
\begin{equation*}
    |x_{n+N_j+R}|\ge |x_R|-|x_{n+N_j+R}-x_R| \ge \delta - \frac{\delta}{4}>\frac{\delta}{4}.
\end{equation*}


 
 Moreover, since by \eqref{bilateral eq separation zero one freq}, $|n-m|>N_p+k_{N_{q}}+|R|$ whenever $n\in A_{M_p}$,  it follows that the $|n-m+N_j+R|>k_{N_{q}}$ for each $0\le j \le p$. Note also that since $\min A_{M_q}>k_{M_q}+|R|$ we have $R-m<-k_{M_q}\leq- k_{N_q}$, and thus,
  \begin{align*}
 \begin{split}
  \|B^m(y)-\sum_{j=-|R|}^{N_{q}+|R|}y_je_j \|_{q}&=\left\|e_{R-m}+\sum_{p=0}^\infty \sum_{n\in A_{M_p},n\neq m}\sum_{j=0}^pe_{n-m+N_j+R}\right\|_{q}\\
% &=
%  \left\|\sum_{s=0}^\infty \sum_{n\in A_{N_s},n\neq m}\sum_{l=0}^se_{n-m+N_l+R}\right\|
&= \left\|\frac1{x_{R}}x_{R}e_{R-m}+\sum_{p=0}^{\infty}\sum_{n\in A_{M_p},n\neq m}\sum_{j=0}^p\frac1{x_{n+N_j+R}}x_{n+N_j+R}e_{n-m+N_j+R}\right\|_{q}\\
 &\le \left\|\frac1{x_{R}}[B^m(x)-\sum_{j=-k_{N_{q}}}^{k_{N_{q}}}x_je_j]_{R-m}e_{R-m}\right\|_q\\
 &+\left\|\sum_{p=0}^{\infty}\sum_{n\in A_{M_p},n\neq m}\sum_{j=0}^p\frac1{x_{n+N_j+R}}[B^m(x)-\sum_{j=-k_{N_{q}}}^{k_{N_{q}}}x_je_j]_{n-m+N_j+R}e_{n-m+N_j+R}\right\|_q\\
 &\leq \Big(\frac{C_q}{\delta}+\frac{4C_q}{\delta}\Big)\left(\|B^m(x)-x\|_{m_q}+\|x-\sum_{j=-k_{N_{q}}}^{k_{N_{q}}}x_je_j\|_{m_q}\right)< \frac{\varepsilon}{2^{{q}}},
  \end{split}
 \end{align*}
where the last inequality follows from \eqref{eq: beta_p (zero one freq) BILAT} and \eqref{bilateral eq: A_p zero one freq}.
 This implies that $y$ is a frequently recurrent vector.
 
\end{proof}



In the same way as in the unilateral case we can deduce a characterization for  bilateral frequently recurrent backward shifts.
 
\begin{theorem}\label{Characterization frequent recurrence bil}
Let $X$ be a Banach space with unconditional basis $(e_n)_{n\in\mathbb Z}$. The following are equivalent: 
\begin{enumerate}
\item $B$ is frequently recurrent and
\item
 there are a sequence $\varepsilon_p\to 0$ and sets of positive lower density $A_p$ such that
\begin{enumerate}
    \item [i)]for every $p\in\mathbb N_0$, $\sum_{n\in A_p}e_{n+p}$ converges 
\end{enumerate}
and such that one of the following conditions holds
\begin{enumerate}
    
    \item [ii)]For every $q\in \zN$ and every $m\in A_q$ we have that
    $$\Big\|\sum_{n\in RAT((A_p)), n\neq m} e_{n-m}\Big\|_q<\varepsilon_q;$$
    
    
\end{enumerate}
\item There is a compatible sequence $(A_p)_p$ formed by sets of positive lower density such that
$$\sum_{n\in RAT((A_p)_p)} e_n \text{ is a frequently recurrent vector}.$$
\end{enumerate}

\end{theorem}
\begin{remark}\label{rem: condition ii' bil}
Condition $ii)$ can also  be replaced by:
    \begin{itemize}
    \item [ii')]
     If $N_j$ is the sequence defined by $N_k=\min_{n>{N_{k-1}}} \bigcup_{p=0}^{k-1} \bigcup_{j=0}^p A_p+N_j$ and $N_0=0$, then 
    for every $q\in\mathbb N_0$ and every  $m\in A_q$ we have that
    $$\|\sum_{p=0}^\infty \sum_{n\in A_p, n\neq m} \sum_{j=0}^p e_{n-m+N_j}\|_q<\varepsilon_q$$
    
    \end{itemize}
    
\end{remark}

\begin{proof}[Sketch of the proof of Theorem \ref{Characterization frequent recurrence bil}.]
As in the unilateral case, the implications $(1)\Rightarrow (2)$ and $(1)\Rightarrow (3)$ can be extracted from the proof of  Theorem \ref{0-1 law frequent recurrence bil}. By the zero-one law Theorem for bilateral shifts we have that $(3)\Rightarrow (1).$

(2)$\Rightarrow$ (1).  
 By Theorem \ref{0-1 law frequent recurrence bil} it suffices to construct a nonzero frequently recurrent vector.
Let $(N_j)_j$ be the sequence associated to $RAT((A_p)_p)$ (Or just $N_j$ if we are assuming ii')). 
 
 Applying the Separation Lemma \ref{separation lemma} and finite invariance there are subsets $A_p'\subseteq A_p$ such  that,  for every $p,q\in\zN$, every $n\in A_p'$ and $m\in A_q'$, and every $0\leq j\leq p$, 
\begin{equation}\label{separation property characterization bil}
|n-m|>p+q,\qquad \text{
 $A_p'\subseteq [N_p+1,\infty)$, }\qquad \|\sum_{n\in A_p'} e_{n+j}\|\leq \frac{1}{2^p}.    
\end{equation}


We apply now Lemma \ref{lem: RAT subconjuntos} to obtain a sequence $(M_p)_p$ such that $(A_{M_p}')_p$ is compatible and $RAT((A_{M_p}')_p)\sub RAT((A_p)_p)$.
 Moreover, if $RAT((A_{M_p}')_p)=(N_j')_j$, $N_{l}'=\min_{n>N_{l-1}'} \bigcup_{p=0}^{l-1} \bigcup_{j=0}^p A_{M_p}'+N_j'$ for every $l$, we may assume that 
\begin{enumerate}
    %\item[a)] $N_{l}'=\min_{n>N_{l-1}'} \bigcup_{p=0}^{l-1} \bigcup_{j=0}^p A_{M_p}'+N_j'$ for every $1\leq l\leq k-1$,
   \item[a)] $M_{l}>\max\{m_{l},N_l'\}$ for every $1\leq l$, (see Remark \ref{rem: cosas del lemma que usamos})  %\label{eq: caracterizacion 1 Los M_l son grandes}
   \item [b)] $\varepsilon_{M_l}<\varepsilon_{M_{l-1}}$ for $1\leq l$, 
    \item [c)] $C_l\cdot\varepsilon_{M_l}<\frac{1}{l}$ for every $1\leq l$,
    \item [d)] $\{N_1',\ldots N_k'\}\subseteq \{N_1,\ldots N_k\}$,
\end{enumerate}
where $C_l,m_l$ are the constants from the unconditionality.

Either if ii) or ii') holds, we have by unconditionality that
\begin{align}
\nonumber \left\|\sum_{\overset{N\in RAT((A_{M_p}')_p)}{N\neq m}} e_{N-m}\right\|_q & =    \left\|\sum_{p=0}^\infty\sum_{\overset{n\in A_{M_{p}}'}{ n\neq m}}\sum_{j=0}^p e_{n-m+N_j'}\right\|_q \leq C_q\cdot \left\|\sum_{p=0}^\infty\sum_{\overset{n\in A_{M_p}}{n\neq m}}\sum_{j=0}^{M_p} e_{n-m+N_j}\right\|_{m_q}\\
 \label{eq: norma de la suma n>m <1/q BIL}   & \leq C_q\cdot\left\|\sum_{p=0}^\infty\sum_{\overset{n\in A_{M_p}}{n\neq m}}\sum_{j=0}^{M_p} e_{n-m+N_j}\right\|_{M_q}
     <C_q\cdot \varepsilon_{M_{q}}%<C_q\cdot \varepsilon_{M_q}
    <\frac{1}{q}. 
\end{align}
Let $y=\sum_{p=0}^\infty \sum_{n\in A_{M_p}'}\sum_{j=0}^p e_{n+N_j'}.$ This vector is well-defined by \eqref{separation property characterization bil} and it is frequently recurrent because if $m\in A'_{M_q}$ then
 $$B^m(y)-\sum_{j=0}^{N_q'} y_j=B^m(y)-\sum_{j=0}^q e_{N_{j}'}= \sum_{p=0}^\infty \sum_{n\in A_{M_p'}, n\neq m} \sum_{j=0}^p e_{n-m+N_j'}$$
 and,   by \eqref{eq: norma de la suma n>m <1/q BIL},
 $$
 \Bigg\|\sum_{p=0}^\infty \sum_{n\in A_{M_p'}, n\neq m} \sum_{j=0}^p e_{n-m+N_j'}\Bigg\|_q<\frac{1}{q}.
 $$
%  We follow the ideas of the previous Theorem to construct a frequently recurrent vector near $y^0=e_R$, where $R=N_0'=\min\{n\in A_0\}$.


% Let $p_0>N'_0=R$ and let $y^1=y^0+\sum_{n\in A_{p_0}} e_{n}$. Consider $N'_1=\min\{n\in A_{p_0}\}$ and $p_1>N'_1$. In particular, $N_{p_1}>N'_1$. Applying $p_1-N'_1$ times the backward shift we obtain that $\sum_{n\in A_{p_1}} e_n +e_{n+N'_1}$ converges  and by the finite invariance property, we may suppose that $\|\sum_{n\in A_{p_1}} e_n+e_{n+N'_1}\|<\frac{1}{4}$.
% Let
% $$y^2=y^1+\sum_{n\in A_{p_1}}\sum_{j=0}^{N'_1}y^1_j e_{n+j}=y^1+\sum_{n\in A_{p_1}} e_n+\sum_{n\in A_{p_1}}e_{n+N'_1}.$$

% If $y^1,\dots,y^k$,  $p_0,\dots, p_{k-1}$ and $N_0',\dots,N_{k-1}'$ have been constructed, let $N'_k=\min \{n\in \bigcup_{0\le j\le k-1} A_{p_j}:n>N'_{k-1}\}$ and let $p_k>N'_k$ (in particular, $N_{p_k}>N'_k$). Therefore,
% $\sum_{n\in A_{p_k}
% }e_{n+p_k}$ converges and applying the backward shift,
% $\sum_{n\in A_{p_k}}e_{n+N'_j}$ converges for every $0\leq j\leq k$.
% Without loss of generality, we may suppose that 
% \begin{align}\label{eq: suma en A_p chica bil}
% \|\sum_{n\in A_{p_k}}e_n+e_{n+N'_1}+\ldots +e_{n+N'_k}\|<\frac{1}{2^{k+1}}.    
% \end{align}
% We define
% $$
% y^{k+1}=y^k+\sum_{n\in A_{p_k}} \sum_{j=0}^{N'_k}y^k_j e_{n+j}=y^k+\sum_{n\in A_{p_k}} \sum_{j=0}^{k} e_{n+N'_j}=\sum_{l=0}^k \sum_{n\in A_{p_l}}\sum_{j=0}^l e_{n+N'_j},
% $$
% and
% $$
% y=\lim_k y^k=\sum_{l=0}^\infty \sum_{n\in A_{p_l}}\sum_{j=0}^l e_{n+N'_j}=\sum_{l=0}^\infty e_{N_l'},
% $$ 
%  where the last equality holds since by the definition of $N_k'$ it follows that $\bigcup_{l\ge 0}A_{p_l}=\{N_j':j\ge 0\}$.
%  Note that $y$ is a well defined nonzero vector by \eqref{eq: suma en A_p chica bil}.
 
%  We claim that $y$ is a frequently recurrent vector. 

% Let $s\in\zN$ and ${m_{s}}$ and $C_{s}$ be the associated numbers to $s$ according to the unconditionality of the basis (see \eqref{incondicionalidad}). Let $m\in A_{p_{m_{s}}}$, then
% \begin{align*}
% B^m(y)-\sum_{j=0}^{N'_{m_s}}y_je_ j
% &=\sum_{k\ne m_s}\sum_{n\in A_{p_k}}\sum_{j=0}^k e_{n-m+N'_j}+\sum_{n\in A_{p_{m_s}}}\sum_{j=0}^{m_s} e_{n-m+N'_j}-\sum_{j=0}^{{m_s}} e_{N'_j}\\
% &=\sum_{k=0}^\infty\sum_{n\in A_{p_k},n\ne m}\sum_{j=0}^k e_{n-m+N'_j}.
% \end{align*}
% Hence noting that $\{N'_j\}\subseteq \{N_j\}$ we obtain by the unconditionality of the basis 
%  and by $ii)$  that 
% \begin{align*}\label{aproxima bien characterization bil}
% \|B^m(y)-\sum_{j=0}^{N'_{m_s}}y_je_ j\|_{s}
% %&=\|\sum_{k=0}^\infty\sum_{n\in A_{p_k}}\sum_{j=0}^k e_{n-m+N'_j}-\sum_{j=0}^{k_0} e_{N'_j}\|_{k_0}\\
% & =\|\sum_{k=0}^\infty \sum_{n\in A_{p_k}, n\ne m} \sum_{j=0}^{k} e_{n-m+N'_j}\|_{s}\le  C_s  \|\sum_{p=0}^\infty \sum_{n\in A_p, n\ne m} \sum_{j=0}^p e_{n-m+N_j}\|_{m_s}<C_{s}\varepsilon_{p_{m_{s}}}.
%  %\leq \|\sum_{p=0}^\infty \sum_{n\in A_p, n \ne m} \sum_{j=0}^p e_{n-m+N_j}\|_{p_{m_{s}}}
% \end{align*}
% Since $\sum_{j=0}^{N'_{m_s}}y_je_ j$ converges to $y$, this shows that $y$ is a frequently recurrent vector.
\end{proof}


Proceeding as in Theorems  \ref{thm: caract S_1(A)} and \ref{cyclic and frequently recurrent vector}, but using now the zero-one law from Theorem \ref{0-1 law frequent recurrence bil} and the characterization for bilateral backward shifts in Theorem \ref{Characterization frequent recurrence bil}, we can also prove the following results.   
\begin{theorem} \label{thm: caract S_1(A) bilat}
Let $X$ be a Fr\'echet sequence space with unconditional basis $(e_n)_{n\in \mathbb Z}$. Then $B$ is frequently recurrent if and only if there are $A\subseteq \mathbb N_0$ with $\underline {dens}(A)>0$ and $x\in X$ supported in $A$ such that $x$ is a frequently hypercyclic vector for $\mathfrak B_1(A)$. 
\end{theorem}

% \begin{proof}
% If there are $A$ and $x\in X$ supported in $A$ such that $x$ is frequently hypercyclic for $\mathfrak B_1(A)$ then we have by unconditionality that $P_{\mathfrak B_1(A)} x:= \sum_{n\in A} \min\{|x_n|,1\}sgn(x_n) e_n$ is frequently hypercyclic for $\mathfrak B_1(A)$. Indeed, for any $y\in \mathfrak B_1(A)$, $q\in \mathbb N_0$, we have,
% $$
% \|B^m(P_{\mathfrak B_1(A)} x)-y\|_q\le C_q \|B^m( x)-y\|_{m_q}.
% $$


% Since $P_{\mathfrak B_1(A)} x$ belongs itself to $\mathfrak B_1(A)$ then it is a frequently recurrent vector. By the zero-one law Theorem \ref{0-1 law frequent recurrence bil} we deduce that the operator is frequently recurrent.  

% Suppose now that $B$ is frequently recurrent. Let  $A=\{N_j\}$, $A_p$ and $\varepsilon_p$ given by Theorem \ref{Characterization frequent recurrence bil}.

% By the Separation Lemma \ref{separation lemma} we may suppose that for every $p,q\in\mathbb N_0$ and every $n\in A_p$, $m\in A_q$ with $n\neq m$ we have that $|n-m|>p+q.$
% By considering a subsequence of $A_p$ and by finite invariance, we may suppose that for every $p\in \zN_0$, $\|\sum_{n\in A_p} e_{n+N_p}\|\leq\frac{1}{(p+1)2^{p+1}}$. This implies that the vector $z=\sum_{p=0}^\infty\sum_{n\in A_p} \sum_{j=0}^{N_p} e_{n+N_j}$ is in $X$. Let $(y^p)_p$ be a dense sequence in $\mathfrak B_1(A)$  such that $supp(y^p)=\{0,N_1,\ldots, N_p\}.$


% Consider now the vector $y=\sum_{p=0}^\infty \sum_{n\in A_p} \sum_{j=0}^p y^p_j e_{n+N_j}.$ By unconditionality this vector is well defined. We see now that $y$ is a frequently hypercyclic vector for $\mathfrak B_1(A)$. Let $q\in \zN_0$ and $m\in A_{m_q}$, then by condition $(2)(ii)$ of Theorem \ref{Characterization frequent recurrence bil} and unconditionality,
% $$\|B^m(y)-y^{q}\|_q=\|\sum_{p=0}^\infty \sum_{n\in A_p,n\ne m}\sum_{j=0}^p y^p_j e_{n-m+N_j}\|_q\leq C_q\varepsilon_{m_q}.$$

% \end{proof}

\begin{theorem}\label{cyclic and frequently recurrent vector bil}
Let $X$ be a Fr\'echet sequence space with unconditional basis $(e_n)_{n\in\mathbb Z}$. Then  $B$ is frequently recurrent if and only if there exists a cyclic and frequently recurrent vector.

In particular, if $B$ is frequently recurrent then
\begin{enumerate}
    \item the set of frequently recurrent vectors is dense lineable and moreover,
    %\item $B\times\cdots \times B:X^n\to X^n$ is frequently recurrent for every $n\in\mathbb N.$
    \item the set of frequently recurrent vectors of $B\times\cdots \times B:X^n\to X^n$ is dense lineable for every $n\in\mathbb N.$
\end{enumerate}
\end{theorem}
% The proof is identical to the proof of Theorem \ref{cyclic and frequently recurrent vector}.
% \begin{proof}
% If $B$ supports a cyclic and frequently recurrent vector, then the operator is frequently recurrent by the zero-one law Theorem \ref{0-1 law frequent recurrence bil}.

% Conversely suppose that $B$ is frequently recurrent. Then there are, by  Theorem \ref{thm: caract S_1(A) bilat}, a set $A\subseteq \mathbb N_0$, with $\underline{dens}(A)>0$, and $x\in X$ supported in $A$ such that $x$ is a frequently hypercyclic vector for $\mathfrak B_1(A)$. By the proof of Theorem \ref{thm: caract S_1(A)} we can suppose that $x\in \mathfrak{B}_1(A)$. Hence, $x$ is frequently recurrent. We claim that $x$ is a cyclic vector.

% Let $y=\sum_{j=-N}^N y_j e_j$ and $\varepsilon>0$. For each $n\in \mathbb N$ consider $A_n:=A-n$ and let $n_1,\ldots, n_r$ such that $[-N,N]\subseteq \bigcup_{l=1}^r A_{n_l}$. 
% Thus, $y$ be can written as $\sum_{l=1}^r y^l$, where each $y^l$ is supported in $A_{n_l}$. Let $M=\|y\|_\infty$. We have then that for every $1\leq l\leq r$, $M B^{n_l}(x)$ is a frequently hypercyclic vector for $\mathfrak B_M(A_{n_l})$, the set of vectors supported in $A_{n_l}$ with coordinate less than $M$. Thus, there are $k_1,\ldots k_r$ such that for every $1\leq l\leq r$, $\|B^{k_l}(MB^{n_l}(x))-y^l\|<\frac{\varepsilon}{r}.$
% It follows that $\|M\sum_{l=1}^r B^{n_l+k_l}(x)-y\|<\varepsilon.$

% To prove $(1)$ just observe that $\{P(B)(x):P\in \zK[x]\}$ is a dense subspace formed by frequently recurrent vectors.
% On the other hand it is well known that if an operator admits a cyclic and frequently recurrent vector then the $n$-fold product is frequently recurrent and with a dense lineable set of frequently recurrent vectors, see  \cite[Proposition 3.8]   {CardeMurfreqrecurrence_arxiv} or \cite[Proposition 5.9]{GriLop22arxiv}.
%  \end{proof}


\section{Final remarks}

An important open question, posed in \cite{BonGroLopPer22JFA}, is whether every frequently recurrent backward shift is indeed frequently hypercyclic. For boundedly complete bases, there is an easy positive answer: by \cite{charpentier2019chaos}, every reiteratively hypercyclic backward shift operator over a boundedly complete basis is frequently hypercyclic, and by \cite{CardeMurBlock} every reiteratively recurrent backward shift operator is reiteratively hypercyclic. However, the proof breaks down for more general spaces. 

Inspired by the results of Subsection \ref{intermediate} we can split this problem into two different questions.

\textbf{Question 1} Let $B$ be frequently hypercyclic for the set of coordinates bounded by $1$. Is $B$ frequently hypercyclic?

\textbf{Question 2} Let $B$ be a frequently recurrent backward shift operator. Is $B$ frequently hypercyclic for the set of coordinates bounded by $1$?



By \cite{GriLop23}, the existence of invariant measures is related to the existence of frequently recurrent vectors in the support of the measure.
It is thus natural to ask if a zero-one type law
holds for invariant measures of weighted shifts: does the existence of a non-trivial invariant measure imply the existence of an invariant measure of full support? 

If $B$ is a backwardshift operator acting on a Banach space $X$ in which $\{e_n\}$ is a boundedly complete and unconditional basis (in this case $B$ is an adjoint operator) then using the zero-one law we can give a positive answer and moreover
$B$ supports a non-trivial invariant measure  if and only if it has an ergodic measure of full support.
This also provides an alternative proof of  \cite[Proposition 2.1]{grivaux2024orthogonality}.

Indeed, according to \cite[Lemma 3.1]{GriLop23} the existence of non-trivial invariant measure implies the existence of a nonzero frequently recurrent vector. By the zero-one law Theorem \ref{0-1 law frequent recurrence} this implies that the operator is frequently recurrent and since $B$ is an adjoint operator then the operator is frequently hypercyclic and chaotic \cite[Theorem 2.1]{charpentier2019chaos}. In particular, it satisfies the frequent hypercyclicity criterion \cite{GroPer11}[See Proposition 9.13 and Theorem 4.8], which in turn implies ergodicity \cite[Theorem 1]{MurPer13}.

For non-adjoint operators or in the case of non-normable Fr\'echet spaces, the situation is more subtle. 
We thus pose the following question.

\textbf{Question 3.} Let $B$ be a backwardshift operator on a Fr\'echet (or Banach) sequence space $X$. Suppose that $B$ has a non-trivial invariant measure. Does $B$ have an invariant measure of full support?



We remark that for an adjoint operator operator on a Fr\'echet space $X$ such that its predual is a quasi-$\ell_\infty$-barreled Hausdorff locally convex topological vector space, sufficient conditions for the existence invariant measures of full support were proved in \cite[Theorem 2.3]{Lop24Locallybounded}, which are related to the existence of frequently recurrent vectors with locally bounded orbit.


Finally we note that the results presented here can be extended to $\mathcal F$-recurrence for families $\mathcal F$ which are finitely invariant and satisfy a separation lemma (like Lemma \ref{separation lemma}, see also \cite{MarMenPui22,BayGri04}). But, for many families, like upper or block families, the results are simpler to prove (see \cite{bonillazerofurstenberg_preprint}). 




\subsection*{Acknowledgments} We would like to express our gratitude towards the reviewer, who read carefully the entire manuscript and beside other helpful comments, pointed out a significant gap in the previous version. Their comments motivated us to analyze in detail the arithmetic structure of the return time sets.


 \bibliographystyle{abbrv} 
\begin{thebibliography}{10}
\bibitem{abakumov2024several}
E.~Abakumov and A.~Abbar.
\newblock On several dynamical properties of shifts acting on directed trees.
\newblock {\em Canadian Journal of Mathematics}, pages 1--30, 2024.

\bibitem{abakumov2024orbits}
E.~Abakumov and A.~Abbar.
\newblock Orbits of the backward shifts with limit points.
\newblock {\em Journal of Mathematical Analysis and Applications}, 537(2),
  2024.

\bibitem{BayGri04}
F.~Bayart and S.~Grivaux.
\newblock {Hypercyclicity: the role of the unimodular point spectrum.
  (Hypercyclicit\'e : Le r\^ole du spectre ponctuel unimodulaire.)}.
\newblock {\em C. R., Math., Acad. Sci. Paris}, 338(9):703--708, 2004.

\bibitem{BayGri06}
F.~{Bayart} and S.~{Grivaux}.
\newblock {Frequently hypercyclic operators}.
\newblock {\em {Trans. Am. Math. Soc.}}, 358(11):5083--5117, 2006.

\bibitem{BayGri07}
F.~Bayart and S.~Grivaux.
\newblock Invariant {G}aussian measures for operators on {B}anach spaces and
  linear dynamics.
\newblock {\em Proc. Lond. Math. Soc. (3)}, 94(1):181--210, 2007.

\bibitem{BayGriMatMen24}
F.~Bayart, S.~Grivaux, E.~Matheron, and Q.~Menet.
\newblock Hereditarily frequently hypercyclic operators and disjoint frequent
  hypercyclicity.
\newblock {\em Ergodic Theory and Dynamical Systems}, page 1–52, 2025.

\bibitem{BayMat09}
F.~Bayart and E.~Matheron.
\newblock {\em Dynamics of linear operators}, volume 179 of {\em Cambridge
  Tracts in Mathematics}.
\newblock Cambridge University Press, Cambridge, 2009.

\bibitem{BayRuz15}
F.~Bayart and I.~Z. Ruzsa.
\newblock Difference sets and frequently hypercyclic weighted shifts.
\newblock {\em Ergodic Theory Dynam. Systems}, 35(3):691--709, 2015.

\bibitem{Bes16}
J.~{B}\`es, Q.~{Menet}, A.~{Peris}, and Y.~{Puig}.
\newblock {Recurrence properties of hypercyclic operators}.
\newblock {\em {Math. Ann.}}, 366(1-2):545--572, 2016.

\bibitem{bonillazerofurstenberg_preprint}
A.~Bonilla, R.~Cardeccia, K.-G. Grosse-Erdmann, and S.~Muro.
\newblock Zero-one laws for weighted shifts under furstenberg families.
\newblock {\em In preparation}.

\bibitem{bonilla2025zeroEdinburgh}
A.~Bonilla, R.~Cardeccia, K.-G. Grosse-Erdmann, and S.~Muro.
\newblock Zero-one law of orbital limit points for weighted shifts.
\newblock {\em Proceedings of the Edinburgh Mathematical Society}, page 1–34,
  2025.

\bibitem{BonGro18}
A.~{Bonilla} and K.-G. {Grosse-Erdmann}.
\newblock {Upper frequent hypercyclicity and related notions}.
\newblock {\em {Rev. Mat. Complut.}}, 31(3):673--711, 2018.

\bibitem{BonGroLopPer22JFA}
A.~Bonilla, K.-G. Grosse-Erdmann, A.~L{\'o}pez-Mart{\'{\i}}nez, and A.~Peris.
\newblock Frequently recurrent operators.
\newblock {\em J. Funct. Anal.}, 283(12):36, 2022.
\newblock Id/No 109713.

\bibitem{CardeMur22}
R.~Cardeccia and S.~Muro.
\newblock Arithmetic progressions and chaos in linear dynamics.
\newblock {\em Integral Equations Oper. Theory}, 94(2):Paper No. 11, 18, 2022.

\bibitem{CardeMur22multipleScand}
R.~Cardeccia and S.~Muro.
\newblock Multiple recurrence and hypercyclicity.
\newblock {\em Math. Scand.}, 128(3):589--610, 2022.

\bibitem{CardeMurBlock}
R.~Cardeccia and S.~Muro.
\newblock Frequently recurrence properties and block families.
\newblock {\em Rev. R. Acad. Cienc. Exactas F\'{\i}s. Nat. Ser. A Mat. RACSAM},
  119(3):Paper No. 60, 27, 2025.

\bibitem{ChaSec12}
K.~Chan and I.~Seceleanu.
\newblock Hypercyclicity of shifts as a zero-one law of orbital limit points.
\newblock {\em J. Operator Theory}, 67(1):257--277, 2012.

\bibitem{ChSe10}
K.~C. Chan and I.~Seceleanu.
\newblock Orbital limit points and hypercyclicity of operators on analytic
  function spaces.
\newblock {\em Math. Proc. R. Ir. Acad.}, 110A(1):99--109, 2010.

\bibitem{ChaErnMen16}
S.~{Charpentier}, R.~{Ernst}, and Q.~{Menet}.
\newblock {$\Gamma$-supercyclicity.}
\newblock {\em {J. Funct. Anal.}}, 270(12):4443--4465, 2016.

\bibitem{charpentier2019chaos}
S.~Charpentier, K.~Grosse-Erdmann, and Q.~Menet.
\newblock Chaos and frequent hypercyclicity for weighted shifts.
\newblock {\em Ergodic Theory and Dynamical Systems}, pages 1--37, 2019.

\bibitem{CosManPar14}
G.~Costakis, A.~Manoussos, and I.~Parissis.
\newblock Recurrent linear operators.
\newblock {\em Complex Anal. Oper. Theory}, 8(8):1601--1643, 2014.

\bibitem{Fur81}
H.~Furstenberg.
\newblock {\em Recurrence in ergodic theory and combinatorial number theory}.
\newblock Princeton University Press, Princeton, N.J., 1981.
\newblock M. B. Porter Lectures.

\bibitem{GriLop23}
S.~Grivaux and A.~L\'{o}pez-Mart\'{\i}nez.
\newblock Recurrence properties for linear dynamical systems: an approach via
  invariant measures.
\newblock {\em J. Math. Pures Appl. (9)}, 169:155--188, 2023.

\bibitem{GriLop25}
S.~Grivaux, A.~L\'{o}pez-Mart\'{\i}nez, and A.~Peris.
\newblock Questions in linear recurrence {I}: the {$T\oplus T$}-recurrence
  problem.
\newblock {\em Anal. Math. Phys.}, 15(1):Paper No. 1, 26, 2025.

\bibitem{GriLopPerQuestions2}
S.~Grivaux, A.~L\'{o}pez-Mart\'{\i}nez, and A.~Peris.
\newblock Questions in linear recurrence {II}: lineability properties.
\newblock {\em Banach J. Math. Anal.}, 19(4):Paper No. 61, 28, 2025.

\bibitem{grivaux2024orthogonality}
S.~Grivaux, E.~Matheron, and Q.~Menet.
\newblock Orthogonality of invariant measures for weighted shifts.
\newblock {\em Pure and Applied Functional Analysis}, 2024.

\bibitem{Gro19bilateral}
K.-G. Grosse-Erdmann.
\newblock Frequently hypercyclic bilateral shifts.
\newblock {\em Glasg. Math. J.}, 61(2):271--286, 2019.

\bibitem{GroPer11}
K.-G. {Grosse-Erdmann} and A.~{Peris Manguillot}.
\newblock {\em {Linear chaos}}.
\newblock Berlin: Springer, 2011.

\bibitem{he2018jf-class}
S.~He, Y.~Huang, and Z.~Yin.
\newblock {$J^\mathcal F$}-class weighted backward shifts.
\newblock {\em International Journal of Bifurcation and Chaos}, 28(06):1850076,
  2018.

\bibitem{Lop24Locallybounded}
A.~L\'{o}pez-Mart\'{\i}nez.
\newblock Invariant measures from locally bounded orbits.
\newblock {\em Results Math.}, 79(5):Paper No. 185, 30, 2024.

\bibitem{lopez2024recurrent}
A.~L{\'o}pez-Mart{\'\i}nez.
\newblock Recurrent subspaces in {B}anach spaces.
\newblock {\em International Mathematics Research Notices},
  2024(11):9067--9087, 2024.

\bibitem{LopMen25}
A.~L\'{o}pez-Mart\'{\i}nez and Q.~Menet.
\newblock Two remarks on the set of recurrent vectors.
\newblock {\em J. Math. Anal. Appl.}, 541(1):Paper No. 128686, 17, 2025.

\bibitem{lopez2025shifts}
A.~L{\'o}pez-Mart{\'\i}nez and D.~Papathanasiou.
\newblock Shifts on trees versus classical shifts in chain recurrence.
\newblock {\em Journal of Differential Equations}, 433:113230, 2025.

\bibitem{MarMenPui22}
O.~Martin, Q.~Menet, and Y.~Puig.
\newblock Disjoint frequently hypercyclic pseudo-shifts.
\newblock {\em J. Funct. Anal.}, 283(1):Paper No. 109474, 44, 2022.

\bibitem{MurPer13}
M.~Murillo-Arcila and A.~Peris.
\newblock {Strong mixing measures for linear operators and frequent
  hypercyclicity.}
\newblock {\em J. Math. Anal. Appl.}, 398(2):462--465, 2013.

\end{thebibliography}


\end{document}
